\section{Primary theorem applications}
\begin{longtable}{p{.26\textwidth}p{.65\textwidth}}
\toprule Input & Exact application boundary\\\midrule\endhead
Selmer parity \cite{dd} & Rational two-isogeny, $p=2$, full Selmer corank;
no finite-$\Sha$ hypothesis.\\
Kato \cite{kato} & Fixed weight-two form and symbol, integral transfer and
fixed-valuation smoothing; rational divisibility is Theorem 13.4(2),
with oddness, rational irreducibility and the unipotent condition.\\
BFH and Hoffstein--Luo \cite{bfh,hl} & Actual signed-prime twist with
trivial character and arbitrary finite splitting set. Only the
nonvanishing clause has the four-prime bound before the extra prime.\\
Waldspurger and Ramsey \cite{waldspurger,ramsey} & Full good-Hecke
eigenspace, same dyadic squareclass, and algebraic vectors by flat
field extension. No arbitrary-ring base change in weight $3/2$.\\
CM norms and heights \cite{ccl,cst,cornutjetchev} & Split level,
conductors coprime to level and the new prime, no extra units,
unaveraged sums and one fixed parametrization.\\
Duality \cite{milne,nekovar} & Actual derived Kummer conditions and their
homotopies. Support change only after the coefficient representation
is unramified at the removed place.\\
Forward finiteness \cite{gz,kolyvagin} & Only at known analytic
nonvanishing/simple zeros or a known nontorsion Heegner trace.\\
Isogeny ratio \cite[proof of I.7.3]{milne} & Finite $\Sha$, dual
sum/difference maps, canonical heights and pullback measures;
no BSD leading-term assumption.\\
Height \cite{pan2022,javanpeykar,gr2014,fowler,cmheight,pazuki} &
Odd irreducible de Rham representation at five; explicit conductor
and modular-Jacobian transfer. CM order class number bounded by
the degree of the original Hom field.\\
Primes \cite{greentao,moebius,gtz,bmor} & Fixed finite-complexity system
after fixing sieve cutoff, and the unconditional explicit progression
estimate.\\\bottomrule
\end{longtable}
The source package includes the specification, natural precision polynomial
and construction sources in \nolinkurl{SUPPLEMENT_MANIFEST.json}.
The larger proof corpus is identified by hashes and stable locators in
\nolinkurl{SOURCE_MAP.json}. A source identity is not a proof verdict.

% End 06-literal-construction.tex


% Begin 07-even-source.tex
\section{Fixed-form arithmetic and the even coefficient source}\label{app:even-source}
This appendix supplies the constructions used in (3.1)--(3.3).
Throughout, $E=E_\ell$ is fixed, with the lattice, parametrization and
period convention of Section~\ref{sec:rc}. Constants may depend on these
data and on the prescribed finite list of local types. They do not depend
on the number or sizes of the later good conductor primes.

\subsection{One retained class, including discriminant \texorpdfstring{$4h$}{4h}}
All numbered primary references in this subsection are to Kato~\cite{kato}.

Fix a non-CM elliptic curve $E/\mathbf Q$ of conductor $N$, its weight-two newform $f$,
$T=T_2E$, and $V=T\otimes\mathbf Q_2$. In the application
$E=E_\ell$ is fixed first. Put
$\Lambda=\mathbf Z_2[[\Gamma]]=\mathbf Z_2[[t]]$, where $\Gamma$
is the real cyclotomic Galois group. Let
$\kappa:G_{\mathbf Q}\to\Gamma\subset\Lambda^\times$ be the quotient map,
and set $A=\Lambda[1/2]$.
We use arithmetic Frobenius throughout. The action on the cyclotomic
group-ring coefficient is the inverse deck action.

\paragraph{The precise output.}
\begin{proposition}[One retained Kato system and its two comparisons]
\label{ex:app-kato-system}
Use the fixed curve, lattice and cyclotomic convention just specified.
Fix $e\ge\max\{3,v_2(N)\}$ large enough for the fixed initial local data;
the smoothing choice is (K5).
Let $G$ be a finite elementary abelian two-group, let
$\psi:G_{\mathbf Q}\to G$ be a continuous homomorphism, and let
$\chi_x=\chi_0(x\circ\psi)$, for $x\in\widehat G$, be a finite
family of quadratic characters with one fixed real sign and fixed
local squareclasses at the primes dividing $2N$.
Since the characters of $G$ separate
its elements, $\psi(c)=1$ for complex conjugation $c$, and
$\psi|_{G_{\mathbf Q_v}}=1$ at every $v\mid2N$. All varying conductor primes
are good odd primes for $E$. Let $F$ be the squarefree least common
multiple of the odd parts of the primitive conductors. There is a
single retained-smoothed class
\[
 Z_F\in 2^{-c_E}H^1\bigl(G_{\mathbf Q,S},
        T\otimes\Lambda[G](\chi_0\psi\kappa^{-1})\bigr),
 \qquad S=\operatorname{prime}(2NF),                         \tag{K1}
\]
where $c_E$ is fixed before $F$ and $G$.
There is also one $F_E\in\Lambda$ with $F_E(0)\ne0$, fixed before
$F,G$, such that at every height-one prime $\mathfrak p$ of $A$,
\[
 \ell_{A_{\mathfrak p}}H^2(C_{S,x})_{\mathfrak p}
 \leq
 \ell_{A_{\mathfrak p}}
 \left(H^1(C_{S,x})/A F_E(Z_F)_x\right)_{\mathfrak p}.       \tag{K2}
\]
Here $C_{S,x}$ is the full rational global Iwasawa complex of the
constituent. For an odd squarefree signed parameter $h$ giving
$\chi_x=\chi_h$, central reciprocity for this same class gives
\[
 v_2\exp^*\bigl(F_E(Z_F)_x(0)\bigr)
 =v_2L_*(h)+\sum_{q\mid F,\ q\nmid h}
        v_2\#E^h(\mathbf F_q)+2e+v_2F_E(0)+\epsilon_h,       \tag{K3}
\]
when the value is nonzero. The numbers $\epsilon_h$ belong to a
bounded set depending only on the fixed curve, period conventions
and local types at two. A zero central $L$-value gives zero dual
exponential. Formula (K3) does not assert that its remaining unit is
one. No part of (K1)--(K3) compares integral lattices of varying newforms.
\end{proposition}
\begin{proof}\leavevmode
\paragraph{The retained route and its comparison.}
All classes in this subsection retain the smoothing integers of (K5).
The comparison used below is the exact identity (E.1), obtained by
integral weighted corestriction and the norm relation of
\cite[Proposition 8.12]{kato}. We do not divide by the smoothing
operators, or compare a chosen generator with a separately normalized
zeta class for the varying newform. In particular the central factor
is $s_F$ of (K6), whose valuation is $2e$ for every $F$.
The use of power-series zero finiteness in paragraph~7 establishes
nonvanishing of the supplied Euler system in both generic sign
components; it is not a comparison of two Euler systems.

\paragraph{1. One original lattice and one symbol combination.}
Kato (8.1.3), p.~180, constructs retained-smoothed classes
\[
 {}_{c,d_s}z_m(f,1,1,\xi,S_m)
 \in H^1\bigl(\mathbf Z[\zeta_m,1/2],L_f(1)\bigr),
 \qquad S_m=\operatorname{prime}(2mN),                      \tag{K4}
\]
where $L_f=V_{\mathcal O_\lambda}(f)$ is the image lattice of the
fixed form. For $\xi\in\mathrm{SL}_2(\mathbf Z)$, section~8.9,
p.~185, explicitly constructs the trace homomorphism into the
integral induced lattice
\[
 \mathbf Z_2[\operatorname{Gal}(\mathbf Q(\zeta_m)/\mathbf Q)]
       \otimes V_{2,\mathbf Z_2}(Y_1(N)).
\]
Section~8.11 projects this lattice onto $L_f$. Auxiliary modular
curve levels used in the trace do not change its target lattice.
Choose a rational identification $L_f(1)\otimes\mathbf Q_2=V$.
One fixed $2^{c_E}$ sends its lattice into $T$.

By Kato~13.6--13.7, p.~227, the symbol span has a nonzero projection
to each Betti sign line. The union of the two rational kernels does
not contain the whole span. Choose an integral finite linear
combination $\xi$ with both projections nonzero, clearing its fixed
denominators once. Norm relations and reciprocity are linear, so
the same combination works at every $m$. Although (K4) is written
with one symbol, it henceforth denotes this fixed combination.

\paragraph{2. Retain one smoothing choice for the whole family.}
Fix $e\geq\max\{3,v_2(N)\}$, enlarged only for fixed initial data,
and take
\[
                 c=d_s=1+2^e\,3N_{\rm odd}F.              \tag{K5}
\]
These integers exceed one, are prime to $6NF$, are one modulo $N$
and $F$, and are one modulo every dyadic quadratic conductor under
consideration. Distinctness or primality of $c,d_s$ is not a
hypothesis of Kato's construction. The exponents $(u,v)$ in (4.2.4),
p.~143, are $(1,1)$ when $k=2,r=r'=1$. The retained-smoothed
version of Theorem~6.6, p.~163, therefore has central scalar
\[
 s_F=(c^2-c)(d_s^2-d_s)=c(c-1)d_s(d_s-1),
                   \qquad v_2(s_F)=2e.                   \tag{K6}
\]
No smoothing factor is divided out. The prime factors of $cd_s$
exclude some Euler-system indices in Kato~13.1; they are not added
to the inverted arithmetic support $S$, as is explicit in (8.1.3)
and in the definition of $H^q$ in section~8.2. In particular their
number creates no local $H^2$ correction.

\paragraph{3. The integral $j_*$ condition survives ramified transfer.}
For a number field $K$, let $X_K=\operatorname{Spec}\mathcal O_K[1/2]$
and $j_K:\operatorname{Spec}K\to X_K$. At finite coefficients
$M=T/2^aT$, the low-degree Leray sequence gives an injection
\[
 H^1(X_K,j_{K,*}M)\longrightarrow H^1(K,M)
\]
whose image is the kernel of all inertia restrictions outside two.
Indeed the higher-direct-image stalk formula and Leray sequence
\cite[Tags 03Q9 and 03QA]{stackscohom} identify the stalks of
$R^1j_{K,*}M$ at a finite place with $H^1(I_v,M)$;
the edge map is restriction to these stalks. The generic-point map
$j_K$ is affine, hence quasi-compact and quasi-separated, as required
by the stalk formula. Kato's section~8.2,
pp.~180--181, defines the integral groups by the inverse limit of
these finite-coefficient groups, so this assertion applies to the
classes (K4) before rationalization.

Let $L/K$ be finite, with absolute Galois groups $H\subset G_K$.
For a class $z\in H^1(H,M)$ and an inertia subgroup $I\subset G_K$,
the Mackey formula is
\[
 \operatorname{res}_{I}^{G_K}\operatorname{cor}_{H}^{G_K}z
 =\sum_{g\in I\backslash G_K/H}
   \operatorname{cor}_{I\cap gHg^{-1}}^{I}
   \operatorname{res}_{I\cap gHg^{-1}}^{gHg^{-1}}(g z).    \tag{K7}
\]
The intersection in each summand is the upper inertia subgroup
for the corresponding place of $L$. If all upper inertia
restrictions vanish, every summand is zero. This does not require
$L/K$ to be unramified. It is a finite double-coset identity at each
finite coefficient precision. Uniqueness under the Leray injection
makes the extensions across $X_K$ compatible in $a$. Thus the
integral $j_*$ condition survives corestriction, including at primes
ramified in a quadratic twisting field.

\paragraph{4. Primitive twisting is an actual Euler system.}
Fix a constituent $\chi=\chi_h$ of primitive conductor
$D=|\Delta_h|=2^a d$, $d$ odd, and a basis $e_\chi$ of
$\mathbf Z_2(\chi)$. Put
$\Sigma=\operatorname{prime}(2Ndcd_s)$.
For every positive integer $m$ with $(m,cd_s)=1$, set
\[
 K_m=\mathbf Q(\zeta_m),\quad
 L_m=\mathbf Q(\zeta_{Dm}),\quad
 y_m=\operatorname{cor}_{L_m/K_m}
             (z_{Dm}\otimes e_\chi).                      \tag{K8}
\]
The character is trivial over $L_m$, so tensoring is an integral
map there. Formula (K7) places (K8) in the precise $j_*$ group
required by Kato~13.1 on its index set
$\operatorname{prime}(m)\cap\Sigma=\{2\}$. The larger definition is
permitted by (8.1.3), Sections~8.9--8.11 and Proposition~8.12:
$S_{Dm}=\operatorname{prime}(2NDm)$ is disjoint from
$\operatorname{prime}(cd_s)$, since $d\mid F$ and (K5) holds.
Each class uses the original $f,\xi,c,d_s$. We will also use this
same construction on the different index set required by
\cite[Section 0.1]{kato1999}, as checked below.

For $m\mid m'$, transitivity of corestriction reduces the norm
of (K8) to Proposition~8.12, p.~186. For a deck automorphism $g$,
equivariance gives
\[
 \operatorname{cor}((gz)\otimes e_\chi)
       =\chi(g)^{-1}g\operatorname{cor}(z\otimes e_\chi).
\]
Consequently the factor at a newly introduced good prime $q$ is
\[
 E_q(\chi)=1-\frac{a_q\chi(q)}q\sigma_q^{-1}
                         +\frac1q\sigma_q^{-2}.           \tag{K9}
\]
In Kato~13.1, pp.~224--225, the polynomial is
$P_q(X)=\det(1-\operatorname{Fr}_qX\mid T\chi)$ and the norm
factor is $P_q(q^{-1}\sigma_q^{-1})$, exactly (K9).
Repeated prime powers introduce no factor because Proposition~8.12
has a product over newly added support primes. Thus the subfamily
on Kato~13.1's index set is an Euler system of the stated lattice,
not a class inferred from proportionality of rank-one Iwasawa
modules. The required norm relations on the 1999 index set are
verified separately below.

For $m=2^n$ with $n\geq a$, the field $L_m$ can be replaced by
$\mathbf Q(\zeta_{d2^n})$: apply repeated-two norm compatibility
from exponent $n+a$ to $n$. This introduces no factor $2^a$.
The construction therefore covers discriminant $4h$ on the same
cofinal tail. At the coefficient level it retains the full character
$\chi_2\chi_{\rm odd}$. A dyadic character nontrivial on $\Gamma$
changes its action by the integral automorphism
$1+t\mapsto-(1+t)$, or $t\mapsto-2-t$. A dyadic sign character
instead changes the finite real-sign component. Neither operation
is a division by a sign projector.

\paragraph{5. One coefficient map produces every constituent.}
For $n$ large enough to contain the dyadic characters, put
$A_{F,n}=\operatorname{Gal}(\mathbf Q(\zeta_{F2^n})/\mathbf Q)$.
Apply Shapiro to $z_{F2^n}$, then the map on group bases
\[
 \Phi_n:\mathbf Z_2[A_{F,n}]\longrightarrow
                  \mathbf Z_2[\Gamma_n\times G],\qquad
 [s]\longmapsto\chi_0(s)[\bar s][\psi(s)].                 \tag{K10}
\]
With inverse deck actions, both sides transform under $g\in G_{\mathbf Q}$
by the inverse of its group image. On the target the extra scalar
is $\chi_0(g)$; since $\chi_0^{-1}=\chi_0$, (K10) is equivariant.
This is an integral coefficient homomorphism. Shapiro restriction
to an inertia subgroup is the direct sum of the upper restrictions;
the same zero restrictions used in (K7) therefore remain zero after
(K10). In particular the map acts on actual integral classes.
The map in (K10) is defined on this finite group for the following
reason. The constituent indexed by the trivial character of $G$ is
$\chi_0$, and
$x\circ\psi=\chi_x\chi_0$ for every $x\in\widehat G$.
All these quadratic Dirichlet characters have conductor dividing
$F2^n$ once $n$ is large enough. If an element of $G_{\mathbf Q}$
acts trivially on $\mathbf Q(\zeta_{F2^n})$, every character of $G$
is therefore one on its image under $\psi$. Character separation
makes that image the identity. Thus both $\chi_0$ and $\psi$
factor through $A_{F,n}$. Explicitly, for $g\in G_{\mathbf Q}$,
\[
 \Phi_n([g^{-1}s])
 =\chi_0(g)\,[\bar g]^{-1}[\psi(g)]^{-1}\Phi_n([s]).
\]
This verifies equivariance on each integral group basis before
evaluation at any character. The quotient to the real tower sends
each basis element to its image; its matrix has entries zero and
one. The transition maps and character evaluations are also maps
of integral coefficient modules. None of these operations is the
idempotent $|G|^{-1}\sum_g x(g)[g]$.


The maps $\Phi_n$ commute with the norm transition maps.
Proposition~8.12 has no new prime when $n$ increases, so their
images define the single class $Z_F$ in (K1).
The finite-coefficient degree-zero groups are finite and satisfy
Mittag--Leffler. The inverse-limit exact sequence
\cite[Corollary 4.1.3, p.~96]{nekovar} therefore has zero
$\varprojlim^1 H^0$ term, so it gives this degree-one class
in completed cochains.
The quotient from the full cyclotomic tower to $\Gamma$ is the
integral quotient of group bases, equivalently corestriction to
the real subfield. No factor $(1+\tau)/2$ occurs.

After evaluation at $x$, (K10) is the coefficient form of the
unnormalized corestriction weighted by $\chi_x$. Factor that
corestriction through the primitive conductor field in (K8).
Transitivity and Proposition~8.12 give the exact identity
\[
              (Z_F)_x=
       \prod_{q\mid F/d}E_q(\chi_x)\,y_{2^\infty}.        \tag{E.1}
\]
It is an identity of the given classes. The same smoothing choice
appears on both sides. There is no group-order denominator and no
comparison with a canonical zeta class of $f\otimes\chi_x$.

For clarity, at a constituent write $\chi_h=\chi_2\chi_{\rm odd}$ and
$|\Delta_h|=2^a d$. Before taking the cofinal limit, the corresponding
scalar coefficient map is exactly
\[
 [s]\longmapsto\chi_2(s)\chi_{\rm odd}(s)[s|_{\Gamma_n}],
                    \qquad n\ge\max\{a,2\}.\tag{E.2}
\]
Thus the dyadic character is part of the module action, including
when the field discriminant is $4h$.

\paragraph{6. Central trace and period normalization.}
First use (E.1) to remove the unused tame primes. The central
specialization of (K8) is the corestriction of $z_D\otimes e_\chi$
from $\mathbf Q(\zeta_D)$ to $\mathbf Q$, with omitted support
$\operatorname{prime}(2Nd)$. For $D=1$ use the trivial character
and the same omitted support. This follows from the same repeated-two
norm identity even when $D=4|h|$. One must do this before using a
Gauss basis; the Gauss sum of an imprimitive character modulo a
larger cyclotomic level can be zero.

Theorem~9.7, p.~189, identifies the dual exponential of $z_D$ with
the retained-smoothed de Rham zeta element. Dual exponential commutes
with corestriction, whose de Rham map is field trace. Thus the
coefficient of $e_\chi$ is the exact unnormalized sum
\[
                    \sum_{g\in\operatorname{Gal}(\mathbf Q(\zeta_D)/\mathbf Q)}
                       \chi(g)\,g(\exp^*z_D).             \tag{K12}
\]
Theorem~6.6, p.~163, evaluates precisely this sum. In weight two,
$r=r'=1$, its Betti sign is $\epsilon=\chi(-1)$ and its power of
$2\pi i$ is zero. If $e_B^\epsilon$ is a fixed rational generator
of that Betti line and
$\delta(f,1,\xi)^\epsilon=b_\epsilon e_B^\epsilon$, then
$b_\epsilon\in\mathbf Q^\times$ is fixed and
\[
     \operatorname{per}_f\bigl((\text{K12})\bigr)^\epsilon
          =s_F L_{\operatorname{prime}(2Nd)}(f,\chi,1)\,
                        b_\epsilon e_B^\epsilon.          \tag{K13}
\]
Choose the original de Rham generator so that its Betti period
on this line is $\Omega_E^\epsilon e_B^\epsilon$, absorbing a
fixed rational comparison in $b_\epsilon$. Formula (K13) says
that the scalar multiplying $e_\chi$ in (K12) is
$s_F b_\epsilon L_{\operatorname{prime}(2Nd)}(f,\chi,1)/
\Omega_E^\epsilon$. Its Galois character is $\chi$.

To descend it, take the invariant basis $\sqrt h\,e_\chi$ of the
quadratic character line. With the usual complex convention
$\Omega_E^-=i|\Omega_E^-|$ and $\sqrt h=i\sqrt{|h|}$ for $h<0$,
division by $\sqrt h$ gives, up to the already fixed rational
period choice, exactly
\[
             s_F b_\epsilon\,\frac{L_*(h)}{h}
       \prod_{p\mid 2N}\text{(the omitted local factor at $p$)}.
                                                               \tag{K14}
\]
The signed integer $h$ is odd, so $h^{-1}$ has valuation zero.
This computation also explains the usual Gauss-basis description:
for the primitive character, $\tau(\chi)^2=\Delta_h$ and
$\tau(\chi)e_\chi$ is invariant; replacing it by $\sqrt h e_\chi$
costs the fixed dyadic factor $\sqrt{\Delta_h/h}$. It does not cost
$\varphi(D)$ or any power of the number of prime divisors of $h$.
The discriminant-$4h$ case has ratio two in this last basis change,
and the direct basis calculation (K14) already includes it.
One can verify the normalization without choosing an idempotent.
Write $\beta=\exp^*z_D$ in the original de Rham line and put
$\beta_\chi=\sum_g\chi(g)g\beta$. Then
$u\beta_\chi=\chi(u)\beta_\chi$ for every $u$ in the finite
cyclotomic group. Thus $\beta_\chi\otimes e_\chi$ is invariant,
and its coordinate in the invariant basis
$\sqrt h\,e_\chi$ is $\beta_\chi/\sqrt h$.
Theorem~6.6 uses precisely $\beta_\chi$, rather than
$\beta_\chi/\varphi(D)$. Theorem~9.7 applies here since
$1\le r=r'=1\le k-1=1$; it has no restriction that $p$ be odd.
The dual form $f^*$ in Theorem~6.6 equals $f$, because the elliptic
newform has rational coefficients.

In the real-period convention,
$L_*(h)=\sqrt{|h|}L(E^h,1)/|\Omega_E^\epsilon|$, up to the
fixed rational period comparison already included in $b_\epsilon$.
For $h>0$,
$\Omega_E^+\sqrt h=|\Omega_E^+|\sqrt{|h|}$; for $h<0$,
$\Omega_E^-\sqrt h=-|\Omega_E^-|\sqrt{|h|}$.
Consequently
\[
 \frac{L(E^h,1)}{\Omega_E^\epsilon\sqrt h}
     =\frac{L_*(h)}{h}
\]
with the same fixed period comparison. This proves the signed
factor in (K14), not only its valuation. For a primitive quadratic
character, $\tau(\chi)^2=\Delta_h$ and
$\tau(\chi)/\sqrt h$ equals a fixed choice of
$\sqrt{\Delta_h/h}$. Changing from its Gauss basis to the displayed
invariant basis multiplies the coordinate by this ratio, which
is one or two up to sign. In particular the discriminant-$4h$
case adds no factor depending on the number of odd conductor primes.


Identifying this generator with the chosen local generator for
$E^h$ at two has bounded valuation: odd squarefree $h$ belongs to
one of finitely many classes in $\mathbf Q_2^\times/\mathbf Q_2^{\times2}$,
and changing a representative by a unit square scales the quadratic
twist differential by a two-adic unit. The local factors at $2N$
also form a fixed finite nonzero list. At a good active prime the
local $L$-factor is one; at an unused good prime (K9) at $t=0$ is
$\#E^h(\mathbf F_q)/q$. This proves (K3), including all its
uniformity claims and the assertion about zero values.

\paragraph{7. Every hypothesis of Kato~13.4(2).}
Apply the theorem to the integral system (K8), after multiplying
by the fixed lattice comparison if needed. Its hypotheses are:
\begin{itemize}
\item Each generic sign component is nonzero. For a finite-order
cyclotomic character $\rho$, corestriction and the $\rho$-sum combine
into the single $(\chi\rho)$-weighted sum on the original class.
Theorems~6.6 and~9.7 apply with the same $\xi,c,d_s$. Their smoothing
factors are nonzero, since $c,d_s>1$ cannot be roots of unity.
Theorem~13.5(2), p.~227, gives only finitely many zero $L$-values
among characters supported on the fixed finite set
$\operatorname{prime}(2Nd)$. There are infinitely many $\rho$ of
each parity. A torsion Iwasawa class can have a nonzero specialization
at only finitely many such characters, since a nonzero annihilating
power series has only finitely many cyclotomic zeros. Hence both
generic sign components are nonzero.
\item The two real eigenspaces of $T\chi$ have rank one.
\item The theorem uses inverse Frobenius in its purity condition.
On $T\chi$ its good eigenvalues have complex absolute value
$q^{-1/2}$, so its integer $w$ is $-1$.
\item $V\chi$ is irreducible; twisting preserves irreducibility of
the non-CM rational Tate representation.
\item The image over $\mathbf Q(\mu_{2^\infty})$ contains an open
subgroup of $\mathrm{SL}_2(\mathbf Z_2)$. Intersecting with
$\ker\chi$ keeps an open subgroup, which contains a nonidentity
unipotent. It fixes a one-dimensional subspace and gives condition~(v).
\end{itemize}
These are all rational hypotheses. The residual irreducibility,
integral cokernel and $p\ne2$ conditions of part~(3) are not used.
Theorem~13.4(1) gives torsion $\mathbf H^2(T\chi)$, and part~(2)
gives, at every height-one prime away from two,
\[
 \ell\mathbf H^2(T\chi)_{0,\mathfrak p}
       \leq\ell(A_{\mathfrak p}/J_{\mathfrak p}),\qquad
 \mathbf H^2_0=\ker(\mathbf H^2\to\mathbf H^2_{\rm loc}). \tag{K15}
\]

The finiteness and Euler characteristic (12.2.1)--(12.2.2), p.~220,
give generic rank one in $\mathbf H^1$. Its localization away
from two is torsion-free. Indeed a torsion element for a scalar
prime would come from invariants of the corresponding finite-dimensional
rational coefficient algebra at that prime; its coefficient action is
abelian, whereas the open commutator subgroup on $V\chi$ has no
fixed vectors. The coefficient exact sequence excludes these
invariants. The scalar specialization need not be a finite-order
character. Thus $\mathbf H^1_{\mathfrak p}$ is free of rank one.
Localization commutes with $\operatorname{Hom}$ for this finitely
presented module. Writing $y_{2^\infty}=a e$ in a local basis gives
$J_{\mathfrak p}=(a)$. Formula (K15) is exactly the index bound for
the supplied class, not for an arbitrarily rescaled generator.

\paragraph{The published prime-two Euler-system input.}
For $p=2$ we use the rational, away-from-two case of
\cite[Theorem 0.8, pp.~317--318]{kato1999}, with the stronger
hypothesis in its condition 0.6(ii), to justify the application of
\cite[Theorem 13.4(2)]{kato}. For our non-CM elliptic representation
the image of $G_{\Q^{\rm ab}}$ contains an open subgroup of
$\mathrm{SL}_2(\Z_2)$, as stated in
\cite[Example 0.6.2, pp.~317--318]{kato1999}, using Serre's open-image
theorem. A quadratic character $\chi$ is trivial on $G_{\Q^{\rm ab}}$.
Thus the same subgroup acts on $V\chi$ and contains a matrix
$\left(\begin{smallmatrix}1&b\\0&1\end{smallmatrix}\right)$ with
$b\ne0$. Its fixed space over $\Q_2$ has dimension one. This supplies
the unipotent in $G_{\Q^{\rm ab}}$, not merely in the cyclotomic
subgroup. The latter subgroup also contains this open
$\mathrm{SL}_2$ subgroup, so $V\chi$ remains absolutely irreducible
over $\Q(\mu_{2^\infty})$. Indeed, a line invariant under both upper
and lower nonidentity unipotents cannot exist, even over an algebraic
closure. No inference from irreducibility over $G_\Q$ alone is needed.

We check the index set and the supplied class in that theorem
explicitly. Put $N_K=2Nd$ and $N'_K=cd_s$. These integers are
coprime by (K5), since $d\mid F$. Section~0.1 of \cite{kato1999}
requires every index $m$ with $N_K\mid m$ and $(m,N'_K)=1$.
The enlarged definition (K8) supplies all these classes. They
originate in the integral $j_*$ group over
$\mathbf Z[\zeta_m,1/2]$ and map to full cohomology over
$\mathbf Z[\zeta_m,1/N_K]$. The varying construction support
$S_{Dm}$ is not an instruction to invert all its primes in this
last scheme. Apply the fixed lattice map $2^{c_E}$ from (K4) if
needed, identically at every level.

For a prime $\ell\nmid N'_K$, if $\ell\mid m$ then
$S_{Dm\ell}=S_{Dm}$, so Proposition~8.12 gives norm factor one.
If $\ell\nmid m$, the condition $N_K\mid m$ implies
$\ell\nmid2Nd$, and hence
$S_{Dm\ell}\setminus S_{Dm}=\{\ell\}$. The equivariance identity
before (K9) gives the factor
$P_\ell(\ell^{-1}\sigma_\ell^{-1})=E_\ell(\chi)$, with
$P_\ell(X)=1-a_\ell\chi(\ell)X+\ell X^2$. These are exactly the
two cases of (0.1.1) in \cite{kato1999}. Smoothing primes are
excluded by $N'_K$ and need not be inverted.

For the class in Section~0.7 of that source take its parameter
$d=1$; this parameter is distinct from our odd conductor part $d$.
For every $n\geq1$,
\[
 S_{DN_K2^n}=S_{D2^n}=\operatorname{prime}(2Nd).
\]
Proposition~8.12 therefore gives
\[
 \operatorname{cor}_{\mathbf Q(\zeta_{DN_K2^n})/
                          \mathbf Q(\zeta_{D2^n})}
             z_{DN_K2^n}=z_{D2^n}.
\]
Transitivity of corestriction, through the two intermediate fields
in the common upper field, now yields
\[
 \operatorname{cor}_{\mathbf Q(\zeta_{N_K2^n})/
                          \mathbf Q(\zeta_{2^n})}y_{N_K2^n}
 =\operatorname{cor}_{\mathbf Q(\zeta_{D2^n})/
                          \mathbf Q(\zeta_{2^n})}
                         (z_{D2^n}\otimes e_\chi)
 =y_{2^n}.
\]
No Euler factor occurs because the supports are equal. No field
degree occurs because this is corestriction of the actual upper
class, not corestriction of a restriction. These equalities commute
with transition maps, and $n\geq1$ is a cofinal tail. Write
$\xi_K$ for the Section~0.7 class, to distinguish it from the
modular-symbol combination $\xi$. It is exactly the full-cohomology
image of $y_{2^\infty}$, multiplied by $2^{c_E}$ if the fixed
lattice map was needed. That factor is a unit at every height-one
prime under consideration, so its coordinate ideal is precisely
$J_{\mathfrak p}$. The preceding
nonvanishing argument supplies a nontorsion class on each sign
component, purity gives $w=-1$, and oddness gives condition
0.6(iv$_{\mathfrak p}$). At the height-one primes used here, $2$ is
invertible and the localized Iwasawa algebra is a DVR. Theorem 0.8
then bounds the coordinate ideal of this actual retained class. Its
strong residual and integral-cokernel assumptions concern primes
containing $2$, which are not used here.

We also account for the difference between the two published
cohomology conventions. Write $C_j$ for the $j_*$ Iwasawa complex of
\cite[Section 12.1]{kato}, and $C_S$ for full cochains on the open
scheme with the odd support inverted, as in \cite[Section 0.5]{kato1999}.
Work in one rational sign component and localize at $\mathfrak p$.
Their quotient is the direct sum of the odd local complexes (K16).
Each differential has nonzero determinant, so its degree-one
cohomology vanishes. Hence $H^1(C_j)\simeq H^1(C_S)$ and
$H^2(C_j)\hookrightarrow H^2(C_S)$. At an odd prime $q$, the
unramified local term has cohomological dimension one after inverting
$2$. The localization triangle identifies the degree-two cohomology
of (K16) with full local $H^2(\Q_q,T\chi)$ in Iwasawa cohomology.
The map from $H^2(C_S)$ is the usual local restriction. Exactness
therefore identifies $H^2(C_j)$ with the kernel of all those odd
local restrictions. Intersecting with the kernel of restriction at
$2$ gives
\[
 \ker\bigl(H^2(C_j)\to H^2_{\rm loc,2}\bigr)
 \simeq
 \ker\bigl(H^2(C_S)\to\!\bigoplus_{q\in S}H^2_{\rm loc,q}\bigr).
\]
This is precisely the comparison of the subscript-zero groups in
the two sources. The $H^1$ identification preserves the actual
Euler-system class and its coordinate ideal $J$. Thus Theorem 0.8
gives (K15) in the convention used here. The argument for (K16)
below is local and does not use (K15). We do not invoke the
odd-prime-only Perrin--Riou argument or an unchecked version of
Rubin's theorem at $2$.

\paragraph{8. Paying for the full global complex.}
At an odd prime put $W_q=H^1(I_q,V\chi)$. The rational cone from
$j_*$ to full cochains is
\[
 [\,W_q\otimes A\xrightarrow{1-\sigma_q^{-1}\operatorname{Fr}_{W_q}}
          W_q\otimes A\,],\qquad\deg=(1,2).              \tag{K16}
\]
Its determinant $D_q$ is nonzero, since the cyclotomic Frobenius
exponent is nonzero. The map is injective even after localization
at a prime dividing $D_q$; it is not being reduced modulo that
prime. Thus the cone contributes only degree-two length
$v_{\mathfrak p}(D_q)$. The long exact sequence shows that full
$H^2$ increases by at most their sum and that $H^1$ is unchanged.
This inequality does not require vanishing of the next $j_*$ group.

At a good inactive prime, $W_q=V\chi(-1)$, so $D_q=E_q(\chi)$
with exactly the convention of (K9). At a good active prime,
ramified quadratic inertia acts by $-1$; its rational cohomology
vanishes and $D_q=1$. This includes every active base prime,
regardless of how many there are. At an old bad prime retain
$H^1(I_q,V\chi)$ itself; replacing it by $V^{I_q}\chi(-1)$
would in general be false.

Choose a fixed integral $F_2$ divisible over $A$ by the
characteristic elements of $\mathbf H^2_{\rm loc}(T\chi)$
for the finitely many allowed local types at two and finite
sign components. Its central value can be chosen nonzero:
derived local control, and the absence of local Iwasawa $H^3$,
identify its rational central $H^2$ with $H^2(\mathbf Q_2,V\chi)=0$.
The last equality follows from local duality and finiteness of
local elliptic torsion. Hence the characteristic elements have
no factor $t$. Kato (12.2.3), p.~220, gives the required finiteness
and torsion of this local Iwasawa module.

At the finitely many old bad primes, choose stable lattices in
the $W_q$ and multiply all the determinants $D_q$ for the finitely
many local twisting types. Clear their fixed scalar denominators.
Their central values are nonzero, since the cokernel of
$1-\operatorname{Fr}_{W_q}$ is $H^2(\mathbf Q_q,V\chi)=0$
by the same local duality. Call this product $F_{\rm bad}$ and set
$F_E=F_2F_{\rm bad}$.

The quotient $\mathbf H^2/\mathbf H^2_0$ injects into the local
$\mathbf H^2_{\rm loc}$, so $F_2$ pays for it in (K15).
Multiplying a class in a free rank-one module by a nonzero scalar
adds that scalar's valuation to its index. Formula (E.1) already
pays for every unused good determinant in (K16); $F_{\rm bad}$
pays for the remaining old ones. Adding these length inequalities
proves (K2).
Here is the length calculation with all contributions separated.
Fix a height-one prime $\mathfrak p$ of $A$ and write
$R=A_{\mathfrak p}$. Let $M$ be the localized degree-one module
for the primitive $j_*$ system, let $y=y_{2^\infty}$, and write
$\operatorname{ind}_R(y)=\ell_R(M/Ry)$.
The preceding freeness argument gives $M\simeq R$.
Put $Q_\chi=\prod_{q\mid F/d}E_q(\chi)$.
All primes in $F/d$ are good: the family has identical local
characters at primes dividing $N$, so a bad prime cannot ramify
in one constituent and be unramified in another.
The cone (K16) has zero degree-one cohomology, hence identifies
$M$ with the localized full degree-one module. Its long exact
sequence and (K15) give
\[
\begin{split}
 \ell_R H^2(C_{S,x})_{\mathfrak p}
 &\le \ell_R\mathbf H^2(T\chi)_{0,\mathfrak p}
       +\ell_R\mathbf H^2_{\rm loc}(T\chi)_{\mathfrak p}
       +\sum_{q\in S,\ q\ne2}v_{\mathfrak p}(D_q)\\
 &\le \operatorname{ind}_R(y)
       +v_{\mathfrak p}(F_2)+v_{\mathfrak p}(F_{\rm bad})
       +v_{\mathfrak p}(Q_\chi)\\
 &=\operatorname{ind}_R(F_EQ_\chi y)
  =\operatorname{ind}_R\bigl(F_E(Z_F)_x\bigr).
\end{split}
\]
The second line uses $D_q=1$ at active good primes, and
$D_q=E_q(\chi)$ at inactive good primes. It does not charge an
extra fixed constant for each active prime.

For completeness, nonvanishing at the center of the fixed
correction factors follows before multiplying over local types.
Local derived control and $H^3_{\rm Iw}=0$ identify
$\mathbf H^2_{\rm loc}/t\mathbf H^2_{\rm loc}$ rationally with
$H^2(\mathbf Q_2,V\chi)$. Local duality identifies the latter
with the dual of $H^0(\mathbf Q_2,V\chi)$, which is zero because
local elliptic torsion is finite. Nakayama at $(t)$ therefore
excludes a factor $t$ from its characteristic element.
At an old odd prime the cokernel of
$1-\operatorname{Fr}_{W_q}$ is $H^2(\mathbf Q_q,V\chi)=0$,
by the same duality. Its determinant at $t=0$ is nonzero.
Taking products over the finite list of local types, followed by
one fixed power of two to clear scalar denominators, gives
$F_E\in\Lambda$ with $F_E(0)\ne0$. This operation is independent
of the binary dimension and all later good conductor primes.


\paragraph{9. Uniformity and the boundary of this proposition.}
The only lattice denominator is fixed before the family. The
smoothing integers may grow, but their central valuation is $2e$.
The multiplier $F_E$ depends only on the original local types.
Transfer, coefficient push, real quotient and norm compatibility
are integral. The proof never divides by a conductor degree,
$|G|$, a class number, or a Fourier projector. It keeps one class
through (K1), (K2), and (K3).
\end{proof}

We now use this same class in the positive complex. The preceding
argument supplies its integral construction, its rational height-one
divisibility and its exact central reciprocity normalization.

\begin{proposition}[One integral even scalar for the family]\label{ex:app-even-scalar}
Retain Proposition~\ref{ex:app-kato-system} and assume, as for the
present curve, that $E[2](\Q)\simeq(\Z/2\Z)^2$. Let $C^+$ be the
fiber of the ordinary global cochains to the ordinary real cochains,
with real degree zero retained. It has a compatible finite free model
under the coefficient changes used here. For
$B=\mathscr O[G]$, it satisfies $C^+\otimes B\simeq B^2[-1]$.
The fixed real invariant generator, the positive lift of the retained
class, and the original determinant basis define the single scalar
$u^{\rm ev}$ by (E.4), with
\[
 u^{\rm ev}\in\Lambda[1/2][G],\qquad
 2^{c_E}u^{\rm ev}\in\Lambda[G].
\]
Here $c_E$ may be enlarged for the fixed data but is independent of
$G$ and the later good conductor primes. The assertion requires the
common real sign and local squareclasses at $2N$ in
Proposition~\ref{ex:app-kato-system}.
\end{proposition}
\begin{proof}\leavevmode
Let $C^+$ be the fiber of ordinary global cochains to ordinary real
cochains. Real $H^0$ is included. Over a finite Artin quotient,
continuous cochains are filtered unions of finite free function
modules, hence flat and therefore free. A residual splitting into
finite cohomology and disks lifts by the nilpotent perturbation
series to a finite free model with actual homotopies. Minimal models
at successive quotients are related by invertible matrices:
a homotopy equivalence reduces to an isomorphism of minimal residual
complexes, and Nakayama applies. Lifting these matrices and comparison
homotopies makes the models compatible. Their acyclic, termwise
surjective comparison cones have acyclic inverse limit.
Here is the compatibility argument. At successive finite coefficient
quotients, reduction of a minimal finite free model is a minimal model
of the reduced cochains. Two such models are homotopy equivalent;
reducing a homotopy inverse modulo the maximal ideal shows that the
comparison is degreewise invertible. Lift these invertible matrices
to the next quotient and change bases. The model differentials then
reduce to the preceding differentials exactly. The comparison maps
to the original cochains agree after reduction up to chain homotopy.
Since reduction of continuous cochains is termwise surjective, lift
that homotopy and adjust the upper comparison by its boundary. The
comparison maps are now exactly compatible and remain
quasi-isomorphisms.

Let $\mathcal K_n$ be their acyclic comparison cones. Their transition
maps are termwise surjective. If $s$ denotes the transition shift on
the product, successive lifting gives the exact sequence of complexes
\[
 0\longrightarrow\varprojlim_n\mathcal K_n
 \longrightarrow\prod_n\mathcal K_n
 \xrightarrow{\,1-s\,}\prod_n\mathcal K_n
 \longrightarrow0.
\]
Products of complexes of modules preserve acyclicity, so the long
exact cohomology sequence makes $\varprojlim_n\mathcal K_n$ acyclic.
The inverse-limit comparison is therefore a quasi-isomorphism from
the bounded finite free limit model to the continuous coefficient
complex. This also identifies its derived coefficient changes with
the finite-quotient constructions.
This proves the needed perfection and coefficient change.

% Replace baseline lines 3310--3335; retains equation tag (E.3).
\paragraph{The residual positive complex.}
Write $K_n=\mathbf Q(\zeta_{2^{n+2}})^+$, and let $S_n$ consist of
the places above two and a fixed finite set of odd rational primes.
Put $r_n=2^n$ and $s_n=\#S_n$. We first compute with the scalar
coefficient $\mathbf F_2=\mu_2$.

The ordinary two-class group of $K_n$ is trivial. Indeed this holds
for $K_0=\mathbf Q$. Assuming it for $K_{n-1}$, let $H$ be the
Hilbert two-class field of $K_n$, with the real places split.
The canonical extension $H/K_{n-1}$ is Galois with a finite
two-group $G_H$. Only the unique dyadic prime can ramify, and its
inertia $I$ has order two. The quotient of $G_H^{\rm ab}$ by the image
of $I$ gives an everywhere unramified abelian two-extension of
$K_{n-1}$ with real places split, hence is trivial by induction.
Thus $G_H^{\rm ab}$ is cyclic, generated by an element of order at most
two. The Frattini quotient of $G_H$ has one generator, so $G_H$ is cyclic and
has order at most two. Since it contains $K_n$, we obtain $H=K_n$.
The $S_n$-class group is a quotient, so its two-part also vanishes.

The signatures of units span all $r_n$ real signs. For $n\ge1$ put
$m=n+2$, $\zeta=\zeta_{2^m}$, and
\[
 u=1+\zeta+\zeta^{-1}
   =\zeta^{-1}\frac{\zeta^3-1}{\zeta-1}.
\]
This is a unit: the numerator and denominator generate the same
dyadic prime ideal in the cyclotomic field. Its real conjugates are
$1+2\cos(2\pi a/2^m)$ for odd $0<a<2^{m-1}$.
The negative ones have $2^m/3<a<2^{m-1}$, and their number is
\[
 2^{m-2}-\left\lfloor\frac{2^{m-1}}3+\frac12\right\rfloor,
\]
an odd integer. Hence $N_{K_n/\mathbf Q}(u)=-1$.
The signature vector of $u$ has augmentation one in the local ring
$\mathbf F_2[\operatorname{Gal}(K_n/\mathbf Q)]$ and is therefore
a unit. Its Galois translates span all signatures. At $n=0$ use $-1$.

Kummer theory now gives
$\dim H^1(G_{K_n,S_n},\mathbf F_2)=r_n+s_n$, and its signature map
to $\mathbf F_2^{r_n}$ is onto. The map in degree zero is the diagonal
injection. The positive complex is the fiber of ordinary global
cochains to the sum of ordinary real cochains, so its long exact
sequence gives $\dim H^1(C_n^+)=r_n+s_n-1$.
The vanishing two-part of the $S_n$-class group identifies global
$H^2$ with $\operatorname{Br}(\mathcal O_{K_n,S_n})[2]$.
The local invariant description has $r_n+s_n$ coordinates with one
sum relation. Projection to the $r_n$ real coordinates is onto,
because $s_n\ge1$, and has kernel of dimension $s_n-1$.
Together with signature surjectivity this gives
\[
 \dim H^1(C_n^+)=2^n+s_n-1,\qquad
 \dim H^2(C_n^+)=s_n-1,
 \qquad H^i(C_n^+)=0\quad(i\ne1,2).
\]
The last assertion uses the usual global-to-real comparison in
degrees greater than two; degree zero vanishes by the diagonal injection.

For each fixed support, $s_n$ is bounded. The dyadic prime is unique;
each fixed odd prime has open decomposition group in the real
cyclotomic $\mathbf Z_2$-extension, since its image in
$\mathbf Z_2^\times/\{\pm1\}$ has nonzero two-adic logarithm.
Use now the finite free model and coefficient change already constructed.
Its scalar reduction is a bounded finite free $\mathbf F_2[[t]]$-complex,
and its finite cyclotomic specializations are derived quotients by
$t^{2^n}$. Smith decomposition shows that torsion disks and their Tor
terms contribute bounded dimensions, while a free cohomology generator
contributes $2^n$. The displayed dimensions therefore give generic
rank one in degree one and zero in all other degrees.

For the present curve, $E[2]$ is the direct sum of two trivial
$\mathbf F_2$-representations. Thus the residual positive complex has
rank two in degree one only. The residue field of the local ring
$B=\mathscr O[G]$ is $\mathbf F_2((t))$; the augmentation ideal of
the finite two-group lies in its maximal ideal. A minimal finite free
$B$-model has zero residual differential, so the just-computed dimensions
force exactly two generators in degree one and no others. Consequently
\[
                    C^+\otimes B\simeq B^2[-1].\tag{E.3}
\]
This uses the local ring $B$, not semisimplicity of $\mathbf F_2[G]$.

The common real sign is needed here. The real cyclotomic character
and $\varepsilon$ are trivial on complex conjugation, so the real
action is the fixed action on $T_2E\otimes\chi_0$, extended to the
group ring. Its invariant module is the fixed free rank-one
$\Z_2$-lattice extended to $\Lambda[G]$, and then to
$\mathscr O[G]$. Thus a generator can be chosen before $G$ varies.
Real $H^1$ is killed by two over these coefficient rings, so the
positive lift below costs the same single factor two for every
binary dimension.
Choose a generator of the fixed real invariant lattice and its
connecting class $a$. Since real $H^1$ is killed by two,
$2^{c_E+1}F_EZ_F$ has an integral positive lift $b$.
For a basis $\theta$ of the original inverse determinant line, define
\[
                         u^{\rm ev}\theta=2^{-c_E}a\wedge b.
 \tag{E.4}
\]
The wedge does not depend on the choice of the positive lift.
Indeed, the difference of two lifts is in the image of real
degree-zero cohomology by the cone exact sequence. That module
is generated by the chosen real invariant generator, so the
difference is a multiple of $a$ and wedges to zero with $a$.
After rationalization set $z=F_EZ_F$ and choose any lift
$\widetilde z$ of $z$. Formula (E.4) then reads
\[
              u^{\rm ev}\theta=2a\wedge\widetilde z.
\]
Thus the factor $2^{c_E}$ used to make the lift integral cancels
exactly. Increasing the fixed clearing exponent does not change
the rational scalar. Character evaluation and derived central
specialization preserve this equality and its cone maps. The
class whose dual exponential occurs in (E.7) is consequently
the same $F_E(Z_F)_h$ as in (K3), and the single remaining factor
two is exactly the $+1$ in (E.7).

\paragraph{From the full cohomology bound to the same determinant scalar.}
Here is the rational divisibility step explicitly. Fix a character
$x$ of $G$, a height-one prime $\mathfrak p$ of $A$, and
$R=A_{\mathfrak p}$. Put $z=F_E(Z_F)_x$, without changing its
normalization. The full global degree-one module is free of rank one
over $R$ by the argument following (K15) and the injectivity of the
maps in (K16). Write it as $Re$, and write $z=\alpha e$.
The higher real cohomology vanishes over $R$, since two is invertible.
Global degree-zero cohomology is zero: an invariant would restrict
to an invariant of the open nonabelian subgroup on the elliptic
representation, on which the scalar coefficient character is trivial.
Thus the positive-cone exact sequence gives
\[
 0\longrightarrow R a\longrightarrow
 H^1(C^+_x\otimes R)\longrightarrow Re\longrightarrow0,
 \qquad
 H^2(C^+_x\otimes R)=H^2(C_{S,x}\otimes R).
\]
The element $a$ here is the image of the fixed real invariant generator
in (E.4), and it is primitive in this positive degree-one module,
because the quotient $Re$ is free. Choose a lift $\widetilde e$ of $e$;
then $(a,\widetilde e)$ is a basis. The equality already proved from
(E.4) becomes
\[
                 u_x^{\rm ev}\theta_x
                  =2\alpha\,a\wedge\widetilde e.
\]
Let $d_1,\ldots,d_r$ be the nonzero Smith entries of a finite free
model of $C^+_x\otimes R$. All its cohomology outside degrees one and
two is zero, its degree-one cohomology is free, and its degree-two
cohomology is torsion. Its inverse determinant lattice, under the
rational identification with $\bigwedge^2H^1$, is therefore generated,
up to an $R$-unit, by
$(\prod_i d_i)a\wedge\widetilde e$. It follows that
\[
 \begin{split}
 v_{\mathfrak p}(u_x^{\rm ev})
 &=v_{\mathfrak p}(\alpha)-\sum_i v_{\mathfrak p}(d_i)\\
 &=\ell_R(H^1(C_{S,x}\otimes R)/Rz)
       -\ell_R H^2(C_{S,x}\otimes R)\ \geq\ 0,
 \end{split}                                                   \tag{E.4a}
\]
where two is a unit and the inequality is exactly (K2).
This uses the full $H^2$ corrected in (K16), rather than only Kato's
$\mathbf H^2_0$.

The intersection of the height-one valuation rings of the normal
domain $A$ gives $u_x^{\rm ev}\in A$ for every $x$. Since $|G|$ is
a power of two, the character idempotents belong to $A[G]$ and these
memberships give $u^{\rm ev}\in A[G]$. This establishes membership
only; it is not a bound for the denominator from character inversion.
The uniform bound comes separately from (E.3): in its free rank-two
$\mathscr O[G]$-module the integral vectors $a,b$ of (E.4) have integral
wedge, so $2^{c_E}u^{\rm ev}\in\mathscr O[G]$ with the fixed exponent
already used for the original lattice. Coefficientwise intersection
in the original group basis gives


\[
              A\cap\mathscr O=\Lambda,\qquad
                  2^{c_E}u^{\rm ev}\in\Lambda[G].                 \tag{E.5}
\]
Indeed an element of $A$ has bounded powers of two as denominators
in its power-series coefficients; membership in $\mathscr O$ makes
their minimum valuation nonnegative. The coefficientwise intersection
retains $c_E$, without a term $v_2|G|$.
\end{proof}

\subsection{Central saturation and the numerical lower bound}
\begin{proposition}[Central calibration of the even scalar]\label{ex:app-central-calibration}
Use the scalar and determinant basis of
Proposition~\ref{ex:app-even-scalar}. At a corank-zero constituent
$E^h$ with finite $\Sha[2^\infty]$, its nonzero central value satisfies
\[
 v_2u_h^{\rm ev}(0)=v_2L_*(h)-2\omega(h)-s(h)+\epsilon_h,
 \qquad |\epsilon_h|\le C_{\rm loc},
\]
with one constant for the fixed local types. At such a constituent
$u_h^{\rm ev}(0)=0$ if and only if $L_*(h)=0$.
Consequently every admissible $h$ with $L_*(h)\ne0$ satisfies
$v_2L_*(h)\ge2\omega(h)+s(h)-c_E-C_{\rm loc}$: the forward
rank-zero theorem in Lemma~\ref{lem:forward-rank} supplies the
preceding corank and finiteness hypotheses in this case. The central saturation corrections are those
in (E.6)--(E.7); no primitivity after inverting $t$ is assumed to
survive central specialization.
\end{proposition}
\begin{proof}\leavevmode
At a corank-zero constituent put
\[
 g=\ell_{\Z_2}E^h(\Q)[2^\infty],\quad
 g_v=\ell_{\Z_2}E^h(\Q_v)[2^\infty],\quad
 r_i=\ell_{\Z_2}H^i(\mathbf R,T_h)\quad(i=1,2).
\]
Let $i$ be the index length of the image of global $H^1$ in the
free rank-one quotient $L_2=H^1(\Q_2,T_h)/H_f^1(\Q_2,T_h)$.
Orthogonal Kummer duality, with its real terms, gives
\[
 \ell H^2_{\rm glob}=s(h)-i+\sum_v g_v+r_2-g.                     \tag{E.6}
\]
Write $K=\ker(H^1_{\rm glob}\to H^1_{\mathbf R})$, let $k$ be its
torsion length and $m$ its free index. If $P=H^1(C_h^+)$, $p_0$ is
its torsion length and $a_s$ the saturation index of the real image
in $P/P_{\rm tors}$, exactness gives
\[
 \ell\operatorname{coker}(H^1_{\rm glob}\to H^1_{\mathbf R})
       =r_1-g+k-m,\qquad k=p_0+a_s.
\]
For the numerical calculation suppose also $L_*(h)\ne0$.
The wedge in (E.4) adds $a_s+1+v_{\rm free}(z)-m$.
Conversion to the inverse determinant adds $p_0-\ell H^2(C_h^+)$.
Substitution of these equalities cancels $m$ and $a_s$ exactly:
\[
 v_2u_h^{\rm ev}(0)
 =v_2\exp^*z-s(h)-\sum_v g_v+2g+1-r_1-e_2,                      \tag{E.7}
\]
where $e_2$ is the valuation of $\exp^*$ on a generator of $L_2$.
Primitivity of the real image after inverting $t$ has not been
specialized to $t=0$.

Here are the real local convention and the exact length sequence used
in (E.6). Put $T=T_2(E^h)$ and $D=E^h[2^\infty]$. Since $E^h[2]$ is
rational, complex conjugation $J$ satisfies $J\equiv1\pmod2$.
The integral projectors $(1\pm J)/2$ give $T=T^+\oplus T^-$.
At the real place use the bounded Kummer complex
\[
 U_{\mathbf R}=\tau_{\le1}R\Gamma(\mathbf R,T)
   =[\,T^+\oplus T^-\xrightarrow{\,J-1\,}T^-\,]
\]
in degrees zero and one, with its map to ordinary real cochains.
Modulo $2^n$ its degree-zero cohomology is
$T^+/2^n\oplus(T^-/2^n)[2]=E^h(\mathbf R)[2^n]$, and its degree-one
image is the real Kummer line $T^-/2$. The latter can also be checked
using a rectangular period lattice: dividing the nonidentity real
component, represented by half the imaginary period, produces the
minus-line cocycle. Thus this line is its exact annihilator under
finite real Tate duality. The ordinary completion
$E^h(\mathbf R)^\wedge_2[-1]$ alone would lose the first summand of
degree-zero cohomology.

At a finite place take the derived Kummer map from
$E^h(\Q_v)^\wedge_2[-1]$. Its reduction has full local degree-zero
cohomology and injective Kummer degree-one map. If $C_{\rm Sel}$
denotes the resulting ordinary global Selmer cone, these assertions
identify
\[
 H^1(C_{\rm Sel}\otimes^{\mathbf L}\Z/2^n)
       =\operatorname{Sel}_{2^n}(E^h/\Q).
\]
At corank zero the discrete Selmer group is finite, the Mordell--Weil
rank is zero, and its order is that of $\Sha[2^\infty]$. The
orthogonal-Kummer specialization of finite Poitou--Tate duality
\cite[5.7.1.4--5.7.1.7]{nekovar}, followed by the compact/discrete
limits, gives
\[
\begin{split}
0\longrightarrow L_2/J_2
 \longrightarrow\operatorname{Sel}_{2^\infty}(E^h/\Q)^\vee
 \longrightarrow H^2_{\rm glob}
 \longrightarrow \bigoplus_{v\in S_f}H^2(\Q_v,T)\oplus H^2(\mathbf R,T)
 \longrightarrow H^0(\Q,D)^\vee\longrightarrow0 ,
\end{split}
\]
where $J_2$ is the image of global $H^1$. There is no real
degree-one quotient: the compact real Kummer image is all of
$H^1(\mathbf R,T)$, and its discrete orthogonal Kummer group is zero.
At finite $v\ne2$ the compact Kummer image is likewise all of local
$H^1$, while at two its quotient is $L_2$. The real degree-two term
is retained. The finite-level images satisfy Mittag--Leffler, and
high-degree localization at the real place terminates the sequence.

This also makes $H^2_{\rm glob}$ finite. The cone sequence then
makes $H^2(C_{\rm Sel})$ finite, and universal coefficients give
\[
 0\to E^h(\Q)/2^n
 \to\operatorname{Sel}_{2^n}(E^h/\Q)
 \to H^2(C_{\rm Sel})[2^n]\to0.
\]
Comparison with the Kummer sequence gives
$\ell H^2(C_{\rm Sel})=s(h)$, with no extra real torsion. Equivalently
the cone sequence is
\[
0\to L_2/J_2\to H^2(C_{\rm Sel})\to H^2_{\rm glob}
 \to \bigoplus_{v\in S_f}H^2(\Q_v,T)\oplus H^2(\mathbf R,T)
 \to H^0(\Q,D)^\vee\to0.
\]
Its alternating lengths are exactly (E.6).

For the positive cone, $H^2_{\rm glob}\to H^2(\mathbf R,T)$ is onto:
any real component's image in $H^0(\Q,D)^\vee$ can be canceled by
a component at two, since restriction
$H^0(\Q,D)\to H^0(\Q_2,D)$ is injective. Thus its exact sequences are
\[
0\to T^+\xrightarrow{a}P\to K\to0,\qquad
0\to\operatorname{coker}(H^1_{\rm glob}\to H^1_{\mathbf R})
 \to H^2(C_h^+)\to H^2_{\rm glob}\to H^2_{\mathbf R}\to0.
\]
If $k$ and $p_0$ are the torsion lengths of $K$ and $P$, $m$ is the
free index of $K$ in global $H^1$, and $a_s$ is the saturation index
of the real line in $P/P_{\rm tors}$, these sequences give
\[
\ell\operatorname{coker}(H^1_{\rm glob}\to H^1_{\mathbf R})
 =r_1-g+k-m,\qquad k=p_0+a_s.
\]
The wedge $2^{-c_E}a\wedge b$ contributes
$a_s+1+v_{\rm free}(z)-m$, and the inverse determinant comparison
contributes $p_0-\ell H^2(C_h^+)$. Substitution leaves
$1+v_{\rm free}(z)+i-s(h)-\sum_vg_v+2g-r_1$.
Finally $v_2\exp^*z=v_{\rm free}(z)+i+e_2$, proving (E.7)
without specializing a primitivity assertion from the punctured
family.
Apply (K3) to the retained class $z=F_E(Z_F)_h$ used in the wedge
above. The primitive-conductor trace and the invariant
$\sqrt h\,e_{\chi_h}$ basis have already been computed in
paragraph~6, including discriminant $4h$. Absorb the fixed
smoothing and old-support factors into $\epsilon_{\rm rec}$.
Then
\[
 v_2\exp^*z
 =v_2L_*(h)+\sum_{q\ {\rm unused}}v_2\#E^h(\mathbf F_q)
                         +\epsilon_{\rm rec},\qquad
 |\epsilon_{\rm rec}|\le C_E.
\]
At unused good primes $g_q=v_2\#E^h(\mathbf F_q)$, so these
terms cancel. At an active good prime the inertia invariants in
$E^h[2^\infty]$ are $E[2]$, of length two for the present curve.
The remaining local types have fixed bounds, including global
torsion by its injection at two.
Thus
\[
 v_2u_h^{\rm ev}(0)=v_2L_*(h)-2\omega(h)-s(h)+\epsilon_h,\qquad
 |\epsilon_h|\le C_{\rm loc}.                                  \tag{E.8}
\]
The central rational positive complex has only rank-two $H^1$, so
this is specialization of (E.4). Corank-zero duality also shows
$u_h^{\rm ev}(0)=0$ if and only if $L_*(h)=0$.
At $L_*(h)\ne0$, Lemma~\ref{lem:forward-rank} supplies the
corank-zero and finite-$\Sha$ hypotheses just used.
Equations (E.5) and (E.8), already with $G=1$, yield
\[
          v_2L_*(h)\ge2\omega(h)+s(h)-c_E-C_{\rm loc}.             \tag{E.9}
\]
No binary-family dimension is needed for this lower bound.
\end{proof}

% Replaces baseline manuscript.tex lines 3416--3510 inclusive.
% Baseline SHA256 525c8c05e7a5acc83d6b1f89c1056fd1321bd7490fd227a63d1a520e941b5197.
\subsection{The actual even form and its phase}
\begin{proposition}[An actual even form with rational normalized phase]\label{ex:app-even-phase}
Fix a sign $\sigma$, unit squareclasses at a finite support $S_0$
containing $2N$, and a positive squarefree reference $D_0$ in that
filter with $L_*(\sigma D_0)\ne0$. For every fixed finite extra
exclusion set $B$ prime to $D_0$, there is an algebraic holomorphic
cuspidal form of weight $3/2$, normalized by $a(D_0)=1$, with the
scalar characters $\psi,\eta$ constructed below and
$\psi/\eta=\chi_{-\sigma}$. It has the prescribed coefficient
support. For every retained positive squarefree $D$ in the original
$S_0$ filter,
\[
 a(D)=z_Dr_D,\qquad z_D\in\mu_{2^\infty},\quad r_D\in\Q,
 \qquad r_D^2=L_*(\sigma D)/L_*(\sigma D_0).
\]
It satisfies the all-index good-prime recurrence (E.12) outside
$S_0\cup B$. Nonzero retained squarefree coefficients have integral
$v_2a(D)$, and $v_2a(D)-\omega(D)$ has a finite attained minimum.
The form and the initial sieve are fixed before this minimum is
chosen. Additional exclusions do not authorize a new minimum.
\end{proposition}
\begin{proof}\leavevmode
Fix a sign $\sigma$ and unit squareclasses at a finite support $S_0$
containing two and every prime dividing $N$. Let $D_0>0$ be squarefree
in these prescribed classes, with $L_*(\sigma D_0)\ne0$.
In (E.10)--(E.11), $D$ ranges over positive squarefree parameters in
these same unit squareclasses at every prime of $S_0$.
Fix also a finite extra exclusion set $B$ with
$\gcd(D_0,\prod_{p\in B}p)=1$.

\paragraph{A nonzero reference functional and its scalar local type.}
Use the global Shimura--Waldspurger packet with underlying
$\operatorname{PGL}_2$ representation $E^{(\sigma)}$ in the
normalization of Waldspurger's Corollary~2. The nonzero central value
at $\sigma D_0$ supplies the member with a nonzero $D_0$-Whittaker
functional. Equivalently, take the nonzero theta lift associated with
$E^{(\sigma D_0)}$ and the additive character of frequency $D_0$;
the change of additive character identifies its underlying packet
with the preceding one. The nonvanishing and the local identification
are the content of Waldspurger's global theta theorem and its
Fourier-functional calculation
\cite[Theorem~1 and Proposition~1]{waldspurgerICM}
and \cite[Proposition~18 and its proof, pp.~466--467]{waldspurger}. The positive real frequency selects the
holomorphic member of lowest weight $3/2$. The lift is in the
unary-orthogonal cuspidal spectrum.

Here is a local construction of a scalar type that retains the
reference functional. At $p\in S_0$, write $\ell=\ell_{D_0}$ and
choose $v$ with $\ell(v)\ne0$. Normalize Haar measure in
\[
 P_Rv=\frac1{\operatorname{vol}(p^{-R}\Z_p)}
       \int_{p^{-R}\Z_p}
          e_p(-D_0b)\pi(n(b))v\,db .
\]
Then $\ell(P_Rv)=\ell(v)$. Let $M(u)$ be the unit action normalized
by the Weil multiplier. Explicitly, if the chosen torus lifts have
cocycle $(u,v)_p$, their theta-cubed multiplier $\gamma$ has
two-power root-of-unity values and satisfies
$\gamma(u)\gamma(v)=(u,v)_p\gamma(uv)$; put
$M(u)=\gamma(u)^{-1}\pi(\widetilde m(u))$.
It is therefore an action of
$\Z_p^\times$ and
$M(u)\pi(n(b))M(u)^{-1}=\pi(n(u^2b))$.
Changing $u$ to its inverse changes none of the following support
conditions. Direct substitution gives
\[
 \ell(M(u)P_Rv)=
 {\bf1}_{\,u^2\equiv1\;(\mathrm{mod}\;p^R)}
                    \ell(M(u)v).
\]
The integral is a finite sum on this smooth vector.

The normalized $M(-1)$ acts by a scalar $\delta\in\{1,-1\}$.
Choose a continuous unit character $\alpha$ of two-power order
with $\alpha(-1)=\delta$. Such a character exists: for odd $p$,
$-1$ lies in the two-primary factor of the finite cyclic group
of residue-field units; for $p=2$ use the factor
$\{1,-1\}$ of $\Z_2^\times$.
Take $R$ large enough that both $v$ and $\alpha$ are constant on
the small neighborhoods appearing above. The projection
\[
 w=\int_{\Z_p^\times}\alpha(u)^{-1}M(u)P_Rv\,d^\times u
\]
satisfies $M(u)w=\alpha(u)w$. Its $\ell$-value is nonzero:
the neighborhoods of $1$ and $-1$ contribute the same nonzero
value. For odd $p$ the solution set is
$1+p^R\Z_p$ together with $-1+p^R\Z_p$.
For $p=2$ and $R\ge3$ it is
\[
 (1+2^{R-1}\Z_2)\ \cup\
 (-1+2^{R-1}\Z_2).
\]
Thus the four residue classes modulo $2^R$ form two neighborhoods;
there are not four independently assignable phases.
For completeness, the noncancellation in this projection is an exact
calculation. Normalize multiplicative Haar measure to give
$\Z_p^\times$ volume one, and put
$K_R=1+p^R\Z_p$ for odd $p$, and $K_R=1+2^{R-1}\Z_2$
for $p=2$. Increase $R$ until $M(k)v=v$ and $\alpha(k)=1$
for $k\in K_R$. On $K_R$ the integrand after applying $\ell$
is $\ell(v)$. On $-K_R$ it is
$\alpha(-1)^{-1}\ell(M(-1)v)=\delta^{-1}\delta\ell(v)$.
It vanishes elsewhere by the displayed additive projection formula.
Consequently
\[
                 \ell(w)=2\operatorname{vol}(K_R)\ell(v)\ne0.
\]
This calculation takes place in characteristic zero; reduction at two
is not involved. For the later support projection choose $m$ so that
the required unit squareclass is a union $A$ of residue classes modulo
$p^m$. Its operator is
\[
 \Pi_A=\frac1{p^m}\sum_{b\bmod p^m}
       \left(\sum_{a\in A}e_p(-ab/p^m)\right)\pi(n(b/p^m)).
\]
On the Fourier functional of integral frequency $D$ it acts as
$\mathbf1_A(D)$. It therefore preserves the reference value and
annihilates every excluded frequency. Conjugation by $M(u)$ permutes
the summands and replaces $A$ by $u^2A=A$, proving the asserted
commutation. Its denominators are fixed before the form is normalized.


The vector $P_Rv$, and then $w$, is fixed by upper unipotents
$n(\Z_p)$, because $D_0$ is a unit. Smoothness makes $w$
fixed by lower unipotents at sufficiently deep level. Together with
the unit eigencharacter, the elementary decomposition of
$K_0(p^m)$ gives a scalar $K_0$ type for sufficiently large $m$.
The Weil normalization contributes only two-power roots of unity,
so the resulting classical character has two-power order.
The forced values at $-1$ give the global parity required for a
weight-$3/2$ scalar character: the product of the finite central
actions is the inverse of the holomorphic archimedean action,
and the theta-cubed multiplier has the same product relation
on rational $-I$. Thus the quotient character satisfies
$\psi(-1)=1$. At $2$ one may, if desired, choose
$\alpha$ of arbitrarily large two-power conductor before increasing
$R$; this does not enlarge the prime support.

Apply the finite additive Fourier idempotent for the prescribed
unit-squareclass subset modulo a sufficiently high power of $p$.
This idempotent keeps the $D_0$ coefficient. Its support condition
is stable under multiplication by unit squares, so it commutes
with the unit action and retains a scalar type, at a possibly
deeper level. At every prime outside $S_0$ take the spherical vector.
The packet there is unramified; its standard squarefree coefficient
$c_p^0$ is $1$ at both valuations zero and one
\cite[Section~I.3, pp.~377--378]{waldspurger}.
This includes the primes dividing $D_0$. Local Whittaker uniqueness
and the tensor-product factorization show that the resulting
global reference coefficient is still nonzero.

\paragraph{The exact full-space comparison.}
Let the resulting scalar character be $\psi$ and put
$\eta=\psi\chi_{-\sigma}$. Both have two-power order.
Let $V$ be the full good-Hecke eigenspace of this scalar character,
at the level just constructed. Keep all old multiplicities. In
Waldspurger's notation its integral-weight eigenform and character are
\[
 \phi_0=f_E\otimes\eta,\qquad
 \chi_0=\psi\chi_{-1},\qquad
 \phi_0\otimes\chi_0^{-1}=f_E\otimes\chi_\sigma .
\]
This also fixes the signs in the packet convention: Waldspurger's
positive index $D$ twists the last representation by $\chi_D$.
The Hecke eigenvalues in $V$ are $\eta(p)a_p(E)$.

The nonzero space implies hypothesis (H1) by Waldspurger's
Proposition~2 \cite[Section~I.2, p.~377]{waldspurger}. In explicit terms, at a principal-series
place its two inducing characters take value $1$ at $-1$.
Corollary~2 \cite[pp.~483--484]{waldspurger} applies to every $f\in V$ and any
positive squarefree pair whose ratio is a square in every
$\Q_p$ at a prime dividing the level. Our $D,D_0$ satisfy
this condition because the level has prime support $S_0$.
The final paragraph of that proof removes (H2) for this
same-$\Q_2$-squareclass comparison. No assumption of
non-supercuspidality at two is required.

\paragraph{The printed coefficient identity and its dyadic hypothesis.}
For clarity, write $\mathcal L_D=L(E^{(\sigma D)},1)$ and let
$N_{\rm W}$ denote the level in Waldspurger's notation, not the
conductor of $E$. In \cite[Corollary~2, pp.~483--484]{waldspurger}
take
\[
 k=3,\quad \chi=\psi,\quad \chi_0=\psi\chi_{-1},\quad
 \rho_0=\rho_E\otimes\eta,\quad n_1=D_0,\quad n_2=D.
\]
Then $\rho_0\otimes\chi_0^{-1}=\rho_E\otimes\chi_\sigma$.
The corollary, in its printed order, is
\[
 a_f(D_0)^2\mathcal L_D\,\psi(D/D_0)D^{1/2}
   =a_f(D)^2\mathcal L_{D_0}D_0^{1/2}.
\]
The passage from the automorphic central argument $1/2$ to the
classical argument $1$ is the same for both terms. Dividing only
by the nonzero reference value gives (E.10). In particular this
is a literal square identity, with neither complex conjugation
nor an unspecified absolute-value normalization.

Every prime of $N_{\rm W}$ belongs to $S_0$ at this point of the
construction. The prescribed local classes make $D/D_0$ a square
in $\mathbb Q_p^\times$ at each such prime, including $p=2$.
In the proof of the corollary, the local coefficient functions
at these primes have equal values at $D,D_0$. Outside the level,
the squarefree arguments have valuation zero or one and the
unramified coefficient function has value one in both cases.
The real coefficient contributes exactly the power
$(D/D_0)^{(k-2)/4}$ before squaring. The quotient of epsilon
factors is computed on pp.~483--484: its Hilbert-symbol product
disappears by reciprocity and its remaining character is
$\chi(D/D_0)=\psi(D/D_0)$. Thus these local terms do not leave
an additional prime-dependent factor in the displayed identity.

The final paragraph on p.~484 explicitly explains the use of
(H2): it is needed in that argument only to compare Fourier
indices in different squareclasses of $\mathbb Q_2^\times$.
Our pair is in one squareclass. Hence the corollary as stated
requires (H1) here, but does not impose (H2), a nonsupercuspidal
condition at two, or a new dyadic conductor restriction. The
nonzero scalar eigenspace supplies (H1) by
\cite[Proposition~2, p.~377]{waldspurger}. Moreover the proof on
p.~483 extends its linear relation to every vector in the full
eigenspace before squaring. The resulting identity therefore
applies simultaneously to old multiplicities, and later to both
Galois-conjugated and character-twisted vectors at a common level.
At that later level the same squareclass test is repeated; its
prime support is still $S_0$.
Putting $n_1=D_0,n_2=D$ in the corollary, with weight $3/2$, gives
\[
 a_f(D)^2=
 \psi(D/D_0)\,
 \frac{\sqrt D\,L(E^{(\sigma D)},1)}
      {\sqrt{D_0}\,L(E^{(\sigma D_0)},1)}
 a_f(D_0)^2.                                      \tag{E.10}
\]
The period cancels, so the quotient in this formula is the quotient
of the specified $L_*$-values. No product over the primes dividing
$D$ or $D_0$ is missing: the unramified factors at squarefree
valuations are exactly those checked above.

For clarity, (E.10) gives proportional coefficient functionals on
the whole space. If $\ell_0(f)=a_f(D_0)$ and $\ell_D(f)=a_f(D)$,
then $\ell_D^2=C\ell_0^2$ as a polynomial identity on $V$.
If $C=0$, then $\ell_D=0$. Otherwise, choose a square root $c$ of
$C$. The polynomial ring on the finite-dimensional vector space
$V$ is an integral domain, so
$(\ell_D-c\ell_0)(\ell_D+c\ell_0)=0$ implies
$\ell_D=c\ell_0$ or $\ell_D=-c\ell_0$.
The scalar is common to every vector. This argument does not
assume that $V$ is one-dimensional.

\paragraph{An algebraic vector and the Galois comparison.}
Ramsey's section model \cite[Definition~3.1 and Theorem~4.2]{ramsey} realizes the half-integral-weight space as
sections of the indicated divisor sheaf on $X_1(4N_{\rm aux})$.
Here $N_{\rm aux}$ is supported on $S_0$ and includes the chosen
local depths; it need not equal the conductor $N$ of $E$.
Over characteristic-zero fields, flat field extension commutes
with these sections. Cuspidality, scalar character, support
conditions and prescribed good-Hecke eigenvalues are algebraic
linear conditions. Although infinitely many coefficients or
primes can occur in their description, a finite subset cuts out
the same subspace, by finite dimensionality.

One need not impose Petersson orthogonality as an additional
transcendental linear condition. The elliptic good-Hecke eigenspace
is automatically orthogonal to unary theta series: at a good prime,
the unary eigenvalues have absolute value $p+1$, whereas
$|\eta(p)a_p(E)|\le2\sqrt p<p+1$. The unary space and its
orthogonal complement are stable under the Hecke operators and
their adjoints. Thus their elliptic eigenspaces are disjoint.
The nonzero reference functional on the complex space is consequently
nonzero on its algebraic model. Choose such an algebraic $g$ and
scale it to $a(D_0)=1$.
This uses flat field extension only, not Ramsey's arbitrary-ring
base-change assertion in weight $3/2$; see \cite[Remark~4.4]{ramsey}.
Here the field-extension assertion can be checked without a
cohomology-vanishing theorem. Work over a number field $K$ containing
the character values and the roots of unity specifying the cusps and
the support idempotents. A finite affine cover of the modular curve
computes global sections as the kernel of the difference of restriction
maps to pairwise intersections. Tensoring with a field extension is
flat and preserves that kernel. Thus the sections in Ramsey's
Definition~3.1 commute with $K\subset\overline\Q\subset\mathbf C$,
also in weight $3/2$. Impose vanishing at the cusps, the diamond
character, the support idempotents, and the kernels of
$T(p^2)-\eta(p)a_p(E)$ inside this finite-dimensional space.
The span in its dual of all the defining linear equations has a finite
basis. These conditions therefore commute with the same field
extensions. Since the reference functional is nonzero after extension
to $\mathbf C$, it is nonzero before extension; a vector in a $K$-basis has
nonzero reference value and may be divided by that value.
This also places every coefficient of the chosen form in one number
field. No assertion about base change to a residue field, nor about
vanishing of $H^1$ in weight $3/2$, is used.


Let $\tau\in\operatorname{Aut}(\overline{\Q}/\Q)$.
Choose an odd integer $e$ describing its action on all values of
$\eta$ and $\psi$, and put $\nu=\eta^{(e-1)/2}$.
Coefficientwise conjugation gives a form $g^\tau$ with character
$\psi^e$ and good eigenvalues $\eta(p)^e a_p(E)$.
One can reduce this conjugation assertion to integral weight by
multiplying by the ordinary theta series: it has rational
coefficients, and its fourth power has trivial scalar character.
The algebraic divisibility condition at the cusps is precisely the
one in Ramsey's section model, so division after conjugation is
again a half-integral-weight form.
The coefficientwise twist $g\otimes\nu$ has character
$\psi\nu^2=\psi^e$ and exactly these good eigenvalues.
The identities hold because $\psi/\eta$ is quadratic.
The twist is modular at a level still supported on $S_0$; this
follows as well by writing it as the finite additive translation
sum using the Gauss sum of the primitive character inducing $\nu$.

Pass to a common such level. Both forms lie in its full
good-Hecke eigenspace, so the proportionality of its $D$ and $D_0$
functionals gives
\[
 \tau(a(D))
 =\frac{\nu(D)a(D)}{\nu(D_0)a(D_0)}
 =\nu(D/D_0)a(D).                                  
\]
There is no vanishing reference denominator: $D_0$ is prime to
$S_0$ and $a(D_0)=1$. No rationality of the central-value quotient
has been assumed in obtaining this equality.

Write the Dirichlet character as a product of its finite
prime-power factors $\eta_p$. For every $p$ dividing its conductor
choose $u_p\in\Z_p^\times$ with $u_p^2=D/D_0$, and set
\[
                 z_D=\prod_{p\mid\operatorname{cond}(\eta)}
                                      \eta_p(u_p).
\]
Then $z_D$ is a two-power root of unity,
$z_D^2=\eta(D/D_0)$, and
\[
 \tau(z_D)=z_D^e,\qquad
 \nu(D/D_0)=z_D^{e-1}.
\]
Changing the selected odd representative of $e$ has no effect on
the last identity, because $D/D_0$ is a local square at the
conductor. The preceding Galois identity now shows that $a(D)/z_D$ is fixed by
every $\tau$, hence lies in $\Q$.
Different choices of the local square roots change $z_D$ by a
sign, which changes only the sign of this rational number.

Finally, $\psi/\eta=\chi_{-\sigma}$ is trivial on $D/D_0$:
its conductor is supported at two and the ratio is a square there.
Thus $\psi(D/D_0)=z_D^2$. Dividing (E.10) by $z_D^2$ gives
\[
 a(D)=z_Dr_D,\qquad z_D\in\mu_{2^\infty},\qquad r_D\in\Q,
 \qquad r_D^2=\frac{L_*(\sigma D)}{L_*(\sigma D_0)}.       \tag{E.11}
\]
The argument also covers the zero coefficient case. In particular,
every nonzero coefficient in (E.11) has integral normalized
two-adic valuation.

\paragraph{Recurrence, later exclusions, and valuations.}
Only now impose the extra finite coprimality sieve $B$.
The zero-or-retain Fourier sieve is a finite additive translation
operator. It preserves cuspidality, keeps every coefficient whose
index is prime to $B$, and commutes with good $T(p^2)$ for
$p\notin S_1=S_0\cup B$. In particular it keeps $a(D_0)=1$.
The squarefree comparison (E.11) was proved before this enlargement,
so it requires no extra squareclass hypotheses at the primes of $B$.
After the sieve it is asserted at the retained indices.
The standard good $T(p^2)$ formula, with $\psi^2=\eta^2$, gives
\[
 a(p^2n)=\{\eta(p)a_p(E)-\psi(p)(-n/p)\}a(n)
                       -p\eta(p)^2a(n/p^2),\qquad p\notin S_1.
                                                        \tag{E.12}
\]
For squarefree $D$, write $a(Dp^{2r})=\eta(p)^r b_r$.
The first multiplier is
$a_p(E)-(\psi/\eta)(p)(-D/p)$; subsequent multipliers satisfy
$b_{r+1}=a_p(E)b_r-pb_{r-1}$. All are integers, including when
$p\mid D$. Iteration over the good primes gives the same assertion
for every retained square shell after removal of the $\eta$ phase.
More explicitly, for $p\notin S_1$ set
$A_{p,0}(D)=1$, $A_{p,1}(D)=a_p(E)-(\psi/\eta)(p)(-D/p)$,
and
$A_{p,r+1}(D)=a_p(E)A_{p,r}(D)-pA_{p,r-1}(D)$ for $r\ge1$.
Then, for every retained index $Dm^2$ with $D$ squarefree,
\[
 a(Dm^2)=z_D\eta(m)r_D
              \prod_{p^r\Vert m}A_{p,r}(D).
\]
The formula includes $p\mid D$, when the initial Legendre symbol
is zero. When applying the recurrence at one prime, the other
prime-square factors do not alter that Legendre symbol. This proves
the formula by induction on the prime factors of $m$, with
$a(n/p^2)=0$ when $p^2\nmid n$. All factors $A_{p,r}(D)$ are
integers. Thus every coefficient used by the extraction and shell
operators of Proposition~\ref{ex:app-coefficient-rules} is a rational
number times a two-power root of unity, including coefficients with
non-squarefree index. In the chosen residue field such a root reduces
to one. A rational coefficient of valuation zero reduces to the
element $1$ of its prime field $\F_2$. Consequently the fresh
coefficient-one reading of that proposition detects precisely equality
in the fixed minimum-depth bound. It uses the same $g$, $D_0$ and
normalization throughout.


Combining (E.9) and (E.11) makes the values
$v_2a(D)-\omega(D)$ at nonzero retained squarefree coefficients a
bounded-below subset of $\Z$. It is nonempty because $D_0$ is retained,
and therefore has an attained finite minimum. The form and the initial
finite sieve are fixed before this minimum is chosen. Restricting
later fresh primes by further exclusions does not authorize choosing
another minimum.
\end{proof}

% End 07-even-source.tex


% Begin 08-odd-source.tex
\section{Actual odd theta forms and their operators}\label{app:odd-source}
We now construct the forms in (3.8), with the same $h_*$, $D_w$, $L$
and actual Kummer stages as in Section~\ref{sec:rc}.
Write $m=|h_*|$. An admissible $D$ is positive, squarefree, prime to
the fixed support and in the prescribed unit classes. Then
$K_D=\Q(\sqrt{Dh_*})$ is imaginary, split at $2N$, has no extra
units, and $Dh_*$ is an odd fundamental discriminant.
The conductor-one point is the unaveraged sum
\[
 P_D=\sum_{[\mathfrak a]\in\operatorname{Pic}(\mathcal O_{K_D})}
       \chi_{h_*}(\Norm\mathfrak a)\,
                         \pi(x_D^{[\mathfrak a]}).               \tag{O.1}
\]
It is distinct from the conductor-$|h|$ ring-class point used in
Section~\ref{sec:limits}.

% Replace immutable baseline lines 3532--3591 inclusive.
% Keeps (O.2). This expands the existing non-CM argument.
\subsection{A fixed witness and one compatible sequence}
\begin{proposition}[One witness and a compatible detecting sequence]\label{ex:app-detector}
For the fixed $E=E_\ell$ under the full-corank-one hypothesis of
Theorem~\ref{hyp:RC}, use the prepared positive unit filter containing
$1$. There exist a negative odd fundamental discriminant $h_*$ with
$L(E^{h_*},1)\ne0$, a positive fundamental discriminant $D_w$ in
that filter with $L'(E^{D_w},1)\ne0$, and a fixed integer $L$ such
that $LP_D$ belongs to the $D$-twist rational line for every
admissible $D$ and $Q=LP_{D_w}$ is nontorsion.
There is one sequence of actual stages with $M\to\infty$, good
primes $\ell_0$, and additive maps
\[
 \lambda_M:E(\F_{\ell_0^2})\longrightarrow\Z/2^M\Z,
 \qquad v_2\lambda_M(\overline Q)=c_w<M,
\]
where $c_w$ is fixed. At these stages Frobenius squared is trivial
on $E[2^M]$, $(D_w/\ell_0)=1$, $(h_*/\ell_0)=-1$, and
$\lambda_M$ has Frobenius sign $+1$.
Any specified finite prime exclusion is permitted. Additional
field prescriptions require a compatible common lift; no preservation
of $c_w$ under arbitrary later field prescriptions is asserted.
\end{proposition}
\begin{proof}\leavevmode
We use the positive unit-squareclass filter containing one, with its
fixed support enlarged before choosing the witness. The rational
two-isogeny parity theorem \cite[Theorem 1.8]{dd} gives root number
$-1$ for $E$ under the full-corank-one assumption. Apply clause (ii)
of \cite[pp.~543--544]{bfh} to its weight-two newform, with every
old prime in the splitting set, to obtain a negative fundamental
discriminant $h_*$ with $L(E^{(h_*)},1)\ne0$. Including two makes
$h_*$ odd and excludes the extra-unit fields. Now add the primes
of $h_*$ to that set and apply clause (i) to obtain a positive
fundamental discriminant $D_w$ with $L'(E^{(D_w)},1)\ne0$.
Splitting at every old prime puts $D_w$ in the filter containing one
and makes it prime to $h_*$. These choices precede every detecting
prime and every theta precision.

The conductor-one genus character for
$K_D=\Q(\sqrt{Dh_*})$ cuts out
$F_D=\Q(\sqrt D,\sqrt{h_*})$. Thus the unaveraged sum $P_D$ in (O.1)
belongs to $E(F_D)$. Its rational character components are the
$D$- and $h_*$-components. In the degenerate case $D=1$ these are
the trivial and $h_*$ characters of the actual quadratic group.
The $h_*$-component has rank zero by the forward rank-zero theorem.
For $D_w$, the unaveraged height formula
\cite[Theorem 1.1]{cst} has Rankin derivative
$L'(E^{(D_w)},1)L(E^{(h_*)},1)\ne0$, so $P_{D_w}$ is nontorsion.
Here the conductor is one, all primes of $N$ split, and the
extra-unit fields have been excluded.

We make the fixed multiplier explicit. Choose distinct good rational
primes $p_1,p_2$ for $E$ and put
\[
 B_j=\prod_{f=1}^{4}\#E(\F_{p_j^f}),\qquad B=B_1B_2.
\]
For every number field of degree at most four, prime-to-$p_j$
torsion injects into a reduction group with residue degree at most
four. The two primes therefore show that $B$ kills its entire
torsion subgroup. Put $B_{\rm cusp}=C_\pi$, where $C_\pi$
annihilates the image of the cuspidal subgroup defined after (4.1).
This includes every Atkin--Lehner translation $\pi(W\infty)$,
which need not vanish merely because $\pi(\infty)=0$. For $\Delta_D=\operatorname{Gal}(F_D/\Q)$,
the integral character operator
\[
 4e_D=\frac4{|\Delta_D|}
                  \sum_{\gamma\in\Delta_D}\chi_D(\gamma)\gamma
\]
is defined also when $D=1$. The difference
$4P_D-4e_DP_D$ is torsion, since its only possible rational
component is the rank-zero $h_*$-component. Hence
$4BP_D=B(4e_DP_D)$ transforms exactly by $\chi_D$.
Choose $L$ to be a common multiple of the previously required fixed
multiplier and $4BB_{\rm cusp}$, for example their product. This
retains every earlier descent and local-condition factor.
Then $LP_D$ is a rational point of
$E^{(D)}$ under the standard twist identification, for every
admissible $D$. In particular
\[
                         Q=LP_{D_w}\in A(\Q),
                 \qquad A=E^{(D_w)},
\]
is nontorsion. No assertion that the original $P_D$ is
$E(K_D)$-rational is required.

Choose compatible division points $Q_n$ with
$2^nQ_n=Q$ and $2Q_{n+1}=Q_n$, and let
$c_Q(g)=(gQ_n-Q_n)_n\in T_2A$ be the Kummer cocycle.
Write $\rho_A$ for the Tate representation and $H=\ker\rho_A$.
On $H$, this cocycle is a homomorphism. Its image is compact,
and hence a closed $\Z_2$-submodule of $T_2A$. Moreover
\[
 c_Q(ghg^{-1})=\rho_A(g)c_Q(h)\qquad(h\in H),
\]
so its rational span is invariant under the Tate image.

The curve $A$ is non-CM. Open image \cite{serre} supplies both
irreducibility of its rational Tate representation and a central
scalar $z\ne1$ in its Tate image. If $c_Q(H)=0$, the rational
cocycle descends to that image. Commuting $z$ with $g$ gives
\[
 (z-1)c_Q(g)=(g-1)c_Q(z).
\]
Since $z-1$ is invertible over $\Q_2$, this is a coboundary.
It contradicts the Kummer injection
$A(\Q)\otimes\Q_2\hookrightarrow H^1(\Q,V_2A)$ and
nontorsion of $Q$. Thus the translation image is nonzero.
Irreducibility makes its rational span all of $V_2A$.
It consequently contains $2^{a_0}T_2A$ for some $a_0$.

Let $F_0/\Q$ be any fixed finite Galois auxiliary extension,
including $\Q(\sqrt{D_w},\sqrt{h_*})$ and any fixed ring-class
data to be retained with the action of complex conjugation.
Put $H_0=H\cap G_{F_0}$. It has finite index $d$ in $H$.
For $h\in H$, $h^d\in H_0$ and
$c_Q(h^d)=d\,c_Q(h)$. Hence $c_Q(H_0)$ still contains
$2^aT_2A$ for some fixed $a$.
This argument uses no independence between $F_0$ and the Tate
or division-point fields.

Fix actual complex conjugation $\tau$. Since $D_w>0$, the real
twist identification identifies its matrix on $T_2A$ with its
matrix $C$ on $T=T_2E$. We have $C^2=1$ and $\det C=-1$.
Its dual $+1$-eigenspace therefore contains a primitive integral
row $b:T\to\Z_2$, with $bC=b$ and $b(T)=\Z_2$.
Choose $t_0\in T$ with $b(t_0)=1$, and choose $h\in H_0$ with
$c_Q(h)=v=2^at_0$. Set $\sigma=h\tau$. If $d_\tau=c_Q(\tau)$,
then $(1+C)d_\tau=0$, because $\tau^2=1$. Consequently
\[
 \rho_A(\sigma)=C,\qquad
 c_Q(\sigma^2)=(1+C)(v+d_\tau)=(1+C)v,\qquad
 b\,c_Q(\sigma^2)=2^{a+1}.
\]
Put $c_w=a+1$. The element $\sigma$ and this finite valuation
are fixed before precision varies. Its restriction to $F_0$
is that of $\tau$: it fixes $\sqrt{D_w}$ and negates
$\sqrt{h_*}$. Since $h$ fixes $\sqrt{D_w}$, the matrix of
$\sigma$ on $T_2E$ is also $C$.

For each $M>c_w$, take the finite Galois extension generated by
$F_0$, $E[2^M]$, $A[2^M]$, and all $2^M$-division points of $Q$.
Restrict this same actual $\sigma$ to the extension. Finite
Chebotarev supplies infinitely many good rational primes $\ell_0$
in its conjugacy class. Choose a place realizing this representative,
and exclude two, all bad primes, the support of $D_wh_*$ and any
other finite set. Thus, with the torsion basis transported by this
place,
\[
 \operatorname{Fr}_{\ell_0}\equiv C\pmod{2^M},\qquad
 (D_w/\ell_0)=1,\qquad(h_*/\ell_0)=-1.
\]
The determinant also gives $\ell_0\equiv-1\pmod{2^M}$.
In particular $\operatorname{Fr}_{\ell_0}^2$ acts trivially on
$E[2^M]$.

Here is the finite-field map without a normalization ambiguity.
Put $n=2^M$ and write $F$ for Frobenius on the reduction of $E$.
For $P\in E(\F_{\ell_0^2})$, choose $R$ with $nR=P$ and set
\[
                    \kappa_M(P)=F^2R-R\in E[n]=T/nT.
\]
Changing $R$ by $n$-torsion does not change this value.
The map is additive, with kernel $nE(\F_{\ell_0^2})$.
It is onto: $F^2-1$ is a nonzero separable isogeny, so for
$S\in E[n]$ one can solve $(F^2-1)R=S$ and take $P=nR$.
Also $\kappa_M(FP)=F\kappa_M(P)$, by choosing $FR$ as a
division point. Therefore $\lambda_M=b\kappa_M$ has
Frobenius sign $+1$.

Good reduction identifies the reduction of the global division
cocycle with this exact map. The chosen place makes the twist
identification defined over $\F_{\ell_0}$, since
$(D_w/\ell_0)=1$. Thus the fixed witness satisfies
\[
 \lambda_M:E(\F_{\ell_0^2})\longrightarrow\Z/2^M\Z,\qquad
 v_2\lambda_M(\overline Q)=c_w<M.                         \tag{O.2}
\]
Choosing increasing $M$ and excluding previous primes gives one
sequence of actual stages. Subsequent finite prime exclusions
do not alter $c_w$. New field prescriptions require a compatible
common lift and, if added after this choice, are not asserted to
preserve this particular constant.
\end{proof}

% Completion audit replacement for baseline lines 3593--3748.
\subsection{Local functions and the theta sum}
\begin{proposition}[The actual odd theta form and its recurrence]\label{ex:app-theta-form}
Fix an actual stage of Proposition~\ref{ex:app-detector} with $M\ge2$.
Use $m=|h_*|$, the same $L$, the Eichler order of level $N$ in the
quaternion algebra ramified at $\ell_0,\infty$, and the local
Schwartz functions specified below, including (O.3).
With $c_0=|\operatorname{GL}_4(\F_3)|=24\,261\,120$, the sum
$H_M$ in (O.4) is well defined modulo $2^M$ in the character ring.
It gives an actual holomorphic form
\[
 g_M=\sum_n a_M(n)q^n,\qquad a_M(n)=H_M(mn),
 \qquad g_M\in M_{3/2}(\Gamma_0(K),\psi),
\]
with level supported only at the specified coefficient primes.
The characters are exactly
$\psi=\eta_B(m/\cdot)$ and $\eta=\eta_B\chi_{h_*}$, so
$\psi/\eta=\chi_{-1}$. For every prime outside that support,
the recurrence (O.6) holds at every integer index, with nonintegral
indices interpreted as zero. Changing the lifts of $P_M$ loses no
precision. A specified further local coefficient sieve retains this
conclusion outside the enlarged support.
\end{proposition}
\begin{proof}\leavevmode
In this appendix $M=M_\theta$; take $M\geq2$.  Retain the
same $m=|h_*|$, $L$, and detecting prime $\ell_0$ as in (O.2).
Since $\operatorname{Fr}_{\ell_0}^2=1$ on $E[2^M]$, evaluation at
arithmetic Frobenius squared identifies the finite-field Kummer
quotient with $T_2E/2^M T_2E$.  Equivariance and $bC=b$ imply
$\lambda_M(\operatorname{Fr}_{\ell_0} P)=\lambda_M(P)$.
The determinant gives $\ell_0\equiv-1\pmod4$; hence
$(m/\ell_0)=(-1/\ell_0)(h_*/\ell_0)=1$.
\paragraph{The finite-field evaluation map.}
Write $k_0=\mathbb F_{\ell_0^2}$ and let $F$ denote its arithmetic
Frobenius. For $P\in E(k_0)$, choose $Q_P\in E(\overline{k_0})$
with $2^M Q_P=P$. Then
\[
 \kappa_M(P)=FQ_P-Q_P\in E[2^M]
\]
is independent of the choice of $Q_P$, since $F$ is the identity
on $E[2^M]$. This is an additive map. It kills $2^ME(k_0)$,
and its kernel is exactly that subgroup: a zero value means that
$Q_P$ itself lies in $E(k_0)$. The finite group
$E(k_0)/2^ME(k_0)$ has order $|E(k_0)[2^M]|=2^{2M}$, as does
$E[2^M]$, so the induced injection is an isomorphism.
Composing $\kappa_M$ with the fixed row $b$ modulo $2^M$ gives
$\lambda_M$. The arithmetic Frobenius over $\mathbb F_{\ell_0}$
commutes with $F$ and acts on the torsion by $C$. Hence
$\kappa_M(\operatorname{Fr}_{\ell_0}P)=C\kappa_M(P)$ and
$bC=b$ proves the asserted Frobenius sign.
For the fixed witness, specialization of a compatible division
point identifies this evaluation with the cocycle value at
$\sigma^2$, whose $b$-value is $2^{a+1}$. Thus the visibility
valuation is fixed as the precision grows.

By \cite[Main Theorem 14.6.1]{voight-quaternion}, there is a quaternion
algebra $B/\mathbb Q$ ramified exactly at $\ell_0$ and infinity.
Let $R$ be an Eichler
order of level $N$.  Supersingular pairs are
$X=B^\times\backslash\widehat B^\times/\widehat R^\times$;
their automorphism groups are $R_z^\times$ for
$R_z=B\cap z\widehat Rz^{-1}$.  This description allows every
$N$ prime to $\ell_0$, including powers of two.  One may choose each
underlying supersingular curve over $\mathbb F_{\ell_0^2}$ with scalar
Frobenius; its cyclic subgroups of order $N$ are then stable.  Thus
all the corresponding coarse moduli points are defined over that
field.  The geometric and adelic descriptions are those of
\cite[Sections 1.1--1.2 and 2.5]{jetchevkane}.

Define $P_M(z)=\lambda_M(L\pi_{\ell_0}(z))$.  Here the specialization is
legitimate because $\ell_0\nmid N$: the modular curve has good
reduction and the map to the good-reduction elliptic curve extends
over the local base.  The good Hecke correspondences also extend.
Writing $\pi$ as the Abel--Jacobi map based at infinity followed by
the elliptic quotient gives, for $p\nmid N\ell_0$,
\[
 \sum_{z'\text{ a }p\text{-neighbor of }z}\pi_{\ell_0}(z')
       =a_p(E)\pi_{\ell_0}(z).
\]
Indeed the quotient intertwines $T_p$ with $[a_p(E)]$ and
$T_p(\infty)=(p+1)\infty$.  For an Atkin--Lehner involution
$W_{p^e}$ the analogous identity has eigenvalue $w_{p^e}$ and the
constant $\pi(W_{p^e}\infty)$, which $L$ kills.  Frobenius acts on
the supersingular set as the right division-uniformizer operation.
Consequently $P_M$ has the asserted good Hecke and Atkin--Lehner
eigenvalues modulo $2^M$, and uniformizer eigenvalue $+1$.
Choose arbitrary integral lifts of this finite function.

\paragraph{The finite Schwartz function.}
All local functions below are zero off their specified supports.
Set $V=B^0$ and $Q_B(x)=\operatorname{Nrd}(x)=-x^2$.
For $p\nmid2Nm\ell_0$ use $1_{R_p^0}$.
For odd $p^e\Vert N$, use the split integral unit support.  The
cyclic level quotient selects one of the two eigensummands and its
root $r\in\mathbb Z_p^\times$.  The projectors
$(x+r)/(2r)$ and $(r-x)/(2r)$ are integral, so reduction and Hensel
lifting show that the selected root modulo any fixed power of $p$
is locally constant.  Weight it by a unit character $\alpha_p(r)$
of finite $2$-power order with
\[
 \alpha_p(-1)=w_{p^e}\chi_{h_*}(p)^e. \tag{O.3}
\]
Such a character exists: at an odd prime the cyclic $2$-primary
part of $\mathbb F_p^\times$ admits a character taking $-1$ to
$-1$, and the trivial character supplies the other parity.
Conjugation by $R_p^\times$ preserves the cyclic quotient and its
selected summand.  Multiplication of $x$ by a unit $u$ preserves
the summands and replaces $r$ by $ur$, giving multiplier
$\alpha_p(u)$.

At two require $e_x=(1+x)/2\in R_2$ and
$x^2=r^2\in\mathbb Z_2^{\times2}$.  The eigenvalues
$a=(1+r)/2$ and $a'=(1-r)/2$ of $e_x$ differ by the unit $r$.
Its projectors $(e_x-a')/r$ and $(a-e_x)/r$ are integral and remain
distinct modulo two.  If $2^e\Vert N$, the cyclic quotient is a
free rank-one $\mathbb Z/2^e$-module.  Of these complementary
projectors exactly one acts as the identity on it, since the only
idempotents of $\mathbb Z/2^e$ are zero and one.  This defines the
selected root without a choice of its parity.  For every odd $u$,
\[
 e_{ux}=u e_x+(1-u)/2.
\]
This affine transformation preserves the two projectors
individually, so the selected root changes to $ur$, even when
the two eigenvalue parities switch.  The same root character rule
therefore applies at two.  A character with either prescribed
parity exists, using the trivial character or the character modulo
four.  If $2\nmid N$ there is no root character at two.

For $p\mid m$, require $x\in M_2(\mathbb Z_p)^0$ and $v_pQ_B(x)=1$.
The reduction of $x$ is a nonzero nilpotent: zero reduction would
force $p^2\mid Q_B(x)$.  All nonzero nilpotents are conjugate to
$X_0=\left(\begin{smallmatrix}0&1\\0&0\end{smallmatrix}\right)$.
The centralizer of $X_0$ in $\operatorname{GL}_2(\mathbb F_p)$ is
$\{\left(\begin{smallmatrix}a&b\\0&a\end{smallmatrix}\right):a\ne0\}$
and has square determinant.  Thus
\[
 \epsilon_p(gX_0g^{-1})=(\det g/p)
\]
is well-defined.  Conjugation by $u$ multiplies it by $(\det u/p)$,
and scalar multiplication by $u\in\mathbb Z_p^\times$ multiplies
it by $(u/p)$, as follows from conjugation of $X_0$ by
$\operatorname{diag}(u,1)$.

At $\ell_0$ require $x\in R_{\ell_0}^0$ with nonzero reduction
$\bar x=t\beta$ in the trace-zero line of $\mathbb F_{\ell_0^2}$,
and use $(t/\ell_0)$.  Order units act trivially on the residue
field.  A division uniformizer acts by Frobenius, which sends
$\beta$ to $-\beta$, giving multiplier $-1$.
Unit scaling gives $(u/\ell_0)$.  Moreover $\beta^2$ is a nonsquare
in $\mathbb F_{\ell_0}$: otherwise $\beta=\pm a$ for
$a\in\mathbb F_{\ell_0}$, contrary to its nonzero trace-zero condition.
Both $-1$ and $\beta^2$ are nonsquares, so $Q_B(x)=-t^2\beta^2$
is a nonzero square.  Since $(m/\ell_0)=1$, extracted indices have
$(D/\ell_0)=1$.

These supports are compact and open, and all weights are locally
constant.  Additional fixed unit-squareclass conditions multiply
them by locally constant characteristic functions.  They are
preserved by $x\mapsto ux$, which multiplies the norm by $u^2$.
Let $\mathcal O$ be the valuation ring at a place above two of the
field generated by the fixed root-character values at primes
dividing $N$.  All other weights are quadratic, so this field is
independent of $M$ and $\ell_0$.
With $f=\prod f_p$, the preceding calculations give
\[
 f(u^{-1}xu)=\chi_{h_*}(\operatorname{Nrd}u)^{-1}f(x)
       \quad (u\in\widehat R^\times),
 \qquad f(vx)=\alpha(v)f(x)\quad(v\in\widehat{\mathbb Z}^{\times}).
\]
The ramified local quadratic character on units at $p\mid m$ is
$(\cdot/p)$; all other order-unit character factors are trivial.
The scaling character has
\[
 \alpha(-1)=
 \Bigl(\prod_{p^e\Vert N}w_{p^e}\chi_{h_*}(p)^e\Bigr)
             \Bigl(\prod_{p\mid m}(-1/p)\Bigr)(-1)=1.
\]
For the last equality, $m\equiv3\pmod4$ and the first product is
$w(E^{h_*})=+1$.  More explicitly,
$w(E)=-\prod w_{p^e}$ and, since $(h_*,N)=1$,
$w(E^{h_*})=w(E)\chi_{h_*}(-N)$.
The latter twist formula also follows by multiplying the local
unramified-twist factors $\chi_{h_*}(p)^{v_pN}$ at $p\mid N$
and the factors $(-1/p)$ at the newly ramified good primes.

Let $\eta_B$ denote the classical character associated with $\alpha$
in the inverse-idele convention.  The formula
$\eta_B(d)=\prod_p\alpha_p(d)$ for positive integers $d$ prime to
the support is verified in the transformation calculation below.

Set $c_0=|\operatorname{GL}_4(\mathbb F_3)|=24\,261\,120$,
so $v_2(c_0)=9$, and define
\[
 H_M(n)=c_0\sum_{z\in X}
 \frac{P_M(z)\chi_{h_*}(\operatorname{Nrd}z)}{|R_z^\times|}
       \sum_{x\in V(\mathbb Q),\ Q_B(x)=n}f(z^{-1}xz). \tag{O.4}
\]
Left multiplication embeds each $R_z^\times$ into
$\operatorname{GL}_4(\mathbb Z)$; reduction modulo three is
injective on finite subgroups.  Indeed, a prime-order element
$1+3^aA$ in its kernel, with $a\geq1$ and $A$ nonzero modulo three,
has nonzero first term in its prime-th power modulo $3^{a+1}$,
or modulo $3^{a+2}$ when its order is three, a contradiction.
Thus $|R_z^\times|\mid c_0$.
The right-conjugation identity proves independence of the right
representative.  On the left, rational reduced norms are positive
and the global quadratic idele character is trivial on them.
The transformation law below shows that each summand is an actual
holomorphic theta series with the same level and character;
$c_0/|R_z^\times|$ is an integer.
Consequently the coefficients lie in the fixed character ring
$\mathcal O$.  Changing any lift of $P_M$ by $2^M$ changes every
coefficient by an element of $2^M\mathcal O$.  In particular
there is no loss of precision from an automorphism denominator.

\subsection{Transformation law and recurrence}
Use the standard global additive character, self-dual measures,
and the Gaussian $e^{-2\pi Q_B}$ at infinity.  Weil's operator
formulas \cite[I.13]{weil} give multiplication by $\mathbf e_p(bQ_B)$ for
$n(b)$, scaled pullback $f_p(x)\mapsto\gamma_{Q_B}(u)f_p(ux)$ for
$m(u)$ with $u$ a unit, and a constant times Fourier transform
for the Weyl element.  The upper unipotents $n(\mathbb Z_p)$ fix
$f_p$, because $Q_B$ is integral on its support.  If $L'_p$ is any
lattice of translations fixing $f_p$, its Fourier transform has
support in $(L'_p)^\vee$.  Choose $A_p$ with
$p^{A_p}Q_B((L'_p)^\vee)\subseteq\mathbb Z_p$.  Fourier conjugation
then proves that $n^-(p^{A_p}\mathbb Z_p)$ fixes $f_p$.
This includes $p=2$. At two take $A_2\ge2$ and use the comparison
with the ordinary theta multiplier given below; no splitting of
the metaplectic cover over all diagonal units is asserted.
For clarity, the local normalizations and the lower-unipotent
argument are as follows. Take
$\mathbf e_\infty(t)=e^{2\pi it}$ and
$\mathbf e_p(t)=e^{-2\pi i\{t\}_p}$, so that
$\prod_v\mathbf e_v(r)=1$ for $r\in\mathbb Q$.
For $Q(x)=x^{\mathsf T}Ax$ put
$B_Q(x,y)=Q(x+y)-Q(x)-Q(y)=2x^{\mathsf T}Ay$ and set
\[
 \mathcal F_Q\phi(y)=\int\phi(x)\mathbf e_p(B_Q(x,y))\,dx_Q,
 \qquad dx_Q=|\det(2A)|_p^{1/2}\,dx.
\]
Here $dx$ is the coordinate product measure with
$\operatorname{vol}(\mathbb Z_p^3)=1$.
At the real place use the same formula
$dx_Q=|\det(2A)|_\infty^{1/2}dx$ with coordinate Lebesgue measure.
Thus $\mathcal F_Q^2\phi(x)=\phi(-x)$.
If $\phi(x+\lambda)=\phi(x)$ for $\lambda\in L'$, translation
in this integral shows that its Fourier transform vanishes unless
$\mathbf e_p(B_Q(\lambda,y))=1$ for every $\lambda\in L'$.
This is precisely the asserted support in $(L')^\vee$.
Since this dual lattice is compact, an integer $A_p$ with
$p^{A_p}Q((L')^\vee)\subset\mathbb Z_p$ exists. The operator
for $n^-(c)$ is a Fourier conjugate of multiplication by
$\mathbf e_p(-cQ)$; consequently it fixes $\phi$ whenever
$c\in p^{A_p}\mathbb Z_p$. A scalar used to normalize the Weyl
operator cancels with its inverse in this conjugation.

All the specified local functions have such a translation lattice.
At an odd level prime the unit-root condition and its character
depend on a finite residue quotient. At two the square-unit
condition is a congruence modulo eight; the integral projectors
selecting the cyclic quotient, and the finite root character,
again depend on a finite residue quotient. At a prime dividing
$m$, valuation one and the nonzero-nilpotent sign depend on finite
residue data. The division-prime function is already a residue
function. Each support is contained in a lattice and is a finite
union of cosets of a sufficiently small lattice. This proves the
Schwartz--Bruhat and translation assertions for the actual weights.

For the reference $Q_0=\sum X_i^2$, the dyadic dual of
$\mathbb Z_2^3$ is $\tfrac12\mathbb Z_2^3$ and
$\operatorname{vol}_{dx_{Q_0}}(\mathbb Z_2^3)=2^{-3/2}$.
In particular its Fourier transform is
$2^{-3/2}1_{\frac12\mathbb Z_2^3}$, not $1_{\mathbb Z_2^3}$.
The lower-unipotent conclusion nevertheless holds at depth two,
because $4Q_0(\frac12\mathbb Z_2^3)\subset\mathbb Z_2$.
These positive measure factors affect neither the diagonal
eigenvalue quotient nor the counting coefficients of the theta sum.


The diagonal factor can be computed directly.  For
$Q_B=\sum_{i=1}^3 a_iX_i^2$, the one-dimensional Weil-index identity
\[
 \frac{\gamma(ua_i)}{\gamma(a_i)}
       =\frac{\gamma(u)}{\gamma(1)}(u,a_i)_p
\]
and multiplication under orthogonal sums give, relative to the
ordinary theta multiplier cubed, the factor
$(u,a_1a_2a_3)_p$.  These identities are \cite[II.25 and II.28]{weil};
the latter obtains the Hilbert symbol from the four-dimensional
quaternion norm form.  For $B=(-a,-b)$ the norm on $B^0$ is
$aX^2+bY^2+abZ^2$, whose determinant is $(ab)^2$.
There is therefore no additional determinant character for $Q_B$.
\paragraph{The normalized diagonal factor.}
The Weil-index identity used here can be derived without choosing
Gauss-sum signs. The norm form of the quaternion algebra $(u,a)$
is $\langle1,-u,-a,ua\rangle$. By multiplicativity under orthogonal
sums and $\gamma(-q)=\gamma(q)^{-1}$, its index is
\[
 \frac{\gamma(1)\gamma(ua)}{\gamma(u)\gamma(a)}.
\]
For a split quaternion algebra the norm form is hyperbolic and its
index is one; for the division algebra its index is minus one by
\cite[II.28, Proposition 4]{weil}. These are exactly the two values
of $(u,a)_p$. Solving this equality gives
\[
 \frac{\gamma(ua)}{\gamma(a)}
   =\frac{\gamma(u)}{\gamma(1)}(u,a)_p.
\]
Applying it to three diagonal entries gives the ordinary
one-dimensional theta factor cubed times $(u,\det Q_B)_p$.
This derivation includes $p=2$. A change of diagonalization changes
the determinant only by a square, so it changes no symbol.
For $Q_B/m$ the determinant becomes $m^{-3}\det Q_B$, whose
squareclass is $m$. Thus the extra factor is precisely $(u,m)_p$.
This establishes the determinant-character contribution in (O.5)
with the same additive character as the original theta multiplier.

\paragraph{An operator construction of the common double cover.}
Here is a direct justification that the same lifts may be used in
the comparison. It also fixes the signs of the diagonal operators.
Work over $k=\mathbb Q_p$, including $p=2$, or $k=\mathbb R$,
with the additive character just specified. For a nondegenerate
ternary $Q$ write $\gamma_Q=\gamma(Q)$ and define
\[
 \begin{split}
 U_Q(b)\phi(x)&=\mathbf e_k(bQ(x))\phi(x),\\
 D(a)\phi(x)&=|a|_k^{3/2}\phi(ax),\\
 W_Q&=\gamma_Q\mathcal F_Q,\qquad
 M_Q(a)=r_Q(a)D(a),\qquad
 r_Q(a)=\frac{\gamma(aQ)}{\gamma(Q)}.
 \end{split}
\]
The Weyl matrix in this convention is
$w=\left(\begin{smallmatrix}0&1\\-1&0\end{smallmatrix}\right)$.
The opposite Fourier-sign convention changes both compared
representations together.

The Fourier formula of \cite[I.14, Theorem~2]{weil}, or its
completion-of-the-square calculation, says for $b\ne0$ that
\[
 \int_{k^3}
 \mathbf e_k\bigl(bQ(z)+B_Q(z,v)\bigr)\,dz_Q
   =|b|_k^{-3/2}\gamma(bQ)\mathbf e_k(-b^{-1}Q(v)).
\]
The integral has the oscillatory meaning of that theorem.
Applying it to a Schwartz function gives the operator identity
\[
 \mathcal F_Q U_Q(b)\mathcal F_Q
 =\gamma(bQ)U_Q(-b^{-1})D(b^{-1})
                   \mathcal F_Q^{-1}U_Q(-b^{-1}).
\]
Thus no choice of an individual square root of a Gauss sum is
needed. The index identities proved above also give
\[
 \gamma(Q)^2=(\det Q,-1)_k\gamma(1)^6,
 \qquad \gamma(1)^4=(-1,-1)_k,
 \qquad r_Q(a)r_Q(b)=(a,b)_k r_Q(ab).
\]
For the middle equality apply the quaternion-norm formula to
$\langle1,1,1,1\rangle$; the first follows by applying the same
formula to $\langle1,-a_i,1,-a_i\rangle$ for each diagonal
entry $a_i$ of $Q$.

Consequently the displayed operators satisfy
\[
 \begin{gathered}
 U_Q(b)U_Q(c)=U_Q(b+c),\qquad
 M_Q(a)U_Q(b)M_Q(a)^{-1}=U_Q(a^2b),\\
 M_Q(a)M_Q(b)=(a,b)_kM_Q(ab),\qquad
 W_QM_Q(a)W_Q^{-1}=M_Q(a^{-1}),\\
 W_Q^2=(-1,-1)_kM_Q(-1),\\
 W_QU_Q(b)W_Q
 =(-b,-1)_kU_Q(-b^{-1})M_Q(-b^{-1})W_QU_Q(-b^{-1})
 \quad(b\ne0).
 \end{gathered}
\]
For the last equality, the preceding Fourier identity leaves
the scalar $\gamma(Q)^2\gamma(bQ)^2$; the index-square formula
reduces it to $(-b,-1)_k$. Every scalar in these relations is
independent of the ternary form. Interpreting a scalar $-1$ as
the central element $\varepsilon$ gives one common double
cover of $\operatorname{SL}_2(k)$. Indeed these relations reduce
every word to one of
$\varepsilon^eU(x)M(a)$ or
$\varepsilon^eU(x)M(a)WU(y)$, with $e\in\{0,1\}$.
The ordinary Bruhat decomposition makes their projections
unique, apart from $\varepsilon^e$; the operators send
$\varepsilon$ to $-\operatorname{id}$, so the central kernel
has exactly two elements. This realizes the rank-three pullback
of Weil's metaplectic operators, rather than assuming a section
over the diagonal units.

For $Q_0=\langle1,1,1\rangle$ these operators are the tensor
products of the three ordinary one-dimensional theta operators.
For either $Q$ or $Q_0$, changing a lift multiplies its operator
by the same $-1$. Hence the quotient of their eigenvalues is
independent of the lift. At two the relation
$M_Q(-1)^2=(-1,-1)_2\operatorname{id}=-\operatorname{id}$
explicitly retains the order-four lift of $-I$.

On a unit $a$ the quotient of the two diagonal index factors is
\[
 \frac{r_Q(a)}{r_{Q_0}(a)}=(a,\det Q)_k.
\]
Together with $f_p(ax)=\alpha_p(a)f_p(x)$, this is exactly
the eigenvalue quotient $\alpha_p(a)(a,\det Q)_p$ used below.
There is no remaining eighth-root factor at two.

\paragraph{The metaplectic comparison at two and the extracted character.}
We specify how the preceding index calculation fixes the multiplier,
including its dyadic part. Let
\[
 Q_0=X_1^2+X_2^2+X_3^2,\qquad
 \phi_{0,p}=1_{\mathbb Z_p^3},\qquad
 \theta(\tau)=\sum_{n\in\mathbb Z}e^{2\pi i n^2\tau}.
\]
Use the same additive character and metaplectic lifts for the
rank-three Weil representations of $Q_B$ and $Q_0$. The latter
theta series is literally $\theta(\tau)^3$; this fixes the meaning
of the ordinary theta multiplier cubed, including its eighth roots
of unity. For either representation the nontrivial central element
of the metaplectic cover acts by $-1$. Consequently the ratio of
two eigenvalues computed using the same lift is independent of
that lift. In particular a central factor in multiplying the
three lifts in the displayed triangular factorization cancels
in this ratio. We compare eigenvalues, and do not choose a group
section of the cover over $\Gamma_0(2^{A_2})$.

Here is the comparison on the three factors. Both upper unipotents
act trivially on the indicated vectors. Choose $A_p$ large enough
that both lower unipotents do also. This is the Fourier-support
argument above for $f_p$; for $\phi_{0,p}$ it requires no positive
depth at an odd prime and depth two at two, since the bilinear dual
of $\mathbb Z_2^3$ for $Q_0$ is $\tfrac12\mathbb Z_2^3$ and
$4Q_0(\tfrac12\mathbb Z_2^3)\subset\mathbb Z_2$.
The Weyl scalar cancels in conjugating an upper unipotent to a
lower one, so it introduces no additional eigenvalue. On a
diagonal unit $u$, the scaling part acts by $\alpha_p(u)$ on
$f_p$ and by one on $\phi_{0,p}$. The quotient of the Weil-index
parts, by the preceding calculation, is
\[
 \frac{\gamma(uQ_B)/\gamma(Q_B)}
      {(\gamma(u)/\gamma(1))^3}
       =(u,\det Q_B)_p.                                      \tag{O.5a}
\]
If the inverse Weil-index convention is used for the diagonal
operator, both numerator and denominator are inverted; the
quotient is unchanged because the Hilbert symbol has order two.
Thus this computation does not leave an undetermined inverse
or a dyadic root of unity. These operator comparisons are precisely
the scaled pullback, Fourier conjugation and scalar relations of
\cite[I.13--I.14 and II.25--II.28]{weil}.

Here $\det Q$ means the determinant of the matrix in
$Q(x)=x^{\mathsf T}Ax$. The associated bilinear Gram matrix is
$2A$. If that convention is used instead, both $Q$ and $Q_0$
acquire the same factor $2^3$ in their determinants, which cancels
in the comparison. In particular no extra squareclass of two
is introduced at the dyadic place.

For a general ternary form $Q$ in this argument, at each prescribed
prime take $A_p\ge1$ and increase $A_p$
also to the conductor of $u\mapsto(u,\det Q)_p$. The eigenvalue
relative to the ordinary theta multiplier, on a matrix whose
upper-left entry is $a$, is then
\[
                    \alpha_p(a)(a,\det Q)_p.                  \tag{O.5b}
\]
Indeed the triangular factorization has just the two trivial
unipotent factors and this diagonal factor, and its lift ambiguity
cancels as explained above. This also verifies multiplication:
the upper-left entry of a product is congruent to the product
of the upper-left entries modulo $p^{A_p}$, so both local
characters in (O.5b) multiply. The argument therefore applies
to every matrix in the indicated congruence subgroup at a prescribed
prime, not only to a choice of generators. At an unramified prime
outside the support, the full spherical invariance established
below treats matrices with nonunit upper-left entry as well.

For $Q_B$ the determinant squareclass is one at every place,
including two. Thus (O.5b) has no determinant factor. For the
extracted form $Q_B'=Q_B/m$ it is $m^{-3}\equiv m$.
At an odd prime $p\mid m$ and a unit $u$, and at two, respectively,
the additional factors are explicitly
\[
 (u,m)_p=(u/p)^{v_p(m)},\qquad
 (u,m)_2=(-1)^{((u-1)/2)((m-1)/2)}.
\]
Here $m$ is odd and $m\equiv3\pmod4$, so the second factor is
$(-1)^{(u-1)/2}$. It is already trivial on $1+4\mathbb Z_2$.
At every odd prime outside $m$ the factor on units is one. There
is consequently no new level prime, and there is no omitted
dyadic character.

Finally take a positive integer $d$ prime to the complete level.
The finite operator used in the classical transformation is the
one at $\gamma^{-1}$, with upper-left entry $d$. The root-character
factor is therefore $\prod_p\alpha_p(d)=\eta_B(d)$.
Hilbert reciprocity gives
\[
 \prod_{p\mid 2m}(d,m)_p
   =\prod_{q\mid d}(m/q)^{v_q(d)}=(m/d).                       \tag{O.5c}
\]
Indeed the infinite symbol is one because $m,d>0$, all other
finite symbols are one, and at an odd $q\mid d$ the symbol is
$(m/q)^{v_q(d)}$. This proves the displayed product without
discarding its dyadic factor. Both sides extend to the same
even Dirichlet character. Since $h_*=-m$, division by the genus
character now gives
\[
 \frac{\eta_B(d)(m/d)}{\eta_B(d)\chi_{h_*}(d)}=(-1/d).
\]
Thus the characters in (O.5), with the stated inverse-idele
convention for $\eta_B$, are exactly the characters supplied
by these local Weil operators.

Take $A_p$ also at least the conductor exponent of $\alpha_p$.
For $\gamma=\left(\begin{smallmatrix}a&b\\c&d\end{smallmatrix}\right)$
with $c\in p^{A_p}\mathbb Z_p$ and $a$ a unit, use
\[
 \gamma=n^-(c/a)m(a)n(b/a).
\]
The scalar is the diagonal theta factor times $\alpha_p(a)$.
In passing to the classical transformation law one applies the
finite operator at $\gamma^{-1}$, whose upper-left entry is $d$.
Thus the classical Dirichlet character is
$\eta_B(d)=\prod_p\alpha_p(d)$, not its inverse.  Its value at
$-1$ is one, as required by the weight-$3/2$ theta multiplier.

There are no additional level primes hidden in this argument.
Outside the prescribed support, $p$ is odd, $B_p$ is split, and
$R_p^0=\mathfrak{sl}_2(\mathbb Z_p)$ is self-dual for the bilinear
form of $Q_B=-a^2-bc$.  Its determinant is a unit.  Hence the
Fourier transform of $1_{R_p^0}$ equals itself; the unramified
Weil constant is one, as also follows from \cite[II.27]{weil} by integrating
over this self-dual lattice.  Upper unipotents, diagonal units,
and the Weyl element fix it.  These generate
$\operatorname{SL}_2(\mathbb Z_p)$, so the exponent there is zero.
At the finitely many prescribed primes the preceding $A_p$ suffice.
Conjugation by any $z\in\widehat B^\times$ is an isometry of $Q_B$,
commutes with the Weil representation, and preserves the
self-dual measure.  The same level and character therefore work
for every summand indexed by $z$.

\cite[III, Theorem 6]{weil} applies to this Schwartz--Bruhat function on
$V(\mathbb A)$ and proves invariance of its rational theta sum.
The normalized archimedean expression is
$\sum_x f(x)e^{2\pi i\tau Q_B(x)}$: the factor $y^{3/4}$ in the
Weil action is removed in passing to the classical form.
It has holomorphic weight $3/2$.  At a rational cusp, the slash
operation replaces the finite Schwartz function by its finite
Weil transform.  Its support meets $V(\mathbb Q)$ in a finite
union of rational lattice cosets, with bounded denominators.
Positive definiteness gives nonnegative rational exponents of
bounded denominator and normal convergence for $\Im\tau>0$.
This proves holomorphy at every cusp.  Theorem 6 alone asserts
automorphic invariance; this last argument supplies the
holomorphy conclusion used here.

\paragraph{Extraction.}
The support at every $p\mid m$ has $v_pQ_B=1$.  All nonzero
indices are integral and divisible by $m$.  Repeat the preceding
construction for $Q_B'=Q_B/m$ with the same finite weights and the
Gaussian for $Q_B'$.  This quadratic form is integral on the
support.  Away from $m$ it differs by a unit, so the good-prime
self-duality argument remains valid.  Its determinant squareclass
is $m^{-3}\equiv m$, giving the extra Dirichlet character
$(m/\cdot)$.  Hence
\[
 g_M(\tau)=\sum_n H_M(mn)q^n\in
      M_{3/2}(\Gamma_0(K),\psi).
\]
Here $K$ has only the stated prime support, and
\[
 \psi=\eta_B(m/\cdot),\qquad
 \eta=\eta_B\chi_{h_*},\qquad
 \psi/\eta=\chi_{-1}. \tag{O.5}
\]
This directly constructs the extracted form without a
half-integral-weight extraction operator.
The inverse in this classical passage can also be checked before
any character notation is introduced. Put
$g_\tau=n(x)m(\sqrt y)$ for $\tau=x+iy$, so $g_\tau i=\tau$.
For $\gamma\in\operatorname{SL}_2(\mathbb Q)$ write
$\gamma g_\tau=g_{\gamma\tau}k_\infty$ with
$k_\infty\in\operatorname{SO}_2(\mathbb R)$. In the adelic group,
\[
 (g_{\gamma\tau},1_f)
   =\gamma\,(g_\tau,\gamma_f^{-1})(k_\infty^{-1},1_f).
\]
This equality also holds with consistent metaplectic lifts.
The global rational lifts are compatible for the two forms:
\cite[II.30, Proposition~5]{weil} gives
$\prod_v\gamma_v(Q)=\prod_v\gamma_v(aQ)=1$ for rational $Q,a$,
and hence the local normalizing factors in $W_Q$ and $M_Q(a)$
have product one. The product formula also gives
$\prod_v|\det(2A)|_v^{1/2}=1$, so the product of the stated
self-dual measures is the measure for which the rational lattice
has covolume one. Apply the rational invariance of
\cite[III.41, Theorem~6]{weil} to the displayed adelic identity.
The archimedean Gaussian for either positive ternary form has
the same weight $3/2$; therefore its factor from $k_\infty^{-1}$
is the same as that for $\theta^3$. The relative finite factor
is evaluated at $\gamma_f^{-1}$, whose upper-left entry is $d$.
Thus the relative classical character is
$\prod_p\alpha_p(d)(d,\det Q)_p$, over the prescribed primes.
This proves the inverse-idele convention directly. It compares
the transformation identities and does not divide by a value
of a theta function.

For the extracted form there is likewise no measure or degree
normalization in its coefficients. With exactly the same finite
weights, replace only the quadratic form and archimedean Gaussian
by $Q_B/m$ and $e^{-2\pi Q_B/m}$. The summand indexed by $x$ now
has exponent $Q_B(x)/m$. The original support makes $Q_B(x)$ an
integer with valuation one at every prime dividing $m$, so these
exponents are nonnegative integers. Coefficient $n$ is exactly
the original sum over $Q_B(x)=mn$, namely $H_M(mn)$.
The self-dual measure changes by $|m|_v^{-3/2}$ locally, but a
theta coefficient is a discrete sum and contains no Haar
measure; this change is used only in the Fourier operators.
Outside the original support it is a unit rescaling of the same
unimodular local form, and at primes dividing $m$ the
lower-unipotent depth can be increased within the existing
support. The determinant changes by $m^{-3}$, whose squareclass
is $m$, so (O.5c) supplies exactly $(m/\cdot)$.
In particular no new level prime, factor of $m$, or dyadic
root of unity is introduced by extraction.


\paragraph{All-index recurrence and its phases.}
Let $p$ avoid the full support and put $s=\eta_B(p)$.
Locally use $R_p=M_2(\mathbb Z_p)$ and $V_p=\mathfrak{sl}_2$.
The $p+1$ neighbors are indexed by lines of $\mathbb F_p^2$.
For integral $x$ with nonzero reduction and $Q_B(x)=n$, the number
of preserved lines is $1+(-n/p)$, including the nilpotent case.
If $x\in pR_p^0$, all $p+1$ lines are preserved; the excess
over the preceding formula is $p$.

A nonintegral $x$ in a neighbor lies in $p^{-1}R_p^0$.
Write $x=y/p$ with $y$ primitive and $Q_B(y)=p^2n$.
Then $\bar y$ is nonzero nilpotent.  In a basis with
$\bar y=X_0$, write
$y=\left(\begin{smallmatrix}a&b\\c&-a\end{smallmatrix}\right)$.
Here $p\mid a,c$, $b$ is a unit, and
$p^2\mid-a^2-bc$ forces $p^2\mid c$.  Exactly the neighbor
belonging to the kernel line makes $y/p$ integral.  The integral
and nonintegral parts therefore exhaust the sum with no genericity
assumption on $n$.

Scaling by $p$ at all other places multiplies the weight by $s$.
Thus the primitive nonintegral contribution is
$s^{-1}H_M(p^2n)-H_M(n)$, and the integral contribution is
$(1+(-n/p))H_M(n)+psH_M(n/p^2)$.  Their sum is
\[
 s^{-1}H_M(p^2n)+(-n/p)H_M(n)+psH_M(n/p^2).
\]
Here and below a coefficient at a nonintegral index is zero.
For completeness, if $N_{ij}$ counts neighbor kernels from
$i$ to $j$ and $w_i=|\operatorname{Aut}(i)|$, counting actual
isogenies and dualizing gives
$N_{ij}w_j=N_{ji}w_i$, or $N_{ij}/w_i=N_{ji}/w_j$.
This proves the self-adjointness with the mass $1/w_i$ appearing
in $H_M$.  A neighbor changes the norm character by
$\chi_{h_*}(p)$.  Moving the neighbor operator onto $P_M$ and
using its actual modular-parametrization identity consequently
gives
\[
 H_M(p^2n)+s(-n/p)H_M(n)+ps^2H_M(n/p^2)
     =s\chi_{h_*}(p)a_p(E)H_M(n)\pmod {2^M\mathcal O}.
\]
This computation explains both powers of $s$ even when $s$ is
not quadratic.

Put $a_M(n)=H_M(mn)$ and
$\eta=\eta_B\chi_{h_*}$.  Since
$(-mn/p)=(m/p)(-n/p)$ and $\eta(p)^2=s^2$, one obtains
\[
 a_M(p^2n)=
 \{\eta(p)a_p(E)-\psi(p)(-n/p)\}a_M(n)
        -p\eta(p)^2a_M(n/p^2)\pmod {2^M\mathcal O}. \tag{O.6}
\]
The proof covers $p\mid n$ and $p^2\mid n$ explicitly.
At zero the coefficients vanish by the local support, and negative
indices vanish by positive definiteness.
An added finite coprimality or unit-squareclass sieve is realized
by the corresponding local characteristic function for $Q_B'$;
the same modularity proof applies.  At primes outside the enlarged
support the three indices $p^2n,n,n/p^2$ have the same retain-or-zero
status.  Hence these sieves preserve the displayed recurrence.
\end{proof}

\subsection{The unaveraged point comparison}
\paragraph{The actual weighted coefficient comparison.}
\begin{proposition}[The unaveraged weighted theta comparison]
\label{ex:app-theta-comparison}
Put $m=|h_*|$ and $F=\mathbf Q(\sqrt{Dh_*})$, where $D$ is a
positive squarefree parameter in the fixed odd filter. Assume
$(D/\ell_0)=1$ at the chosen actual stage, so that $\ell_0$ is
inert in $F$. Keep the local Schwartz functions, the actual function
$P_M(z)=\lambda_M(L\pi_{\ell_0}(z))$, and the integer $c_0$
of (O.4). Then a $2$-power root of unity $\zeta_M(D)$ satisfies
\[
 H_M(mD)=c_0\,2^{\omega(N)}\zeta_M(D)
                   \lambda_M(\overline{LP_D})\pmod {2^M}.\tag{O.7}
\]
The root of unity belongs to the field generated by the fixed
characters at the primes of $N$. No class number is divided out.
\end{proposition}


\begin{proof}
The filter makes $Dh_*$ an odd fundamental discriminant, with
$F$ split at every prime dividing $2N$. Its unit group is
$\{\pm1\}$. The condition $Q_B(x)=mD$ is equivalent to
$x^2=Dh_*$. Together with $(1+x)/2\in R_{z,2}$ and the other
integrality conditions, this says exactly that
\[
 \mathcal O_F\longrightarrow R_z,
 \qquad \sqrt{Dh_*}\longmapsto x
\]
is an optimal embedding. Optimality follows because the source
order is already maximal. All nonzero local weights on such an
embedding are roots of unity. The valuation-one condition at each
prime of $h_*$ holds because $(D,h_*)=1$; at $\ell_0$ the residue
of $x$ is nonzero. The fixed unit filters hold by admissibility.

First classify the local embeddings. Away from $N\ell_0$ there
is one local orbit: an $\mathcal O_{F,p}$-stable lattice is locally
a fractional ideal, hence principal, and this also covers ramified
primes dividing $Dh_*$. At $p^e\Vert N$, the algebra is split.
The two integral idempotents decompose both lattices of the
Eichler lattice pair. Its quotient is cyclic of order $p^e$,
so its elementary exponents lie wholly in one summand. The two
possibilities give exactly two orientations. This applies at two
using the integral spectral projectors specified in the construction.
At $\ell_0$, the maximal order of the division algebra has residue
field $\mathbf F_{\ell_0^2}$. The two residue-field embeddings
of the unramified quadratic order give the two orientations.
Consequently there are $2^{\omega(N)+1}$ local orientation tuples.

Fix one tuple. Choose a rational embedding
$\iota:F\hookrightarrow B$ and write $x_0=\iota(\sqrt{Dh_*})$.
Such an embedding exists because $F$ is a field at the two ramified
places of $B$, namely $\ell_0$ and infinity
\cite[Proposition 14.6.7, $(v)\Rightarrow(iii)$]{voight-quaternion}.
Choose $z_0$ realizing
the selected integral local embeddings. The adelic representatives
with this tuple are exactly
\[
       \widehat F^{\times}z_0\widehat R^{\times}.
\]
Their quotient by the rational centralizer $F^{\times}$ is
\[
 F^{\times}\backslash\widehat F^{\times}/
      (\widehat F^{\times}\cap z_0\widehat R^{\times}z_0^{-1})
 =F^{\times}\backslash\widehat F^{\times}/\widehat{\mathcal O}_F^{\times}
 =\operatorname{Pic}(\mathcal O_F).
\]
Thus this is an actual class-group torsor, not an average or a
quotient of a class-group orbit.

For an ideal-class representative $a\in\widehat F^{\times}$,
the element $a$ centralizes $x_0$, and therefore
\[
 f((az_0)^{-1}x_0(az_0))=f(z_0^{-1}x_0z_0).
\]
On the other hand,
\[
 \chi_{h_*}(\operatorname{Nrd}(az_0))
 =\chi_{h_*}(\operatorname{N}_{F/\mathbf Q}a)
                    \chi_{h_*}(\operatorname{Nrd}z_0).
\]
The first factor is exactly the genus character in (O.1).
To see that it is well-defined on the class group, for $D>1$
the biquadratic field $\mathbf Q(\sqrt D,\sqrt{h_*})$ has
discriminant $D\,h_*\,(Dh_*)=(Dh_*)^2$. Its relative
discriminant over $F$ is therefore one. The character obtained
by restricting $\chi_{h_*}$ to $G_F$ is unramified and, under
class field theory, is $a\mapsto
\chi_{h_*}(\operatorname{N}_{F/\mathbf Q}a)$. For $D=1$ it
is the trivial character. Ideal representatives can be chosen
prime to $h_*$, giving exactly the elementary norm formula in
(O.1). Reversing the Artin convention has no effect because this
character is quadratic.

Here is the precise specialization of the reduction theorem, including
its orbit and orientation data. In the notation of
\cite[\S\S3.1--3.3]{cornutjetchev}, take
\[
 F_{\rm CJ}=\Q,\quad K_{\rm CJ}=F=\Q(\sqrt{Dh_*}),\quad
 B_{\rm CJ}=M_2(\Q),\quad S=\{\ell_0\}.
\]
The field $F$ is imaginary because $D>0$ and $h_*<0$.
It embeds into $M_2(\Q)$ by its regular representation. The source
algebra has no finite ramification, so H.1 holds and the integer
in H.2 is $1+0+1=2$. Moreover $\ell_0\nmid Dh_*N$ and
\[
  (Dh_*/\ell_0)=(D/\ell_0)(h_*/\ell_0)=1\cdot(-1)=-1
\]
by the present support condition and (O.2). Thus H.3 holds in its
inert, rather than ramified, case. The resulting target algebra is
the definite algebra $B$ ramified precisely at $\ell_0$ and infinity.

Take a source Eichler order $R^{\rm src}=R'\cap R''$ of level $N$.
In \cite[\S3.2]{cornutjetchev} the source subgroup is
$H=\widehat{R^{\rm src}}^{\times}$ and the target subgroup is
$H_S=\widehat R^{\times}$. Indeed, away from $\ell_0$ the local
orders are transported by the fixed algebra isomorphisms. At $\ell_0$
the source subgroup is $\operatorname{GL}_2(\Z_{\ell_0})$,
whose determinant is $\Z_{\ell_0}^{\times}$; its inverse image
under the target reduced norm is exactly the unit group of the
maximal division order. The target level is the prime-to-$S$ part
of $N$, which is $N$ itself. This construction allows arbitrary
prime-power exponents in $N$, including its exponent at two.

Both fine conductors are one. To check this assertion using the
definition in \cite[\S3.3]{cornutjetchev}, an embedding of the maximal
order $\mathcal O_F$ into $R'\cap R''$ embeds it into each maximal
order. The intersection of either maximal order with $F$ is an order
containing $\mathcal O_F$, hence is $\mathcal O_F$. Conversely,
if both intersections are $\mathcal O_F$, their intersection with
$R'\cap R''$ is $\mathcal O_F$. The same argument applies on the
target. Thus $c'=c''=c=c_S=(1)$ in that theorem. Its parity
condition at $\ell_0$ is $0-0\equiv v_{\ell_0}(N)=0\pmod2$.
The local lattice pair there has distance zero. Since $F_{\ell_0}$
is the unramified quadratic extension, its maximal order fixes a
unique vertex of the local tree. Consequently
\[
 N_{\ell_0}(0,0,0)=1,\qquad
 |\operatorname{Gal}(F[c]/F[c_S])|=1,
\]
which also follows directly from
\cite[Lemma~2.1(ii)]{cornutjetchev}. The multiplicity $\kappa$ in
\cite[Theorem~3.1]{cornutjetchev} is therefore exactly one.

It is essential to retain the sign fiber in that theorem. At $\ell_0$
the quotient
\[
 F_{\ell_0}^{\times}\backslash B_{\ell_0}^{\times}/R_{\ell_0}^{\times}
       \simeq\Z/2\Z
\]
is measured by the valuation of the reduced norm modulo two:
norms from $F_{\ell_0}^{\times}$ have even valuation, whereas a
division uniformizer has norm valuation one.
By \cite[\S3.6]{cornutjetchev} this sign is precisely the choice of
residue-field embedding of the unramified quadratic order, with
origin fixed by the chosen reference embedding. Theorem~3.1 gives
a bijection onto one such sign, not onto both. Away from $\ell_0$
its map on local orbit data is a bijection, so we may also restrict
to the chosen orientations at the primes of $N$. The resulting
map is an equivariant bijection of the source class-group torsor
onto the target torsor for that full orientation tuple. Corollary~3.2
gives the same bijection after fixing any supersingular pair in the
target.

Choose the reduction place compatibly with the place used in (O.2).
Sections~4.9--4.10 of \cite{cornutjetchev} identify the projection
of this adelic map with geometric reduction of the corresponding
enhanced elliptic curve on $X_0(N)$. Their local formula is the
identity under the chosen identifications away from $\ell_0$, and
at $\ell_0$ replaces a matrix by a division-uniformizer power
whose exponent is the valuation of its determinant. Thus it is
the same sign and the same reduction map as above. Choose $z_0$
to be the image of the actual reference Heegner point $x_D$ in
(O.1). Writing
\[
 \chi_F([a])=\chi_{h_*}(\operatorname{N}_{F/\Q}a),\qquad
 z_a=[az_0],
\]
equivariance gives the exact finite-sum identity
\[
 \sum_{[a]\in\operatorname{Pic}(\mathcal O_F)}
       \chi_F([a])P_M(z_a)
 =\lambda_M\!\left(
       \overline{L\sum_{[a]\in\operatorname{Pic}(\mathcal O_F)}
                         \chi_F([a])\pi(x_D^{[a]})}\right)
 =\lambda_M(\overline{LP_D})\quad\text{in }\Z/2^M\Z.
\]
Here the supersingular pairs and their parametrization values are
defined over $\mathbb F_{\ell_0^2}$, so every term is in the domain
of the same $\lambda_M$. Additivity, specialization, and multiplication
by the fixed integer $L$ commute with this finite sum. If the Artin
action is written with inverses, reindexing $[a]\mapsto[a]^{-1}$
leaves it unchanged since $\chi_F$ is quadratic. Changing the
reference point within its torsor multiplies the sum only by
$\chi_F$ of that reference class.

The weights are constant on each embedding orbit used in the
following count. Indeed a unit of the definite order $R_z$ has
positive integral reduced norm with integral inverse, so its norm
is one. The conjugation law for $f$ therefore gives weight one
for its action. Along the class-group torsor the Schwartz factor
is constant because $a$ centralizes $x_0$, and the remaining
norm factor is exactly $\chi_F([a])$. Hence the finite-sum identity
just proved applies with the single phase
$\chi_{h_*}(\operatorname{Nrd}z_0)f(z_0^{-1}x_0z_0)$.


For each fixed $z$, the stabilizer of an embedding in $R_z^\times$
is $\mathcal O_F^\times=\{\pm1\}$. Each embedding orbit thus
has $|R_z^\times|/2$ elements. The denominator $|R_z^\times|$
in (O.4) leaves the factor $1/2$ per embedding class. The
vectors $x$ and $-x$ are not identified by this stabilizer:
at $\ell_0$ their residues define the two different embeddings
of the residue field, whereas conjugation by an order unit acts
trivially on that field. They already lie in the two different
$\ell_0$ orientation tuples counted above. The
contribution of the reference orientation tuple is
\[
 \frac{c_0}{2}\,\zeta_M(D)
                         \lambda_M(\overline{LP_D}),\qquad
 \zeta_M(D)=\chi_{h_*}(\operatorname{Nrd}z_0)
                               f(z_0^{-1}x_0z_0),
\]
up to the harmless quadratic reference-class factor just described.
All nonzero factors in this displayed phase are fixed-character
values or signs. Thus it is a $2$-power root of unity; its field
does not grow with $D$, $M$, or $\ell_0$.

It remains to account for the other orientations; they are not
all in the image of the one sign component of the reduction map.
At $p^e\Vert N$, use the local Atkin--Lehner normalizer. It
reverses the selected root, so multiplies the Schwartz weight by
\[
              \alpha_p(-1)=w_{p^e}\chi_{h_*}(p)^e.
\]
It multiplies the norm-weighted parametrization value by the same
sign: its reduced norm has valuation $e$, and the actual
parametrization transforms by $w_{p^e}$ after $L$ kills the fixed
cuspidal translation. The two signs multiply to one. This
normalizer has trivial components at the other primes, so it
changes no other orientation or local weight.

At $\ell_0$, a division-algebra uniformizer exchanges the two
residue embeddings. It sends the trace-zero residue $\beta$ to
$-\beta$, multiplying its weight by $(-1/\ell_0)=-1$. Its
reduced norm has valuation one and hence contributes
$\chi_{h_*}(\ell_0)=-1$. Its action on supersingular pairs is
Frobenius. Since $\lambda_M$ has Frobenius sign $+1$ and the
parametrization is defined over $\mathbf Q$, this introduces no
further sign. Again the two factors multiply to one.
The uniformizer operation used here may be described with arithmetic
or geometric Frobenius: they have the same action on the supersingular
pairs and their parametrization values over $\mathbb F_{\ell_0^2}$.
The identity $\lambda_M(FP)=\lambda_M(P)$ from (O.2) consequently
applies to the very operation that exchanges the residue embeddings.
Its norm has valuation one, and $\chi_{h_*}$ is unramified at
$\ell_0$, so its norm-character value is exactly
$\chi_{h_*}(\ell_0)=-1$, independently of the uniformizer's unit part.
Likewise, at $p^e\Vert N$, the norm of the Atkin--Lehner normalizer
has valuation $e$ and $\chi_{h_*}$ is unramified at $p$. Its
norm-character value is exactly $\chi_{h_*}(p)^e$. These observations
fix the norm signs in both orientation changes.


It follows that every one of the $2^{\omega(N)+1}$ tuples has
the same contribution as the reference tuple. Multiplying by
their number gives precisely $c_0 2^{\omega(N)}$, as asserted.
This counts individual embeddings with their stabilizer; a
conjugate-pair formulation of the geometric correspondence must
not be used to insert or remove another factor of two.

Finally, $c_0/|R_z^\times|$ is an integer for every $z$. Thus a
change of the chosen integer lift of $P_M(z)$ by $2^M$ changes
the actual coefficient by an element of $2^M\mathcal O$.
Every preceding equality is consequently valid at the original
retained precision. No denominator depending on the class number,
the discriminant, or the stage has entered.
\paragraph{Precision in the orientation count.}
The factor $1/2$ in the orbit count does not require dividing a
congruence modulo $2^M$ by two. Indeed, $-1\in R_z^\times$ implies
$2\mid |R_z^\times|$, and the chosen integer $c_0$ is divisible by
$|R_z^\times|$. Hence $c_0/2$ is an integer. Choose lifts of all
$P_M(z)$ before summing. Each orbit of embeddings contributes
$(c_0/|R_z^\times|)(|R_z^\times|/2)P_M(z)=(c_0/2)P_M(z)$,
with its prescribed unit weight. Replacing any lift by a multiple
of $2^M$ changes this expression by a multiple of $2^M$.
The comparison with the reduced point may therefore be made at
precision $M$, separately for each orientation. Summing the
$2^{\omega(N)+1}$ orientations gives the stated multiplier
$c_0 2^{\omega(N)}$ without any loss of precision.
In particular, write $c_\theta=9+\omega(N)$, the valuation of
the fixed integer $c_0 2^{\omega(N)}$. At the fixed witness $D_w$,
(O.2) and (O.7) imply
\[
 v_2a_M(D_w)=c_\theta+c_w
       \qquad\text{whenever }M>c_\theta+c_w.
\]
The phase is a unit and the congruence error has strictly larger
valuation, so this is an equality, not merely a lower bound.
Discarding the finitely many earlier stages is sufficient. Thus the
visibility used in choosing the odd minimum retains a fixed witness
and a fixed valuation; neither the class number nor the changing
prime $\ell_0$ introduces a new loss of precision.

\end{proof}

\begin{proposition}[Visibility and the native genus index]\label{ex:app-visibility}
Use the actual odd form and unaveraged point of
Propositions~\ref{ex:app-theta-form} and~\ref{ex:app-theta-comparison}.
There is one constant $C_{\rm det}$, independent of the actual stage
and admissible $D$, for which (O.8) holds. Explicitly, for a known
simple-zero $D$,
$v_2a_M(D)\ge\min\{M,j_{\rm gen}(D)-C_{\rm det}\}$.
For every admissible $D$, the strict inequality $v_2a_M(D)<M$
implies $\operatorname{ord}_{s=1}L(E^D,s)=1$ and
$j_{\rm gen}(D)\le v_2a_M(D)+C_{\rm det}$.
If an admissible odd-sign $D$ is not a simple zero, its coefficient
vanishes modulo $2^M$. No genus index of a torsion point is defined
or used in this assertion.
\end{proposition}
\begin{proof}\leavevmode
For a known simple-zero $D$, write $P_D=m_DG+T$ in its native
rank-one free lattice. The same $L$ kills $T$.
An integral finite-field Kummer map cannot decrease
$v_2(m_D)$, so (O.7) gives the lower bound by the native index.
The rational free line has index at most two in that native
lattice, by conjugation and norm. Absorb this and the fixed
initial convention in $C_{\rm det}$. The first inequality below
holds for a known simple-zero $D$; the visibility implication holds
for every admissible $D$:
\[
 \begin{split}
 v_2a_M(D)&\ge\min\{M,j_{\rm gen}(D)-C_{\rm det}\},\\
 v_2a_M(D)<M&\ \Longrightarrow\
  \ord_{s=1}L(E^D,s)=1,\quad
       j_{\rm gen}(D)\le v_2a_M(D)+C_{\rm det}.
 \end{split}                                                    \tag{O.8}
\]
If an odd-sign $D$ is not a simple zero, the height formula gives a
torsion genus point, killed by $L$, and its coefficient is zero
modulo $2^M$. Combined with the first inequality on simple-zero
parameters, this gives the all-$D$ truncated bound used below,
without assigning a rank-one lattice index to a torsion point.
\end{proof}

\subsection{Integer divisions and actual trace operators}
\begin{proposition}[Arithmetic realization of the coefficient rules]\label{ex:app-coefficient-rules}
Use the rational full-two-torsion graph conventions of
Definitions~\ref{ex:sec7-typed-graph} and~\ref{ex:sec7-symbol}.
For an even source of Proposition~\ref{ex:app-even-phase}, use its
fixed minimum depth. For an odd source, first fix a finite operator
budget $|U|\le W$, take $R\ge2$, and choose an actual stage satisfying
the depth inequalities of Proposition~\ref{ex:app-odd-minimum}.
Then the integer divisions (O.9)--(O.11) are valid at every Fourier
index, with the stated truncation for the odd source.
The fresh coefficient-one reading is
$\overline{(t(\operatorname{Fr}_q)\theta H_U)[1]}=\mathcal A(Uq)$.
On all tested nonempty configurations, these actual symbols depend
only on their typed graph and satisfy the rules of
Lemma~\ref{ex:sec7-coefficient-rules}: a pair difference
contracts the endpoints and adds their types and every external
incidence; whole-prime omission uses the same source, phase, depth
and stage; and an isolated primary vertex of neutral type and
neutral external incidences makes the symbol zero.
Every contraction has nonempty output. Stationary character tests
are carried along in the same product rule. This statement asserts
neither values at untested odd-stage weights nor an identity at an
empty contraction output.
\end{proposition}
\begin{proof}\leavevmode
All new primes below avoid the full coefficient support $S$,
including every late sieve and the current $\ell_0$.
Here $S$ is the actual coefficient support.  The coefficients of
$g$ vanish at all indices not prime to $S$.
For the odd source this is forced at $2,N,h_*,\ell_0$ by the
local Schwartz support and the extraction at $|h_*|$.
At every additional coefficient-sieve prime it is imposed by an
exact zero-or-retain Fourier sieve, which commutes with the good
recurrence outside the enlarged support.  This is the vanishing
used in the condition $(m,S)=1$ in the unary series below.
Later stationary prime-selection tests may enrich the common
label alphabet without resieving the intrinsic odd source;
such tests do not enlarge this coefficient-support condition.
Use the minimum-depth inequalities proved in Appendix~\ref{app:endpoints}.
The general bookkeeping assigns weight one to a split prime of
$\Q(E[2])$ and one-half to a prime with transposition Frobenius, called
simple. In our actual full-two-torsion family every prime is split,
$w(U)=|U|$, and the depth $\delta$ is an integer. Thus only matching
parity, meaning $e_U=0$, is needed here; the displayed general formula
also records where the unused simple-prime term would enter.
For a finite set $U$, put
\[
 n_U=\lceil\delta+w(U)\rceil,\quad e_U=n_U-\delta-w(U),\qquad
 B_U=2^{-n_U}\prod_{p\in U}(U_p-\eta(p)V_p)g,
\]
where $U_p(\sum a(n)q^n)=\sum a(pn)q^n$ and $V_pg(\tau)=g(p\tau)$.
These are operations on the actual form $g$. We verify the level
at a new odd prime using the classical theta multiplier. Write the
old level as $4A$, with $p\nmid4A$, and put
$\beta=\left(\begin{smallmatrix}p&0\\0&1\end{smallmatrix}\right)$.
For $\gamma=\left(\begin{smallmatrix}a&b\\c&d\end{smallmatrix}\right)
\in\Gamma_0(4Ap)$ one has
\[
 \beta\gamma\beta^{-1}=
       \begin{pmatrix}a&pb\\c/p&d\end{pmatrix}\in\Gamma_0(4A).
\]
The square-root factor in the theta multiplier is unchanged on
evaluating the latter matrix at $p\tau$. The quotient of its
arithmetic factor by that of $\gamma$ is
$(c/p\,/d)/(c/d)=(p/d)$. Cubing the multiplier leaves this
quadratic factor unchanged. Hence $V_pg$ has level $4Ap$ and
character $\psi(p/\cdot)$.

At this level the normalized sum
\[
 U_{p^2}h(\tau)=\frac1{p^2}\sum_{r\bmod p^2}
                             h((\tau+r)/p^2)
\]
preserves the same weight and character. Indeed, put
$\alpha_r=\left(\begin{smallmatrix}1&r\\0&p^2\end{smallmatrix}\right)$.
For $\gamma\in\Gamma_0(4Ap)$ choose
$r'\equiv(b+rd)/(a+rc)\pmod{p^2}$. Then
\[
 \alpha_r\gamma=\gamma'_r\alpha_{r'},\qquad
 \gamma'_r=
 \begin{pmatrix}
  a+rc&(b+rd-(a+rc)r')/p^2\\
  p^2c&d-cr'
 \end{pmatrix}\in\Gamma_0(4Ap).
\]
The map $r\mapsto r'$ is a permutation, since $a+rc$ is a unit
modulo $p$ and the fractional linear map has determinant one.
Writing $d'=d-cr'$, its nebentypus factor is unchanged since
$d'\equiv d\pmod{4Ap}$. Its theta arithmetic factor is unchanged
as well: $d'\equiv d\pmod c$, $d'\equiv d\pmod4$, and
$(p^2c/d')=(c/d')=(c/d)$. These are the usual Kronecker-symbol
identities for the theta multiplier on $\Gamma_0(4)$.
Its square-root factor, evaluated at $(\tau+r')/p^2$, is
$(c\tau+d)^{1/2}$. Thus the sum transforms with the same
theta multiplier cubed and the same character. The formulas also
hold for the triangular matrices, and holomorphy at the cusps is
preserved by the finite rational translates.
On Fourier expansions $U_{p^2}V_p=U_p$. Consequently both terms
of $U_p-\eta(p)V_p$ have level exponent at most one at $p$ and
character $\psi(p/\cdot)$. Iteration introduces no higher exponent
at previously added primes.

The same computation applies to the unary corrections. Start with
the odd-character unary theta series at the old support, including
its fixed coprimality sieve. Its $J$-series is its transform by
$V_J$, for squarefree $J\subset U$. A prime $p\mid J$ therefore
has exponent at most one and the extra character $(p/\cdot)$;
a prime $p\nmid J$ is absent from its original level. Multiplication
by the ordinary theta series introduces no odd level prime. This
proves the precise level and oldform alternatives used in (O.12).


Write an index as $Jrm^2$ with $J\subset U$ and $r$ squarefree,
prime to $U$. Products of prime sets are understood here.
The originating squarefree parameter in the coefficient of $g$ is
$D=(U\setminus J)r$, not the squarefree part $Jr$ of the tested index.
Remove the phase $\eta(p)^s$ in (O.6) and put
\[
 b_{-1}=0,\quad b_0=1,\quad b_1=a_p-\kappa,\quad
 b_{s+1}=a_pb_s-pb_{s-1},\quad
 \kappa=(\psi/\eta)(p)(-D/p).
\]
For $p\in U\setminus J$ the factor $b_s-b_{s-1}$ is odd;
for $p\in J$ the factor $b_{s+1}-b_s$ is even.
Consequently, for the even source,
\[
 v_2B_U[Jrm^2]\ge w(r)
              +\tfrac12\#\{p\in J:p\text{ is simple}\}-e_U.      \tag{O.9}
\]
The right side is an integer on the allowed filter and is
nonnegative. In our full-two-torsion application all primes are
split and $w=\omega$.
For the odd source, fix an integer $R\ge1$ and a weight bound $W$
first. At one sufficiently advanced actual stage, simultaneously for
all $U$ with $w(U)\le W$ and every Fourier index, (O.9) holds with
its right side replaced by its minimum with $R$. Thus the untruncated
bound on a prescribed bounded range of squarefree weights follows by
choosing $R$ and then the stage. For the unary-removal and two-trace
tests below, choose $R\ge2$ after fixing the operator budget.
For a split prime $a_p\equiv p+1\pmod4$. The recurrence gives the
same residue on each divided shell,
\[
        (b_{s+1}-b_s)/2\equiv
     c_p(D):=\{a_p-(\psi/\eta)(p)(-D/p)-1\}/2\pmod2.              \tag{O.10}
\]

\paragraph{The two integral shell calculations.}
Here every prime is split in the rational two-torsion field, so
$a_p\equiv p+1\pmod4$. For $p\mid D$ one has $\kappa=0$.
Reducing the recurrence modulo two gives $b_{s+1}=b_{s-1}$,
with $b_0=1$ and $b_1=0$. Thus $b_s-b_{s-1}$ is odd for
all $s\ge0$, including $s=0$ with $b_{-1}=0$.
For $p\nmid D$, one has $\kappa\in\{1,-1\}$, so every $b_s$
is odd. Write $d_s=b_{s+1}-b_s$. The recurrence modulo four gives
$d_s\equiv p d_{s-1}\pmod4$ for $s\ge1$.
Each $d_s$ is even, and multiplication of an even residue modulo
four by an odd number fixes that residue. Consequently
\[
 d_s\equiv d_0=a_p-\kappa-1\pmod4\qquad(s\ge0).
\]
This proves (O.10) without dividing a residue-field identity by two.
For an index $Jrm^2$, each $p\in J$ therefore supplies an actual
factor two, whereas the factors for $p\in U\setminus J$ are integral
units modulo two. The squarefree input has
$\omega(D)=|U|-|J|+\omega(r)$. Its depth bound plus these $|J|$
factors gives
$\delta+\omega(D)+|J|=n_U+\omega(r)$ before division by
$2^{n_U}$. All removed character factors are units.
For the odd source the same argument is a congruence modulo
$2^M\mathcal O_{\rm coef}$ before that division. It therefore
retains precision $M-n_U$ afterward; the subsequent division
in (O.11) retains precision $M-n_U-1$.
Choose $i\in\{0,1\}$ with $\eta\chi_{-4}^i$ odd and let
\[
 \Theta_J=\sum_{\substack{m>0\\(m,S)=1}}
                      \eta(m)\chi_{-4}(m)^i m q^{Jm^2}.
\]
Poisson summation gives this genuine unary weight-$3/2$ form.
Its scalar character is $\eta\chi_{-4}^i(-J/\cdot)$ and every
$p\mid J$ enters the level at exponent one.
In matching parity the shell recurrence gives
$B_U[Jm^2]-B_U[J]\eta(m)\in2\mathcal O_{\rm coef}$.
Thus
\[
                 H_U=\frac{B_U-\sum_JB_U[J]\Theta_J}{2}          \tag{O.11}
\]
is integral; in opposite parity take $H_U=B_U$.
This is division by the integer two, not reduction modulo the
maximal ideal of a ramified coefficient ring. For odd-stage
forms, the coarse bound and bounded-weight stabilization proved
below pay every fixed denominator before choosing the stage.

Write $\mathfrak m_{\rm coef}$ for the maximal ideal of the chosen
coefficient valuation ring $\mathcal O_{\rm coef}$. In the following
trace argument a reduced expansion means coefficientwise reduction
modulo $\mathfrak m_{\rm coef}$. The preceding divisions by integer
powers of two have already been justified in $2^a\mathcal O_{\rm coef}$;
reduction modulo the maximal ideal alone would not justify them.
No saturation of the integral Fourier module is assumed, and no
identification of its abstract quotient with all reduced expansions
is needed.
Multiply (O.11) and the undivided corrections by
$\theta=1+2\sum_{m>0}q^{m^2}$. They are actual weight-two forms.
For a fixed stage and $U$, choose a finite complete coefficient-field
extension containing all constituents and retain their old
multiplicities and diamond translates in the finite-dimensional
characteristic-zero space. Use modular Galois representations on
cuspidal constituents and the two character summands on Eisenstein
constituents. Holomorphy alone does not exclude the latter.
Diamond translates remain integral:
their correction coefficients have the form
\[
 \tfrac12 B_U[J]\{\psi_U(u)-\eta(u)\chi_{-4}(u)^i(-J/u)\},
\]
whose character ratio is quadratic. Good Hecke operations preserve
integrality by
$(T_qf)[n]=f[qn]+q(\langle q\rangle f)[n/q]$.
Finitely many Fourier coordinates inject this space into a finite
coordinate space. The module generated by these translates is
therefore finite over the coefficient DVR, and closed.
Chebotarev and continuity make scalar trace and determinant
operators limits of good Hecke and diamond operations, so they
preserve this module. No integral eigenspace projector is used.

\paragraph{Reading an added vertex.}
For the actual full-two-torsion family, all primes are split,
$\delta\in\mathbb Z$, and $n_U=\delta+|U|$.  Put $F_U=\theta H_U$
and let $q\notin S\cup U$ be a fresh prime, where $S$ includes
every late sieve and the current moving theta prime.
Write $u=\prod_{p\in U}p$ and let $\mathcal A(D)$ denote the
normalized valuation symbol on the tested range: it is one exactly when
$v_2a(D)=\delta+\omega(D)$ on the filter. For an odd source, that
range has a fixed weight bound chosen before the valid actual stage,
as specified before (F.7). No finite-stage value at larger weights is
asserted; the compatible limit tables used later are a separate object.
Expanding the operators, every term containing a $V_p$ has zero
coefficient at $q$, and no unary index $Jm^2$ with $J\subset U$
equals $q$.  Therefore
\[
 B_U[q]=2^{-n_U}a(uq),\qquad
 H_U[q]=2^{-n_U-1}a(uq).
\]
At coefficient one the diamond term in the good Hecke formula
vanishes.  Since $F_U\equiv H_U\pmod{2\mathcal O_{\rm coef}}$,
\[
 \overline{(t(\operatorname{Fr}_q)F_U)[1]}
   =\overline{H_U[q]}
   =\overline{\,2^{-\delta-|U|-1}a(uq)\,}
   =\mathcal A(uq)\in\mathbb F_2.
\]
Here residues lie in the coefficient residue field, with
$\mathbb F_2$ its prime subfield.  Every permitted phase is a
two-power root of unity and reduces to one.  For an odd source,
fix the finite operator budget $|U|\le W$ and $R\ge2$ first;
the stage choice in (O.9) may be taken with
$M_\theta\ge\max(0,\delta+W)+R$.
Thus $M_\theta-n_U-1\ge1$, so the normalized finite-stage error
is still divisible by the integer two.  The same reading holds
with value zero off the fixed or moving coefficient support.

The diamond correction above is a sum of integral multiples of
$\theta\Theta_J$; its fresh coefficient $q$ is zero after reduction.
It introduces neither an extra phase nor a denominator in this
reading.  All readings and subsequent omissions use this same
$g,\eta,\psi,\delta$ and actual stage; only the set $U$ changes.

For any fixed Fourier index $n$, Chebotarev approximation lets one
test a finite product of trace operators using distinct fresh primes
$q_1,\ldots,q_a$ avoiding $n$ and the level. The good Hecke formula
then reads the coefficient at $nq_1\cdots q_a$ of the original vector.
Different indices may use different fresh primes; the conclusion is
an identity of coefficientwise reduced expansions.

The determinant identity uses the actual diamond formula, not merely
the reduction of its eigenvalues. For a good prime $r$,
$d(\operatorname{Fr}_r)=r\langle r\rangle$. The formula above shows
that $(d(\operatorname{Fr}_r)-1)F_U$ reduces to a sum of integral
multiples of the undivided corrections $\theta B_U[J]\Theta_J$:
the scalar $r\psi_U(r)-1$ lies in $\mathfrak m_{\rm coef}$.
A fresh trace kills every such correction. Indeed, to test at $n$
choose $q\nmid nJ$; then $nq$ cannot have squarefree part $J$, and
$\theta\equiv1\pmod{\mathfrak m_{\rm coef}}$. This remains true
after further fresh traces. Continuity and the closed integral
module therefore give, on the reduced trace orbit of $F_U$,
\[
 d(u)t(v)=t(v),\qquad
 t(u)t(v)=t(uv)+t(uv^{-1}),\qquad t(1)=0.
\]
The middle equality follows from the characteristic-zero matrix
identity
$\operatorname{tr}(A)\operatorname{tr}(B)
=\operatorname{tr}(AB)+\det(B)\operatorname{tr}(AB^{-1})$,
using the determinant identity just proved; the last follows from
$\operatorname{tr}(1)=2$.

Two fresh primes outside $S\cup U$ contribute two to the squarefree
weight in (O.9). Thus $B_U[nq_1q_2]$ is divisible by four, no unary
correction occurs there, and $H_U[nq_1q_2]$ remains divisible by two.
Here $R\ge2$ and the actual-stage budget were fixed before the test.
It follows that $t(g)t(h)F_U=0$ after reduction. Applying the middle
identity to $t(gh)t(h)$ gives $t(gh^2)F_U=t(g)F_U$.
Invariance under all squares, and continuity, give factorization
through the elementary quadratic quotient unramified outside
$S\cup U$. Over $\Q$ this quotient is finite: its squareclasses
are generated by $-1$, $2$, and the primes in that support.
This does not assert abelianity of the characteristic-zero
Galois representations.

At $p\in U$, first keep the summands of $F_U$ separately.
The block $\theta B_U$ and each unary block with $p\in J$ have
level exponent at most one and primitive quadratic character at $p$.
The unary blocks with $p\notin J$ are degeneracies of forms at a
level omitting $p$. This distinction is preserved when retaining all
old multiplicities. In particular, the latter blocks cannot acquire a
Steinberg constituent merely by being viewed at a larger level.

Enlarge the characteristic-zero coefficient field to realize all the
cuspidal newform constituents and the Eisenstein character pairs of
these blocks. On a constituent and all its degeneracies, define $t(g)$
and $d(g)$ to be the scalar trace and determinant of its two-dimensional
Galois representation, acting as the identity on its multiplicity
space. Distinct characteristic-zero constituents with the same good
eigenvalues have the same trace and determinant by Chebotarev, so
these definitions agree on coinciding summands. In particular they
give well-defined operators on their sum. No decomposition modulo
two, and no residual irreducibility assertion, enters this construction.

For a cuspidal eigenform $f$, let $\Pi_f$ denote its unitary
adelic normalization: at a good prime $q$ its Satake parameters
$\alpha_q,\beta_q$ satisfy
$a_q=\sqrt q(\alpha_q+\beta_q)$ and
$\alpha_q\beta_q=\chi(q)$.
Apply \cite[\S\S0.3--0.7, Theorem~A, pp.~409--411]{carayol}
to $\Pi_f^\vee$, and set $\rho_f=\sigma(\Pi_f^\vee)^\vee$.
Here $F=\Q$, the real weight is $2$, and the coefficient place
lies above $2$. The degree of $F$ is odd, so the additional
finite discrete-series-place hypothesis for even degree does
not apply. The local prime $p$ is odd.
Carayol's reciprocity sends geometric Frobenius to a uniformizer,
and his $\omega_p$ is $|\cdot|_p$.
His normalized-induction formula is
\[
 \sigma(\operatorname{Ind}(\xi_1,\xi_2))
 =\xi_1^{-1}\omega_p^{-1/2}\oplus
   \xi_2^{-1}\omega_p^{-1/2}.
\]
Thus if $\Pi_{f,p}=\operatorname{Ind}(\mu_1,\mu_2)$,
$\rho_f$ on arithmetic Frobenius has eigenvalues
$\sqrt p\,\mu_1(p),\sqrt p\,\mu_2(p)$ whenever these
characters are unramified. In particular its good polynomial
is $X^2-a_qX+q\chi(q)$. Both the automorphic contragredient and
the dual on the Galois side are part of this convention.

For a block with ramified character, its local conductor is
exactly one: it is at most the level exponent and at least the
conductor of its determinant.
Casselman's newvector theorem and conductor calculation
\cite[Theorem~1, pp.~302--306]{casselman} show that a special
representation has conductor one only for an unramified twist;
a ramified twist has conductor at least two. A supercuspidal
representation has conductor at least two. Consequently
$\Pi_{f,p}=\operatorname{Ind}(\mu_1,\mu_2)$ with $\mu_1$
unramified and $\mu_2$ of conductor one, after interchanging
the characters if necessary. Its newvector space is a line.
On inertia the ramified character is the prescribed quadratic
character.

We check the Fourier normalization of the operator on this line.
Use normalized induction, so that
\[
 v\left(\begin{pmatrix}a&x\\0&d\end{pmatrix}g\right)
 =\mu_1(a)\mu_2(d)|a/d|_p^{1/2}v(g).
\]
Put $K_0(p)=\{g\in\mathrm{GL}_2(\Z_p):g_{21}\in p\Z_p\}$.
Use its character $k\mapsto\mu_2(k_{22})$.
This is the lower-right-character convention. Casselman's
upper-left-character convention is carried to it by
$h_p=\left(\begin{smallmatrix}0&1\\-p&0\end{smallmatrix}\right)$,
which normalizes $K_0(p)$ and interchanges its diagonal entries
under conjugation. We label the inducing characters so that
$\mu_1$ is unramified, using the interchange isomorphism for
irreducible principal series.
Normalize Haar measure by $\operatorname{vol}(K_0(p))=1$.
The character-weighted double-coset operator for
$K_0(p)\operatorname{diag}(p,1)K_0(p)$, with value one at
$\operatorname{diag}(p,1)$, is
\[
 H_p=\sum_{b\in\Z/p\Z}
       \Pi_{f,p}\left(\begin{pmatrix}p&b\\0&1\end{pmatrix}\right).
\]
Take the newvector supported on $B K_0(p)$ with $v(1)=1$.
Evaluating at $1$ gives
\[
 (H_pv)(1)=p\mu_1(p)p^{-1/2}=\sqrt p\,\mu_1(p).
\]
This is the Fourier operator $U_p^{(2)}$. Indeed, in the
Whittaker model use an additive character of conductor $\Z_p$.
Writing $a(n)=\operatorname{diag}(n,1)$ and
$n(x)=\left(\begin{smallmatrix}1&x\\0&1\end{smallmatrix}\right)$,
\[
 a(n)\begin{pmatrix}p&b\\0&1\end{pmatrix}=n(nb)a(np)
 \quad(n\in\Z_p),\qquad
 (H_pW)(a(n))=pW(a(np)).
\]
Normalize the archimedean Whittaker function by
$W_\infty(a(y))=y\exp(-2\pi y)$ for $y>0$.
Comparison of its positive-$n$ term with $y f(iy)$ gives the
finite Whittaker coefficient $a_n/n$. The displayed equality
therefore sends $a_n$ to $a_{pn}$. Hence
$U_p^{(2)}=\sqrt p\,\mu_1(p)$ on the new line, exactly the
arithmetic-Frobenius eigenvalue on the unramified Galois line.
This scalar is algebraic; the square roots only express the
unitary normalization. Old multiplicities at primes other than
$p$ carry the same scalar. Casselman's oldspace dimension formula
\cite[Corollary, p.~306, and Theorem~4, p.~313]{casselman}
does not introduce an additional multiplicity at this conductor-one
prime. For Eisenstein constituents, the two Galois characters and
the divisor-sum Fourier coefficients give the same equality directly;
Carayol's cuspidal theorem is not applied to them.

Choose an inertia element $\tau$ on which the quadratic character
is $-1$, and a Frobenius lift $\phi$. On a ramified block, a basis
of the two character lines gives
\[
 \rho(\tau)=\begin{pmatrix}1&0\\0&-1\end{pmatrix},\qquad
 \rho(\phi)=\begin{pmatrix}\lambda&0\\0&\mu\end{pmatrix},
 \qquad U_p^{(2)}=\lambda,\quad d(\phi)=\lambda\mu.
\]
For $\rho(g)=\left(\begin{smallmatrix}a&b\\c&e\end{smallmatrix}\right)$,
the two inertia differences are
\[
 t(g\tau)-t(g)=-2e,\qquad
 t(g\tau\phi)-t(g\phi)=-2e\mu.
\]
It follows in characteristic zero that
\[
 U_p^{(2)}\{t(g\tau\phi)-t(g\phi)\}
       =d(\phi)\{t(g\tau)-t(g)\}.                       \tag{O.12}
\]
On a block coming from level prime to $p$, the representation is
unramified and both differences vanish. This proves the identity
on its entire old multiplicity space even if $U_p^{(2)}$ is not
scalar or semisimple on that space. It therefore proves (O.12)
on the actual sum containing $F_U$.

These are identities on characteristic-zero forms before any
reduction. To relate them to the integral trace module above, choose
finitely many Fourier coordinates that inject its ambient space.
Every generator of that module has integral coordinates; hence the
module is a submodule of a finite free module over the complete
coefficient DVR. It is finite and closed. Simultaneous Chebotarev
approximation in the finitely many Galois representations makes
$t(g)$ and $d(g)$ limits of good Hecke and diamond operators on this
space. They preserve this integral module, so the identity can be
evaluated on the integral $F_U$ and then reduced. This does not
require integral projectors onto constituents. In particular no
division by the two in the displayed inertia differences is used.


Apply (O.12) at coefficient one.  Put $V=U\setminus\{p\}$,
$v=\prod_{r\in V}r$, so $n_U=n_V+1$.  For a fresh prime $r$ set
\[
 C_p(vr)=\frac{a_p-(\psi/\eta)(p)(-vr/p)-1}{2}\in\Z,
 \qquad c_p(vr)=C_p(vr)\bmod2.
\]
The two surviving operator terms at $pr$ and the good recurrence give
\[
 H_U[pr]\equiv
   \eta(p)C_p(vr)\,2^{-n_V-1}a(vr)\pmod{2\mathcal O_{\rm coef}}.
\]
For the even source this is an equality before reduction.  For the
odd source the error is in $2^{M_\theta-n_U-1}\mathcal O_{\rm coef}$,
which is contained in $2\mathcal O_{\rm coef}$ at our chosen stage.
Choose fresh $q,q_\tau,r_0,r_1$ representing respectively
$g,g\tau,g\phi,g\tau\phi$ in one sufficiently precise joint quotient
for the $U$ and $V$ forms and all retained characters.
The omitted form uses the same $g,\eta,\psi,\delta$ and stage,
and its level omits the entire rational prime $p$.
Thus the two $V$ symbols at $vr_0,vr_1$ agree.  Their Legendre
signs at $p$ are opposite, so $c_p(vr_0)+c_p(vr_1)=1$ in $\F_2$.
The coefficient-one reading and the reduced trace identity now give
\[
             \mathcal A_U(Uq_\tau)+\mathcal A_U(Uq)=\mathcal A_V(Vr_0)\qquad\text{in }\F_2.
\]
Here the subscript records the operator set, not a new coefficient
source or minimum. The factor $\eta(p)$ reduces to one; no factor
two is cancelled in (O.12). For every retained quadratic character
$\chi$ unramified at $p$, the representative of the omitted vertex
satisfies
\[
 \chi(g\phi)=\chi(g)\chi(\phi)
            =\chi(\operatorname{Fr}_q)\chi(\operatorname{Fr}_p).
\]
For a fixed-support character this is the product of the two vertex
labels. For the character of a third vertex $r\in V$, it adds the
two bits $(r/q)$ and $(r/p)$ in $\F_2$. Reciprocity then gives the
reverse bits from the same endpoint types. All additional stationary
tests were included in the joint quotient, so they obey this formula
as well. This proves the full product-label omission rule, including
every external incidence.

A fresh vertex with trivial label in the elementary extension for
$S\cup U$ has the same reduced reading as $t(1)$, hence gives zero.
An isolated primary vertex has precisely this fixed label and these
zero incidence bits. A contracted cluster of primary vertices has
the same property when its external incidences are neutral. The
result is still a nonempty configuration, including the one-vertex
case; no contraction deletes the last vertex.

Finally the fixed-$U$ factorization gives dependence on the entire
finite typed graph. For two disjoint prime realizations of the same
graph, choose a fresh bridge copy with the prescribed types and
matching incidences to both realizations. CRT and Dirichlet provide
the primes successively, and reciprocity determines their reverse
bits. Replace the vertices of the first copy by bridge vertices one
at a time, then the bridge vertices by those of the second copy.
At each step the other vertices form the same fixed $U$, and the
changed vertex has the same elementary Frobenius data for $S\cup U$;
its symbol therefore agrees. If the original realizations overlap,
first replace one by a fresh copy avoiding both original prime sets,
with incidences matching that first realization. A vertex and its
own replacement are never tested together, so their mutual bit can
be chosen freely. The resulting two disjoint realizations admit the
preceding bridge argument without conflicting prescriptions.
Only graphs of the original bounded size are evaluated; extra bridge
incidences are prime-selection conditions. Thus these are rules for
the actual arithmetic graph symbols, at the already chosen valid
stage. The simultaneous finite graph budget and stabilization are
treated below.
\end{proof}

% End 08-odd-source.tex


% Begin 09-normalization-endpoints.tex
\section{Nonvanishing, normalization and the final family}\label{app:endpoints}
The uses of the coefficient and address propositions have the following
order. First apply only their even-source case, whose depth is already
fixed by Proposition~\ref{ex:app-even-phase}. This gives the rough
companion and the genus lower bound. These in turn establish the odd
minimum and its finite-stage inequalities. Only then apply the odd
case of Proposition~\ref{ex:app-coefficient-rules} and the two-symbol
address construction. In particular the proof of the rough companion
does not use the later odd minimum.
We give the choice order underlying (3.3)--(3.9). This also separates
the two valuation arguments: the even upper transfer uses the common
Kato scalar, while the final contradiction uses the common Heegner
coordinate of Section~\ref{sec:limits}.

\subsection{An absolute bound of five new primes}
The theorem of Hoffstein--Luo \cite[p.~436]{hl} supplies fundamental
discriminants with nonzero central value and at most four prime
factors. Its finite splitting set is arbitrary. The bound four is
absolute; no uniform size threshold for those discriminants is needed.

\begin{proposition}[Nonvanishing with five primes and finite local prescriptions]\label{ex:app-five-prime}
Let $A/\Q$ be a fixed modular elliptic curve for this application,
let $S$ contain $2,3$ and all its bad primes, and let $B$ be a finite
prime exclusion set. In every nonempty signed unit-squareclass filter
at $S$ of even root number, there is a signed odd squarefree parameter
$u$, prime to $S\cup B$, with $L(A^u,1)\ne0$ and $\omega(u)\le5$.
Only the intermediate parameter supplied by Hoffstein--Luo is required
to be fundamental; the final field discriminant may be $4u$.
The number five is independent of the prescribed finite set. No
uniform bound on the sizes of the selected primes is asserted.
\end{proposition}
\begin{proof}\leavevmode
Let $A/\Q$ be any fixed modular elliptic curve for this application,
$S$ contain two, three and all its bad primes, and $B$ be any finite
exclusion set. Fix a nonempty signed unit filter of even root number.
By CRT and Dirichlet choose one signed prime $q_0$ in its classes,
outside $S\cup B$. Apply the theorem to the actual curve $A^{q_0}$,
putting
\[
                 S'=S\cup B\cup\{|q_0|\}\cup
                               \{p:p\mid N_{A^{q_0}}\}
\]
in its splitting set. The theorem gives a fundamental $d$ with
$\chi_d(p)=1$ for all $p\in S'$ and at most four prime factors.
The functional equation forces $d>0$, since $A^{q_0}$ has even
root number and every conductor prime splits. Splitting at two
gives $d\equiv1\pmod8$. Thus $q_0d$ has all original unit
classes, is prime to $B$, is squarefree, and has at most five
prime factors. If $q_0d\equiv3$ or $7\pmod8$, its field
discriminant is $4q_0d$; only $d$ was required to be fundamental
in the theorem. This is compatible with the coefficient
construction (E.2).
\end{proof}

Here $h$ is restricted to the prepared positive odd-sign unit
filter containing one, as in (3.3), so $h\equiv1\pmod8$.
The multiset of dummy labels and its size are fixed before $h$;
its realizing primes are chosen after $h$ and avoid its support.

For the rough companion the curve itself varies:
write $a$ for the bounded dummy product and take $A=E^{ha}$.
Choose $q_0$ so $k_0=aq_0$ is negative, $1$ modulo eight and
a square at every prime of $2Nh$, with the required exclusions.
Then $k_0$ is odd fundamental, prime to $N_{E^h}$, and
\[
          w(E^{hk_0})=w(E^h)\chi_{k_0}(-N_{E^h})=+1.
\]
This proves the sign directly for the actual application.
Apply Hoffstein--Luo to $E^{hk_0}$ with every conductor,
old, $h$, dummy and forbidden prime in its splitting set.
The resulting $k'=k_0d$ adds at most five new primes.
The constants and sizes in the analytic theorem may depend on
this varying curve and the whole finite set; the bound five does not.

\subsection{The even upper transfer}
\begin{proposition}[A uniformly bounded partitionable companion]\label{ex:app-rough-companion}
Fix $E$ and the finitely many prepared even sources, local types and
unit seeds used for (3.3). There are constants $A,C_E$, fixed before
$h$ and before any additional finite prime exclusion set, with the
following property. Let $h$ be positive squarefree in the prepared
odd-sign unit filter containing $1$, prime to its fixed support; in
particular $h\equiv1\pmod8$. There is a negative odd fundamental
discriminant $k'$, prime to $2Nh$ and the additional exclusions,
square at every place dividing $2Nh$, for which $hk'$ is partitionable
in the sense of Definition~\ref{ex:sec7-seed-partition},
$L_*(hk')\ne0$, and
\[
 \omega(k')\le A,\qquad
 v_2L_*(hk')\le2\omega(h)+2\omega(k')+s(hk')+C_E.
\]
Here $A$ is the fixed dummy-label count plus five. Dummy types and
their number are fixed first, but their distinct prime realizations
are selected after $h$ and the exclusions. The proof uses the
one-symbol stationary-address conclusion of
Lemma~\ref{ex:sec7-one-symbol-address}, whose proof is given in
Proposition~\ref{ex:app-simultaneous-address}.
\end{proposition}
\begin{proof}\leavevmode
For one fixed even source write (E.11) as
\[
             v_2L_*(\sigma D)=2v_2a(D)+\kappa_e.
\]
The lower bound (E.9) and a nonzero reference coefficient imply
that the nonempty set of integers
$v_2a(D)-\omega(D)$ has a finite attained minimum $\delta_e$.
A unit of the normalized symbol means precisely equality to this
minimum. It therefore gives
\[
 L_*(\sigma D)\ne0,\qquad
 v_2L_*(\sigma D)=2\omega(D)+C_{e,\rm unit},
 \quad C_{e,\rm unit}=2\delta_e+\kappa_e.                         \tag{F.1}
\]
This constant belongs to the fixed source, not to a curve later
introduced solely to obtain nonvanishing.

For the negative types used in the rough companion, the unsigned
index is $7$ modulo eight. Its total $(-1/D)$ label is nonzero, so
every unit seed is nonempty. A seed of least positive size has no
dependence on internal edge bits: a nonzero edge derivative
would be a smaller nonempty unit. Its labels therefore suffice.
If a list with the same total label contains the first $r-1$
seed anchors and a nonempty remainder, contract that remainder
into the last anchor. Its total label is forced to be the
remaining seed label. This proves partitionability regardless
of how many old $h$-primes are in the remainder.

For a partitionable $d$ already known to have Selmer corank zero,
the one-symbol address construction keeps every prime of $d$
stationary and gives arbitrarily large cubes $d_x$, $d_0=d$,
with unit symbol at each $x\ne0$. The scalar in (E.5) is common
to that cube. By (E.8) and (F.1),
\[
      v_2u_x^{\rm ev}(0)\le B_e:=C_{e,\rm unit}+C_{\rm loc}
                                      \quad(x\ne0).
\]
Choose $M-c_E>B_e$. The constant coefficient of
$2^{c_E}u^{\rm ev}\in\Lambda[G_b]$ has sign-value bits of total
Boolean degree at most $2^{M-1}-1$. For a larger cube dimension,
the Boolean comparison in Section~\ref{sec:limits} gives
$x\ne0$ with
$v_2(u_x^{\rm ev}(0)-u_0^{\rm ev}(0))\ge M-c_E$.
If the base were zero or had valuation greater than $B_e$, this
would contradict the upper bound at $x$. Therefore
\[
       v_2L_*(d)-2\omega(d)-s(d)
                       \le C_{e,\rm unit}+2C_{\rm loc}.          \tag{F.2}
\]
The required dimension may depend on $d$; this bound does not.
It precedes the odd minimum and uses no Heegner transfer.

Prepare these sources and nonempty seeds for every total type at
the fixed support. There are finitely many. A bounded supply of
dummy anchors and reserves works for all of them. For each $h$
realize that supply by distinct primes avoiding the entire fixed
support and $h$. Put every dummy and every other already fixed
forbidden prime in the exclusion set of the preceding
nonvanishing construction. It gives $k'$ with
$\omega(k')\le A$, where $A$ is the fixed dummy count plus five.
For the total type that actually occurs, $hk'$ is partitionable.
Nonvanishing gives corank zero by the forward theorem, so (F.2)
applies. Taking the maximum over the finitely many prepared
constants proves
\[
 v_2L_*(hk')\le2\omega(h)+2\omega(k')+s(hk')+C_E.                 \tag{F.3}
\]
No new even minimum is chosen after $h$ or its exclusion set.
Those excluded primes are not added as relaxed places in a
Selmer complex.
\end{proof}

\subsection{The odd minimum is chosen after its lower bound}
\begin{proposition}[The odd minimum and every finite operator budget]\label{ex:app-odd-minimum}
Fix the original odd filter, the detecting sequence of
Proposition~\ref{ex:app-detector}, and one nonprincipal ultrafilter
on its stages. The arithmetic lower bound (3.5), the native comparisons
(3.4), and Proposition~\ref{ex:app-rough-companion} give (F.4)--(F.5)
before any odd minimum is chosen. The bounded-weight visible depths
therefore have an attained finite integer minimum $\delta_o$.
For every fixed weight cutoff, the simultaneous inequality (F.6)
holds on a set in this same ultrafilter. For every finite operator
budget and $R\ge2$, one may choose a stage where all the integer
divisions and residual errors required in
Proposition~\ref{ex:app-coefficient-rules} are valid at every Fourier
index.
On a fixed squarefree-output range the unit symbol defined before
(F.7) is exhaustive, and its value $1$ implies (F.7), with the one
constant $C_o=\delta_o+C_{\rm det}$. No minimum is changed on
omitting a vertex or adding stationary tests, and no single finite
stage is claimed to serve every unbounded weight range.
\end{proposition}
\begin{proof}\leavevmode
The arithmetic result (3.5) and the two native-lattice height
comparisons (3.4) give, using (F.3),
\[
 2j_{\rm gen}(h)\ge2\omega(h)+s(h)-C_{\rm gen},\qquad
 C_{\rm gen}=2A+C_R+C_E-v_2L_*(h_*)+C_H.                         \tag{F.4}
\]
The variable $s(hk')$ cancels. This estimate precedes the theta
prime, precision, odd minimum, final field and binary dimension.

Combining (F.4) with (O.8) gives on all admissible $D$
\[
            v_2a_M(D)\ge\min\{M,\omega(D)-C_{\rm coarse}\},
 \qquad C_{\rm coarse}=C_{\rm gen}/2+C_{\rm det}.                 \tag{F.5}
\]
At visible coefficients (O.8) first gives a simple zero, allowing
(F.4); at the other odd-sign parameters the coefficient is zero
modulo $2^M$. This avoids a circular use of the future minimum.

Fix one nonprincipal ultrafilter on the actual stages with
$M=M_\theta\to\infty$. Consider the integers $r$ for which a
bounded-weight sequence of visible $D_M$ has
$v_2a_M(D_M)-\omega(D_M)=r$ on a set in the ultrafilter.
The fixed witness $D_w$ and (O.2),(O.7) make this set nonempty.
Equation (F.5) bounds it below. Let its attained minimum be
$\delta_o$. For every fixed $W$, on a set in the same ultrafilter,
\[
 v_2a_M(D)\ge\omega(D)+\delta_o
             \quad\hbox{for every visible }D,\ \omega(D)\le W.
 \tag{F.6}
\]
Otherwise choose a violating $D_M$ at each bad stage. One of the
finitely many integers between $-C_{\rm coarse}$ and
$\delta_o-1$ would occur on an ultrafilter-large set, a contradiction.
Invisible coefficients satisfy the required truncated inequality
once $M>W+\delta_o$.

For a fixed finite operator budget $|U|\le W_U$, put
$N_0=\max(0,\delta_o+W_U)$ and fix $R\ge2$.
Choose an integer $T\ge W_U+1$ with
$T-C_{\rm coarse}\ge N_0+R$, then choose a stage in the
ultrafilter-large set where (F.6) holds through weight $T$, with
$M\ge N_0+R$ and $M>T+\delta_o$.
For the index $Jrm^2$ put $D=(U\setminus J)r$.
The integral shell recurrence has an error in $2^M\mathcal O$.
If $\omega(D)>T$, (F.5) gives
\[
 v_2\bigl(2^{n_U}B_U[Jrm^2]\bigr)\ge N_0+R.
\]
If $\omega(D)\le T$, (F.6) for visible coefficients and
$v_2a_M(D)\ge M$ for invisible ones instead give
\[
 v_2\bigl(2^{n_U}B_U[Jrm^2]\bigr)
       \ge\min\{M,n_U+\omega(r)\}.
\]
Consequently $v_2B_U[Jrm^2]\ge\min\{R,\omega(r)\}$ for every
Fourier index and every $U$ in the budget at this one stage.
At $r=1$ the normalized originating coefficient
$2^{|J|-n_U}a_M(U\setminus J)$ is integral, so (O.10) proves
the unary congruence modulo the integer two.
The additional division in (O.11) leaves error valuation at least
$R-1\ge1$.  In particular $M-n_U-1\ge1$ for every reading and
first-shell omission in the budget.
It does not select one stage for infinitely many budgets.

More precisely, fix a squarefree-output cutoff $L$. On a set of
stages in the same ultrafilter, (F.6) holds through $L$ and
$M_\theta>\delta_o+L$. Define $\mathcal A_{M_\theta,L}(D)$ only
for positive squarefree $D$ with $\omega(D)\le L$: it is zero off the full coefficient filter,
one when $v_2a_{M_\theta}(D)=\delta_o+\omega(D)$ on that filter,
and zero otherwise. This is exhaustive on the tested range:
visible coefficients satisfy (F.6), while invisible ones have
valuation at least $M_\theta>\delta_o+\omega(D)$.
In particular a unit in this range is visible and satisfies
\[
 \ord_{s=1}L(E^D,s)=1,\qquad
        j_{\rm gen}(D)\le\omega(D)+C_o,\quad C_o=\delta_o+C_{\rm det}.
 \tag{F.7}
\]
For the preceding operator budget take $L\ge W_U+1$ and include
its valid-stage requirements; all fresh-vertex and omission outputs
then lie in this range. The constant $C_o$ is independent of $L$
and of the stage. No value is assigned by this definition at
unbounded weights of one finite stage. No minimum is changed when
a vertex is omitted or a stationary condition is later added.
\end{proof}

\subsection{Nonempty rational seed and a common finite stage}
\begin{proposition}[Nonempty seeds after fixing the complementary twist]\label{ex:app-nonempty-seeds}
Choose the negative even unit $k$ with the fresh ramification and
exclusions specified in Section~\ref{sec:rc}, and fix
$K=\Q(\sqrt{k})$ before the address dimension. Then $E/K$ has full
two-primary Selmer corank one. Transport the even coefficient source
to $E^k$, normalize at index $1$, and use its positive square-unit
filter. On that filter its coefficients have
$v_2a_{\rm new}(D)\ge\omega(D)$, and the source has a nonempty
$K$-split unit seed.
If the intrinsic odd limit symbol is a unit at the empty
configuration, $E$ already has a simple zero. Otherwise it too has
a nonempty unit seed. In this remaining case both seeds can be
read through the actual common alphabet, whose projections onto
the two intrinsic alphabets are separately surjective. No
surjectivity onto their product is required.
\end{proposition}
\begin{proof}\leavevmode
Choose the final negative even unit $k$ as in
Section~\ref{sec:rc}, with fresh ramification outside the old odd
label and torsion fields. Fix $K=\Q(\sqrt k)$ before the binary
dimension. The forward theorem gives corank zero for $E^k$;
the sum-and-difference isogeny, on compact Selmer modules tensored
with $\Q_2$, gives full corank one for $E/K$.
The intrinsic odd symbol omits the primes of $k$.
Later $\ell_0$ avoids them. Fresh ramification implies that every
old odd label still lifts to a $K$-split label.
Additional stationary tests can have relations; use their actual
common alphabet and its separate surjective projections.

\paragraph{Transport to the fixed complementary twist.}
Put $d_k=|k|$, so that $k=-d_k<0$. The old source has sign $-1$;
its Fourier coefficient at the positive index $d$ tests $E^{-d}$.
Choose the new source for $E^k$ with sign $+1$, normalize its
coefficient at $1$ to be $1$, and impose a support containing the
old support and all primes dividing $k$. For the following comparison
$D$ is positive and squarefree, is prime to that support, and lies
in the square unit classes of $1$ required by the new filter.
In particular $D$ is a square in every old support completion, so
$d_kD$ and $d_k$ lie in the same old filter. Waldspurger comparison
on the two sources gives
\[
 \begin{split}
  v_2 a_{\rm new}(D)
    &=v_2 a_{\rm old}(d_kD)-v_2 a_{\rm old}(d_k)\\
    &\ge \omega(D).
 \end{split}
\]
The equality is a valuation equality; the two sources may have
different root-of-unity phases. The normalized central-value ratios
agree, since both are
\[
 \sqrt D\,\frac{L(E^{kD},1)}{L(E^k,1)}.
\]
Fixed period constants cancel in each ratio. The old unit condition
at $d_k$, and $\omega(d_kD)=\omega(d_k)+\omega(D)$, prove the
inequality. The signed twist is $kD$ throughout; no Fourier
coefficient at a negative index is used. The final family already
satisfies these unit classes. In the nonempty-seed argument the
fresh prime $p'$ is chosen to match all fixed labels of $p$, so
$pp'$ satisfies the same filter, including its dyadic condition.

For a unit of this normalized source equality holds, so
$v_2L_*(h_xk)\le2\omega(h_x)+C_{e,\rm final}$ with one fixed
period-normalization constant.
Its $K$-split seed must be nonempty. Suppose otherwise. The integral
coefficient valuations then improve the preceding lower bound to
$\omega(D)+1$ on every nonempty product of allowed $K$-split primes.
For a fresh good $K$-split prime $p$ outside the support, apply the
trace identity to $f=\theta H_{\{p\}}$:
\[
 U_p^2t(\phi)=U_p^3+d(\phi)U_p,\qquad
 (t(\phi)f)[p^2]=f[p^3]+(d(\phi)f)[p].
\]
Here and in this calculation all congruences are modulo the ideal
$2\mathcal O_{\rm coef}$ before passage to the residue field.
Write $a_p=a_p(E^{(k)})$.
We check every theta convolution term in the three coefficients.
Put
\[
\begin{gathered}
 B=B_{\{p\}}=\frac{(U_p-\eta(p)V_p)g}{2},\qquad
 b_1=B[1]=\frac{a(p)}2,\quad b_p=B[p],\\
 H=H_{\{p\}}=\frac{B-b_1\Theta_1-b_p\Theta_{\{p\}}}{2}.
\end{gathered}
\]
The inherited depth is zero, and $H$ is integral by (O.11).
Also $H[0]=0$, since the cuspidal source and the unary series have
zero constant term.
The ordinary lower bound gives $b_1\in\mathcal O_{\rm coef}$.
Under the assumed absence of a nonempty $K$-split unit seed one
has the stronger assertion
\[
                         b_1\in2\mathcal O_{\rm coef}.          \tag{G.4a}
\]
Indeed, if $p$ lies in the coefficient filter, the improved bound is
$v_2a(p)\ge2$; otherwise the actual filter gives $a(p)=0$.
No condition that $p$ itself belong to that filter is imposed.

Let $\delta_B,\delta_1,\delta_p$ be the scalars by which $d(\phi)$
acts on the full weight-two blocks
$\theta B,\theta\Theta_1,\theta\Theta_{\{p\}}$, respectively.
They are integral units.  The determinant is the cyclotomic character
times the weight-two nebentypus on every cuspidal constituent and
on the Eisenstein character pair, independently of old multiplicity.
For this positive-twist source $\psi/\eta=\chi_{-1}$ and $i=0$.
The half-integral-weight characters of $B$ and $\Theta_{\{p\}}$
are therefore both
\[
                 \psi(p/\cdot)=\eta(-p/\cdot).
\]
Multiplication by the same $\theta$ preserves this equality.
Consequently $\delta_B=\delta_p$ for the actual Frobenius lift
$\phi$, including its action on the ramified character.  The ratio
of the characters of $\theta\Theta_{\{p\}}$ and
$\theta\Theta_1$ is quadratic.  Their cyclotomic determinant factors
are identical, so $\delta_p/\delta_1\in\{1,-1\}$ and
$\delta_p-\delta_1\in2\mathcal O_{\rm coef}$.  We must not simply
equate these two scalars.  On the actual characteristic-zero
module the exact identity is
\[
 d(\phi)f=\delta_p f+
       \frac{(\delta_p-\delta_1)b_1}{2}\,\theta\Theta_1.
\]
Now $H[p]=0$, since $B[p]=b_p$, $\Theta_{\{p\}}[p]=1$ and
$\Theta_1[p]=0$.  Integrality of $H$ implies $f[p]\in2\mathcal O_{\rm coef}$.
Also $\theta\Theta_1[p]=2\sum_{0<j^2<p}\Theta_1[p-j^2]$ is even.
The displayed determinant identity therefore already proves
$(d(\phi)f)[p]\in2\mathcal O_{\rm coef}$ using only the ordinary
integrality of $b_1$.  The following direct expansion checks the
individual cross terms under the no-seed assumption in force here.

For $n_j=p-j^2>0$, one has $p\nmid n_j$ and
$\Theta_{\{p\}}[n_j]=0$.  We claim that
\[
                           B[n_j]\in2\mathcal O_{\rm coef}.    \tag{G.4b}
\]
If $n_j$ meets the old support, the coprimality sieve makes this
coefficient zero.  In particular $2\in S$, so every odd $j$ gives
an even $n_j$ and vanishes in this way.  Otherwise write
$n_j=r m^2$, with $r$ squarefree.
Since $p\nmid n_j$, the $V_p$ term is zero and
$B[n_j]=a(prm^2)/2$.  The integral good-prime shell recurrences
express $a(prm^2)$ as an integral multiple, up to unit character
factors, of $a(pr)$.  If $r>1$, the ordinary coefficient lower
bound gives $v_2a(pr)\ge\omega(pr)=1+\omega(r)\ge2$, or the
coefficient is zero off the filter.  This uses no assertion that
the primes dividing $r$ split in $K$.  If $r=1$, the same shell
recurrences express $B[m^2]$ as an integral multiple of $b_1$,
and (G.4a) applies.  This proves (G.4b).  Also every coefficient
of $b_1\Theta_1$ is even by (G.4a).

Now expand $\theta=1+2\sum_{j>0}q^{j^2}$ before discarding any
term.  Linearity of the determinant operator on the three actual
blocks gives the exact expression
\[
\begin{split}
 (d(\phi)f)[p]
 ={}&\frac{\delta_B B[p]-\delta_p b_p}{2}\\
 &+\sum_{0<j^2<p}
 \bigl(\delta_B B[p-j^2]
       -\delta_1 b_1\Theta_1[p-j^2]
       -\delta_p b_p\Theta_{\{p\}}[p-j^2]\bigr).
\end{split}                                                     \tag{G.4c}
\]
There is no constant-theta $\Theta_1[p]$ term since $p$ is not
a square, and $\Theta_{\{p\}}[p]=1$.
The first line is zero because $B[p]=b_p$ and
$\delta_B=\delta_p$.  Every summand in the second line is even
by (G.4a)--(G.4b) and the vanishing of its last term.  In particular
representations $p=j^2+m^2$, which can occur when $p\equiv1\pmod4$,
are covered by the case $r=1$ and are not omitted.  Thus
$(d(\phi)f)[p]\in2\mathcal O_{\rm coef}$.

The other two readings have no determinant operator.  Since $H$
is integral, for every positive integer $N$,
\[
 f[N]=H[N]+2\sum_{0<j^2\le N}H[N-j^2]
                 \equiv H[N]\pmod{2\mathcal O_{\rm coef}}.     \tag{G.4d}
\]
This includes $N=p^3$ and every term in its theta convolution.
Choose a fresh $K$-split prime $p'$ matching $\phi$ in a sufficiently
precise finite Galois quotient for the retained integral module,
and matching all fixed rational characters.  The good Hecke formula
at $p'\ne p$ identifies $(t(\phi)f)[p^2]$ modulo two with
$f[p^2p']$, hence with $H[p^2p']$ by (G.4d).
Neither unary series has support at $p^2p'$, and the recurrence at
$p\mid pp'$ has $\kappa=0$.  Consequently
\[
 H[p^2p']=\frac{a(p^3p')-\eta(p)a(pp')}{4}
        =\frac{\eta(p)(a_p-1)a(pp')}{4}.
\]
The product $pp'$ lies in the square filter even if $p$ alone does
not.  Its two primes split in $K$, so the no-seed assumption gives
$v_2a(pp')\ge3$.  This last coefficient is therefore even.
The trace identity and (G.4c)--(G.4d) now imply
$H[p^3]\in2\mathcal O_{\rm coef}$.  The following computation
uses precisely this unary-subtracted coefficient.

\paragraph{The coefficient behind the mod-eight contradiction.}
Write $a_p=a_p(E^k)$. The transported source tests positive twists of $E^k$. Thus
$\psi/\eta=\chi_{-1}$ and, at the squarefree index $1$,
$\kappa=(\psi/\eta)(p)(-1/p)=1$. Also $\eta$ is odd, so
the unary correction uses $i=0$. With $a(1)=1$ the recurrence gives
\[
 a(p^2)=\eta(p)(a_p-1),\qquad
 a(p^4)=\eta(p)^2\{a_p(a_p-1)-p\}.
\]
The inherited depth is zero. Hence
\[
 B_{\{p\}}[p]=\frac{\eta(p)(a_p-2)}2,\qquad
 B_{\{p\}}[p^3]=\frac{\eta(p)^2((a_p-1)^2-p)}2.
\]
Only $\Theta_{\{p\}}$ contributes at $p^3$, with coefficient
$\eta(p)p$. Therefore the actual integer division gives
\[
 H_{\{p\}}[p^3]
   =\frac{\eta(p)^2(a_p-1)(a_p-p-1)}4.
\]
This expression is integral because full rational two-torsion gives
$a_p\equiv p+1\pmod4$. The positive-twist character convention is
essential in this computation. Under the assumed absence of a
nonempty unit seed, the trace identity makes this coefficient zero
in the residue field. Since its phase is a unit and $a_p-1$ is odd,
it follows that $a_p\equiv p+1\pmod8$. The contradiction is obtained
only after this conditional deduction.
Choose $p$ also with $(k/p)=1$, $p\equiv1\pmod4$ and
$(\ell/p)=-1$, outside the finite support.  These prescriptions
are compatible: $K$ has a ramified prime outside the ramification
of $\mathbf Q(i,\sqrt\ell)$, so
$K\cap\mathbf Q(i,\sqrt\ell)=\mathbf Q$.  The corresponding
quadratic residue classes are therefore nonempty and Dirichlet's
theorem supplies infinitely many such primes.  Since $k$ is a
square modulo $p$, the reductions of $E^{(k)}$ and $E_\ell$ are
isomorphic.  A half of a nonzero two-torsion point with root $e_i$
on a split cubic requires both differences $e_i-e_j$ to be
squares.  For the roots $0,\ell,-3\ell$, one of these required
differences is, respectively, $-\ell,\ell,-4\ell$.  All three are
nonsquares under the chosen conditions.  Thus none of the three
nonzero points of order two has a half over $\mathbf F_p$.
If the two-primary subgroup had order greater than four, its
structure as a finite abelian group with exactly four elements
killed by two would give an element of order four, and hence such
a half.  Its order is consequently exactly four.  This gives
$v_2(p+1-a_p)=2$, contradicting the congruence modulo eight above.
This proves the required nonempty even seed. If the intrinsic odd
symbol already has a unit at the empty configuration, its detection
formula gives a simple zero for $E$ and the converse is finished.
Otherwise take a least nonempty odd unit seed. Thus the remaining
two-symbol address argument has two nonempty seeds.
\end{proof}

\begin{proposition}[Compatible finite graph tables on one ultrafilter]\label{ex:app-odd-stabilization}
Use $\delta_o$ and the ultrafilter of
Proposition~\ref{ex:app-odd-minimum}. First fix the abstract type of
the bounded elementary label groups and their projection data on
a set in this same ultrafilter. For every finite cutoff $W$,
the valid odd stages give finitely many possible symbol tables on
the compatible typed graphs with at most $W+1$ vertices.
One table occurs on a set in that ultrafilter. These selected tables
agree on their common restrictions, and their union defines the
intrinsic odd symbol on all compatible finite typed graphs.
It retains its bounded-weight witness and the nonempty-output
coefficient rules. After fixing $k$, its stationary tests may be
included in a common alphabet with separately surjective
projections on the same ultrafilter, without changing $\delta_o$.
Every specified finite collection of tests has a common actual
stage. No stage common to all cutoffs is selected.
\end{proposition}
\begin{proof}\leavevmode
Here the finite graph argument has a definite arithmetic order.
The moving theta prime contributes only a bounded number of quadratic
label coordinates. The root characters at the old level were fixed
before $M_\theta$; growing Kummer precision introduces no growing-order
character on the variable primes. Thus the elementary label groups,
their projections and their finite addition tables have bounded size.
Choose their abstract type on the same ultrafilter as $\delta_o$.
For each fixed weight cutoff $W$, restrict first to the
ultrafilter-large set of stages where the preceding budget and
precision conditions hold for every $|U|\le W$, with $R=2$.
On those stages the coefficient reading, factorization, contraction
and isolation rules hold on all configurations of at most $W+1$
vertices, with the stated nonempty-output restriction.  They therefore define actual symbol truth tables on
the compatible typed graph data.  There are finitely many such
tables; select the one occurring on a set in the same ultrafilter.
Two selected tables agree on their common restriction by taking
the finite intersection of their valid-stage sets. Their union defines
the odd limit symbol on compatible finite typed graphs of every size;
it is not a total valuation symbol at any one finite stage.
Every seed or contraction identity uses only finitely many such graphs
and is inherited at a cutoff containing all of them.
The bounded-weight witness remains a witness in these tables.
No stage common to all cutoffs is selected.
After fixing $k$, include its tests and choose the common projection
data on the same filter, without redefining the intrinsic odd minimum.
Stabilize the finite alphabet and its two projections in this way.
\end{proof}

\begin{proposition}[Stationary and simultaneous address constructions]\label{ex:app-simultaneous-address}
Use the definitions and realization freedoms in
Definitions~\ref{ex:sec7-typed-graph}, \ref{ex:sec7-symbol},
\ref{ex:sec7-address} and Lemma~\ref{ex:sec7-realization},
and the coefficient rules of Proposition~\ref{ex:app-coefficient-rules}.
Fix $b\ge1$. Either take one symbol on a nonempty set of always
active terminals with a nonzero polynomial in their pair bits, or
take one or two symbols on terminals with nonzero linear activations,
including a terminal of every nonzero activation, such that each
symbol separately has a nonzero terminal polynomial at every
$x\ne0$. In the latter case the symbols may read different
compatible projections of the common alphabet and fixed stationary
external data. Require the active data to satisfy each prescribed
filter at every nonzero address.
Starting from fixed compatible typed terminal and stationary
data, one can adjoin finitely many primary auxiliaries with nonzero
linear activations and choose only incidences involving them so that
all the symbols equal $1$ at every $x\ne0$. All original data stay
fixed, no auxiliary is present at zero, and the required neutral
external incidences and splitting conditions are retained.
Only nonempty configurations and nonempty contraction outputs are
used. A nonzero iterated difference, hence a partitionability
certificate, suffices for the required nonzero terminal polynomial.
For the two final symbols the zero-total packets below supply these
terminal hypotheses using only separate projection surjectivity.
For the even stationary application, the original prime representatives
may be fixed at the outset. For a stabilized odd symbol, first fix the
abstract data and the finite budget, then choose an actual stage
carrying them, and only then choose prime representatives. Previously
fixed representatives may be retained when compatible with the full
type at that stage; arbitrary earlier primes are not asserted to
satisfy a later moving condition.
\end{proposition}
\begin{proof}\leavevmode
\paragraph{Total labels of the final seeds.}
Both final symbols have projected total label zero. The intrinsic
odd filter contains $1$: its fixed unit squareclasses, positive sign,
and moving condition $(D/\ell_0)=1$ are those of $1$. The even source
for $E^k$ is likewise normalized in the unit filter containing $1$.
Every nonzero coefficient, and hence every unit seed, lies in its
respective filter. Additional stationary tests at primes dividing
$k$ refine the common alphabet but are forgotten by the intrinsic
odd projection. They do not add conditions to that odd symbol.

Take repeated packets with total zero in the common elementary
label group, including sufficiently many copies of every label.
For either projection separately, keep the first $r-1$ anchors of
its unit seed and contract the nonempty remainder into the last
anchor. Since the total projected label is zero, the remainder has
the last seed label. This argument uses surjectivity of each
projection separately. It does not require that a fixed old odd
seed admit lifts whose total is zero in the common alphabet.
The whole packet already has that total, and contractions preserve
it. Thus the final twists satisfy the stationary tests as well as
both coefficient filters.
Concretely, choose a packet with an even number of copies of every
common label, and enough copies to provide the anchors of either
seed under its respective projection.  Its total label is zero.
For every nonzero activation $\lambda$, take a disjoint copy of
this whole packet and give all its vertices activation $\lambda$.
At every $x\ne0$ there are exactly $2^{b-1}$ active packets.
The active list is therefore nonempty, has total common label zero,
and contains enough anchors for each projected seed separately.
Its terminal polynomial for each symbol is nonzero by the
contraction certificate, regardless of the initially fixed bits.
It also contains a terminal of each required activation.  This
constructs the terminal input of the address lemma explicitly.
No simultaneous lifting of a prescribed pair of seed labels is
needed.  Every added network vertex is primary, so the total
common label remains zero at each address.

\paragraph{Network operators and their degree.}
We spell out the simultaneous argument, including the product over
all addresses.  Write $\partial_e$ for the finite difference in a
terminal pair bit $e$.  Products of these operators commute and
$\partial_e^2=0$.  The contraction rule identifies a forest product
with its contracted graph.  A cycle product is zero: after a
spanning path is contracted its last edge is internal and changes
neither a type nor any external incidence.  No contraction removes
the last vertex.

A network of primary auxiliary vertices acts on any of the symbols
by an operator in the $\partial_e$.  To see this directly, expand in
the bits incident to the auxiliary vertices.  Contract the chosen
forest edges.  A component with no terminal becomes an isolated
primary vertex, and its contribution vanishes by primary isolation.
A component meeting several terminals gives their forest difference.
If $Q$ is the grounded auxiliary Laplacian and $c_a$ is its attachment
column at terminal $a$, the degree-zero and degree-one terms are
\[
 \det(Q)\,1+
 \sum_{a<b}c_a^{\mathsf t}\operatorname{adj}(Q)c_b\,
                    \partial_{ab}.
\]
All other terms have difference degree at least two.  The displayed
coefficients are the rooted-forest minors, by Cauchy--Binet on the
incidence matrix.  They are polynomial identities, so apply to
singular $Q$ as well.  Distinct networks have no edges between their
auxiliary vertices.  Their operators multiply: expanding both
forests gives their union, and unions containing a terminal cycle
give zero by the preceding observation.  This also shows that the
operator is the same for both projected symbols.

Fix a target $x_*\ne0$.  Choose a basis
$\lambda_1,\ldots,\lambda_b$ of the dual address space with
$\lambda_i(x_*)=1$ for all $i$.  Such a basis is obtained from one
functional taking value one by adding it to a basis of the space
of functionals vanishing at $x_*$.  Group $i$ has three primary vertices, each activated
by $\lambda_i$.  In its physical three-dimensional coordinate space
let $e_i=(1,1,1)$ and choose $u_i,v_i$ spanning its orthogonal
complement.  The vector $e_i$ has norm one.  Prescribe the symmetric
matrix $Q$ by its bilinear matrix $I_b$ on the $e_i$ and
\[
 H(T,B)=\begin{pmatrix}0&T\\T^{\mathsf t}&B\end{pmatrix}
\]
on the remaining coordinates.  In physical coordinates each
$h_i=Qe_i=e_i$ is supported on its own group.  Attach this column to
a terminal activated by $\lambda_i$, which exists by hypothesis.
For an active set $J$ the sum of these active attachment columns is
$Q_J\mathbf1_J$.  Thus the prescribed principal matrix $Q_J$ really
is the grounded graph Laplacian on that address, including its
diagonal entries.

For a selected terminal pair $a,b$ active at $x_*$, let their
activations be $\alpha,\beta$.  Attach an extra column $c$, dual to
$u_1$, to both.  If $\alpha\ne\beta$, attach the same column to a
terminal of activation $\alpha+\beta$; it is absent at $x_*$.  The
extra columns sum to zero at every address since
$\alpha+\beta+(\alpha+\beta)=0$.  If $\alpha=\beta$, the two
columns already cancel.  Hence they do not alter $Q_J$.

For the singular network set $T=I+P$, where $P$ is the cyclic shift,
and give $B$ one diagonal entry one.  On every proper group subset,
the restricted cyclic shift is a union of paths, so $T_J$ is
invertible and $\det Q_J=\det(T_J)^2=1$.  On all groups, the kernels
of $T$ and $T^{\mathsf t}$ are generated by the all-one vector.
The one nonzero diagonal entry of $B$ forces the lower component of
a vector in $\ker H$ to vanish.  Thus $\ker Q$ is precisely the line
generated by $r=\sum_i u_i$, and
$\operatorname{adj}(Q)=rr^{\mathsf t}$.  The anchor columns pair
to zero with $r$, while $c^{\mathsf t}r=1$.  Only $a,b$ receive
active extra columns at $x_*$.  Consequently its network operator
$S_{x_*,e}$, for $e=ab$, has the expansions
\[
 S_{x_*,e}|_{x_*}=\partial_e+\text{terms of degree at least two},
 \qquad
 S_{x_*,e}|_x=1+\text{terms of positive degree}\quad(x\ne x_*).
\]
For $b=1$, $T=0$ and $B=(1)$ give exactly the same calculation.

For a varying network set $T=I+L$, with $L$ the consecutive lower
shift, and $B_\varepsilon=\varepsilon E_{bb}$ for
$\varepsilon\in\{0,1\}$.  Every $T_J$ is invertible and both
operators $V^\varepsilon_{x_*,e}|_x$ have constant coefficient one.
The upper-left block of $H^{-1}$ is
$T^{-\mathsf t}B_\varepsilon T^{-1}$.  Its change on the selected
columns is $((T_J^{-1})_{b1})^2$, interpreted as zero when an endpoint
is absent.  This is one exactly when the entire path is present,
which means $x=x_*$.  Anchor columns have zero change.  Therefore
\[
 \begin{aligned}
 (V^0_{x_*,e}+V^1_{x_*,e})|_{x_*}
       &=\partial_e+\text{terms of degree at least two},\\
 (V^0_{x_*,e}+V^1_{x_*,e})|_x
       &=\text{terms of degree at least two}\quad(x\ne x_*).
 \end{aligned}
\]

\paragraph{One choice for both symbols at every nonzero address.}
For each $x\ne0$ let $P_{1,x},P_{2,x}$ be the two nonzero terminal
polynomials, of degrees $d_{1,x},d_{2,x}$, and let $H_{s,x}$ be their
highest homogeneous parts.  Work in the ordinary polynomial ring
in the terminal pair variables for this address.  Since it is a
domain, $H_{1,x}H_{2,x}\ne0$.  Choose a coefficient-one monomial
\[
 M_x=\prod_e X_e^{n_{x,e}},\qquad n_{x,e}\in\{0,1,2\},\qquad
 \sum_e n_{x,e}=d_{1,x}+d_{2,x}.
\]
For $n_{x,e}=2$ install one singular network $S_{x,e}$; it acts on
both symbol factors.  For $n_{x,e}=1$ install one varying network
$V^{\varepsilon_{x,e}}_{x,e}$.  Use disjoint auxiliary sets for all
these networks.  Let $\Gamma_y(\varepsilon)$ be the resulting graph
at address $y$, and consider the finite sum in $\mathbf F_2$
\[
 \sum_{\varepsilon}
       \prod_{y\ne0}\prod_{s=1}^2 A_s(\Gamma_y(\varepsilon)).
                                                               \tag{G.1a}
\]
The total degree available to all terminal polynomials is
$D=\sum_{y\ne0}(d_{1,y}+d_{2,y})$.

Each singular network at $x$ has zero constant at $x$, so its action
costs at least one difference in each of the two factors there, hence
at least two in (G.1a).  For a varying network, summing over its
$\varepsilon$ cancels the all-constant term.  Its first-degree
terms are independent of $\varepsilon$ away from its target, and
their difference at the target is $\partial_e$.  Thus every
surviving term costs at least one difference, and a cost-one term
acts at its target on exactly one of the two factors.  The sum of
all these minimum costs is exactly $D$.  Any higher-degree action,
or any additional action at a wrong address, exceeds $D$ and
annihilates at least one terminal polynomial.  Hence only the
minimum-cost terms survive.

At address $x$, a variable with exponent two is therefore
differentiated once in each factor; a variable with exponent one
is differentiated in exactly one factor, summed over the two
possibilities.  The resulting scalar is precisely
$[M_x](H_{1,x}H_{2,x})=1$.  Lower-degree parts cannot contribute.
Taking the product over all $x\ne0$ proves that (G.1a) equals one.
At least one choice of the finitely many $\varepsilon$ consequently
makes every factor in that product one.  This is one graph for both
symbols at every nonzero address.  Boolean reduction of
$H_{1,x}H_{2,x}$ before selecting $M_x$ would destroy this
exponent-two accounting and is not used.
Here is the stationary and one-symbol specialization, so that the
application in Proposition~\ref{ex:app-rough-companion} uses the same
proved construction. If all terminals are always active, attach
each column $h_i=Qe_i$ to any one fixed terminal. Since $h_i$ is
supported on its own group, their restricted sum is
$Q_J\mathbf1_J$ at every address. The two additional columns $c$
at a selected terminal pair cancel directly; no balancing terminal
of activation $\alpha+\beta$ is needed. The same singular program
therefore isolates that pair difference at its target address.
For one symbol, choose a coefficient-one top monomial of each
nonzero $P_x$ and put a separate singular program at $x$ for each
variable in that monomial. Its own programs already spend
$\deg P_x$ contraction degrees at $x$; any extra degree kills the
term. Thus exactly the selected top coefficient survives, and it
is one. A constant nonzero $P_x$ needs no program. This proves the
stationary one-symbol case as well as the moving one-symbol case.
In all cases the primary realization of
Lemma~\ref{ex:sec7-realization} changes only incidences with
the newly adjoined vertices. These are finite constructions for
each $b$, with no assertion of one infinite prime array.

\end{proof}

\begin{proposition}[One actual stage for the final family]\label{ex:app-final-family}
In the nonempty-seed case of Proposition~\ref{ex:app-nonempty-seeds},
fix the final $k,K$ and the stabilized common alphabet. For every
sufficiently large $b$, the address construction is realized at one
actual theta stage by a finite set $Q$ of distinct good odd primes
split in $K$, avoiding the fixed support and the specified finite
exclusions. With nonzero activations $\lambda_q$ it gives exactly
the family in (3.9),
\[
 h_x=\prod_{q\in Q}(q^*)^{\lambda_q(x)},\qquad h_0=1,
 \qquad q^*=(-1)^{(q-1)/2}q.
\]
All fixed local types are retained, and for every $x\ne0$,
\[
 \operatorname{ord}_{s=1}L(E^{h_x},s)=1,\quad L_*(h_xk)\ne0,
 \qquad j_{\rm gen}(h_x)\le\omega(h_x)+C_o,\quad
 v_2L_*(h_xk)\le2\omega(h_x)+C_e.
\]
The constants precede $b$. The finite graph budget and then the
actual theta precision are chosen after $b$. The theta prime remains
a coefficient prime and exclusion, not a relaxed Selmer prime.
\end{proposition}
\begin{proof}\leavevmode
For each chosen $b$, this construction has a finite budget $W_b$.
Enlarge it to bound every graph product and squarefree output in its
finite auxiliary tests, as well as every operator set. Contractions
only decrease the number of vertices. Thus every final $h_x$ lies in
the tested range of (F.7); this enlargement changes no fixed constant.
Only now take one actual theta stage in the finite intersection
of the selected table sets through $W_b+1$, the valid-stage set
for the full budget $|U|\le W_b$ with $R=2$, and the required
precision and finite exclusion sets.  In particular (F.6) holds
through the threshold $T$ used for that budget.
Realize its prescribed forward residue bits by distinct primes using CRT and Dirichlet;
quadratic reciprocity supplies the reverse bits.
There is no uniform bound on $W_b$, no intersection over all
budgets, and no product-surjectivity assumption on the alphabets.
The resulting one family satisfies (3.9), with constants fixed
before $b$. The theta prime $\ell_0$ is a coefficient prime and
an exclusion, not an active relaxed Selmer prime or a moving
arithmetic prime $r_i$.
\paragraph{Finite intersection and the upper bound.}
The required stage conditions form a finite list once the address
program is fixed. Let $\mathcal U$ be the chosen ultrafilter and
let $A_1,\ldots,A_t\in\mathcal U$ be the selected truth-table
sets used by that program, including all intermediate graphs.
Let $B\in\mathcal U$ impose (F.6) through the threshold $T$
for its operator budget and impose
$M\ge\max(0,\delta_o+W_b)+2$ and $M>T+\delta_o$.
Finite prime exclusions are included in $B$.
Then $B\cap\bigcap_{i=1}^t A_i\in\mathcal U$, hence is nonempty.
Choose a stage there before realizing the prime labels. The
coefficient rules hold for every realization at this stage, not
merely for a previously selected finite list of Fourier indices.
The address construction makes each nonzero address a unit of
both sources. For the odd source this says exactly
\[
 v_2a_M(h_x)=\delta_o+\omega(h_x)<M.
\]
The strict inequality follows from the tested-range precision
condition. Applying (O.8) gives the simple zero and
$j_{\rm gen}(h_x)\le\omega(h_x)+\delta_o+C_{\rm det}$.
The normalized even source gives
$v_2L_*(h_xk)\le2\omega(h_x)+C_{e,\rm final}$.
Thus one may take $C_o=\delta_o+C_{\rm det}$ and
$C_e=C_{e,\rm final}$, both fixed before $b$. Only the number
of tests, the threshold $T$, and the chosen stage depend on the
address dimension.
\end{proof}

\subsection{One coordinate in the final contradiction}
The family in Proposition~\ref{ex:app-final-family} is now used with
the same coordinate of Lemma~\ref{ex:arith-common-coordinate}
and its central calibration in
Proposition~\ref{ex:arith-central-calibration}. The denominator and
coefficient comparisons are those of
Propositions~\ref{ex:limits-fixed-field-clearing}
and~\ref{ex:limits-coefficient-transfer}; the zero-base assertion is
Lemma~\ref{ex:arith-base-zero}.
For this fixed $K$, (3.4), (F.7), and the even bound just proved give
\[
 2j_{\rm rc}(h_x,k)\le4\omega(h_x)+C_J,\qquad
 C_J=2C_o+C_{e,\rm final}-v_2L_*(h_*)+C_{H,\rm final}
\]
at every $x\ne0$.
The finite support comparison uses the same class, duality map,
bases and coefficient stages. Its integral correction
\[
 \eta=\prod_q(1+v_q(1+g_q)),\qquad l_q/e_q^2=1+2v_q,
 \qquad \widetilde u=\eta u_{\rm rel}
\]
is formed before characters. The actual integral determinant
triangle gives $\widetilde u_x=\epsilon_xu_{{\rm prim},x}$ with
$\epsilon_x\in\Lambda^\times$, rather than independently chosen
units in $\mathscr O^\times$. Therefore the central identity is
\[
 \begin{split}
 v_2\widetilde u_x(0)={}&2j_{\rm rc}(h_x,k)-4\omega(h_x)
       -s(h_x)-s(h_xk)\\
       &-\sum_{p\mid k}t_p+\epsilon'_x,\qquad
 0\le t_p\le2,\quad |\epsilon'_x|\le C_{\rm cal}.
 \end{split}                                                    \tag{F.8}
\]
The nonnegative finite $\Sha$ terms give one upper bound
$v_2\widetilde u_x(0)\le B$ for all $x\ne0$, independent of $b$.
At the unknown base no finiteness of $\Sha$ is assumed.
The retained central class and corank-one argument instead give
$\widetilde u_0(0)=0$ if the base derivative vanishes.

The fixed-field construction of Section~\ref{sec:limits} has
\[
 \begin{gathered}
 c_{\rm det}=a+4(h_K+2c_E+2t_K),\quad
 L_T(M)=16(d_K+5)M,\\
 f_x\widetilde u_x=d_x,\qquad
 c=v_2f_0(0)\le c_{\rm det}.
 \end{gathered}
\]
Choose $M>\max(c_{\rm det},B+2c_{\rm det},0)$, then
$b>(L_T(M)+2)(2^M-1)$ large enough for the address construction,
then its finite graph budget and actual theta stage $M_\theta$.
The residual comparison tests coefficients of $f$ from
$-L_T(M)$ through zero and only the constant coefficient of $d$.
It gives $x\ne0$ with
\[
       v_2(\widetilde u_x(0)-\widetilde u_0(0))\ge M-2c>B.
\]
The estimate uses both $v_{\mathscr O}(f_x)\le c$ and
$\widetilde u_x\in\Lambda[1/2]$; neither a bound on the whole
numerator nor a finite-Artin inverse of $t$ is asserted.
This contradicts (F.8) under the supposed base zero.
The Rankin factorization and $L(E^k,1)\ne0$ give the asserted
simple zero of $E$.

This completes the proof of the restricted converse,
Theorem~\ref{hyp:RC}, from the displayed classical inputs and
the arithmetic, coefficient and determinant constructions.
The conversion from that converse to pole witnesses, the integer
model, the natural compiler and finally the ordinary rational value
is the separate argument in the main text.

% End 09-normalization-endpoints.tex


\section{De Rham comparison and finite descent in the height argument}
\label{app:de-rham-descent}

\begin{lemma}\label{lem:de-rham-descent}
Let $K/\mathbb Q_\ell$ and $F/K$ be finite extensions, let $A/F$
be an abelian surface, and let $C_\ell/\mathbb Q_\ell$ be a finite
coefficient field. Suppose a continuous representation
$r:G_K\longrightarrow\mathrm{GL}_2(C_\ell)$ satisfies
\[
 V_\ell(A)\otimes_{\mathbb Q_\ell}C_\ell
       \simeq r|_{G_F}\oplus r|_{G_F}.
\]
Then $r$ is de Rham. At every coefficient embedding its
Hodge--Tate weights are $0,-1$, with
$\mathbb Q_\ell(1)$ assigned weight $-1$.
\end{lemma}

\begin{proof}
The smooth proper variety $A$ has a smooth proper rigid
analytification. Apply \cite[Corollary 1.8]{scholze2013}, with the
corrections in \cite{scholze-erratum}. The passage from algebraic
to analytic cohomology uses \cite[Theorems 9.1 and 9.3]{scholze2013}:
Theorem 9.1 identifies coherent cohomology, and hence the filtered
de Rham hypercohomology, while Theorem 9.3 identifies the
$\mathbb Z/\ell^n\mathbb Z$-cohomology, compatibly in $n$.
Its invertibility condition holds because $F$ has characteristic
zero. Taking inverse limits and then inverting $\ell$ identifies
the algebraic and analytic $\ell$-adic cohomology, compatibly with
the action of $G_F$. Thus
$H^1_{\mathrm{et}}(A_{\overline F},\mathbb Q_\ell)$ is de Rham.

The canonical, Galois-equivariant identification
\cite[Part I, Theorem 12.1(a)]{milne-av}
\[
 V_\ell(A)\simeq
 H^1_{\mathrm{et}}(A_{\overline F},\mathbb Q_\ell)^\vee
\]
gives the same property for the Tate module. De Rham representations
are closed under duals and subquotients
\cite[Theorem 5.2.1 and Proposition 6.3.3]{brinon-conrad}.
Extension of the coefficient field is, on underlying
$\mathbb Q_\ell$-representations, a finite direct sum.
The displayed decomposition therefore makes $r|_{G_F}$ de Rham.

Here finite-extension descent requires no semistable descent
assertion. To see it explicitly, pass to a finite Galois extension
$F'/K$ containing $F$, and let $W$ be the underlying
$\mathbb Q_\ell$-representation of $r$. The vector space
\[
 D_{F'}=(B_{\mathrm{dR}}\otimes_{\mathbb Q_\ell}W)^{G_{F'}}
\]
has its natural semilinear $\mathrm{Gal}(F'/K)$-action, and its
invariants are
$D_K=(B_{\mathrm{dR}}\otimes_{\mathbb Q_\ell}W)^{G_K}$.
Finite Galois descent gives $F'\otimes_K D_K\simeq D_{F'}$,
also on each filtration step. Since $W|_{G_{F'}}$ is de Rham,
$\dim_K D_K=\dim_{F'}D_{F'}=\dim_{\mathbb Q_\ell}W$.
This proves the assertion over $K$; it is the finite-extension
case of \cite[Proposition 6.3.8]{brinon-conrad}.

The two Hodge numbers of $H^1_{\mathrm{dR}}(A/F)$ are $2,2$.
Consequently the Tate module has weights $0,0,-1,-1$.
At each coefficient embedding, its two isomorphic summands
have the same weight multiset, so each summand has weights $0,-1$.
The filtration-compatible descent preserves these weights.
\end{proof}

The Hodge--Tate comparison in \cite[Chapter III, Theorem 4.1]{faltings1988}
alone would not prove the de Rham assertion. The preceding argument
uses the de Rham comparison itself and does not infer de Rhamness
from Hodge--Tate weights.

% Begin references.tex
% Original exposition of the prime-pattern argument, prepared for Version 5.
% Requires amsmath, amssymb, amsthm, hyperref and the manuscript theorem environments.
% Publication-status qualifications in this file are intentional.
